\chapter{Formal geometry and the compact comparison}
\label{chap:dirac-igusa-pro-object}

The denominator algebra supplies a target for a geometric Hall
construction. To reach it, one must first construct coefficients
on moduli of objects, maps on extension flags and compatible
orientation data. The product of their signed dimensions is a
consequence to compare, not a definition of those operations.
This chapter begins with the formal deformation model, then
constructs or specifies the maps needed by the compact problem.

Let \(S\) be a smooth projective K3 surface, \(E\) a smooth
elliptic curve and \(X=S\times E\). Fix the polarization and
stability data, the elliptic origin and the support conditions of
Appendix~\ref{app:stability-support}. Singular K3 limits, singular
elliptic curves, imprimitive classes, the compact Hall cusp and
multiple-component degenerations need their own coefficient and
specialization maps. Results below are used on the stated smooth
and bounded recipients.

\section{The source and its finite comparisons}

Three established reference constructions will recur: the
normalized Borcherds product \(D_5=\Delta_5/64\), the denominator
algebra \(\mathfrak g_{\Delta_5}\) with denominator \(D_5(2Z)\),
and the regularized theta lift. A geometric Pfaffian, source
coalgebra or compact state must be constructed independently
before being compared with these targets.

At a stage \(R\), the proposed geometric datum is
\[
 \begin{split}
 \mathfrak D^{\mathrm{DI}}_{X,R}=
 (&\mathcal A_{X,R}^{E_3},F_{X,R}^{\mathrm{hyb}},
   \Gamma_{X,R},\Pi_X,\Phi_R,o_R,H_R,P_R^\Pi,\\
  &C_{X,R},\Theta_{\mathrm{Kos},R},
   \mathscr L_{\mathrm{Pf},R},\operatorname{pf}_{X,R},
   \epsilon_{o,R},\mathrm{Rec}_R).
 \end{split}
\]
The entries denote, respectively, a field complex when that
refinement is required, the local and wrapped carrier, charges
and their Gram map, critical coefficients and orientation,
Hall operations and primitives, a source coalgebra and its
associative comparison, a Pfaffian line and section, Weyl
transport and primitive recognition. Their precise definitions
and morphisms are developed below. A compatible cofiltered
system of such stages is the Dirac--Igusa pro-object
\(\{\mathfrak D^{\mathrm{DI}}_{X,R}\}_R\). An effective
geometric inverse limit is an additional construction.

\subsection*{Formal charges and geometric supports}

The Mukai--Gram map \(\Pi_X\) is quadratic. Its polarization
\(B\) defines an additive normal-ordered central lift
\(\overline\Pi_X\), as in Appendix~\ref{appendix:mukai-gram-dictionary}.
Pointwise lattice lifts of Gram triples and the root assignment
\[
 \alpha(n,l,m)=2nf_2-lf_3+2mf_{-2}
\]
are explicit. They supply no effectivity or semistability statement
for a compact sheaf. The raw zero fibre
\(K_{\Pi,0}=\Pi_X^{-1}(0,0,0)\) is not an additive subgroup,
so it cannot be used as a quotient charge group.

After an actual additive regrading has been fixed, summing finite
vector spaces over its fibres is exact and retains even and odd
parts separately. It also preserves a supplied filtration and
triangular decomposition. The tensor-coalgebra comparison uses
that additive grading; the raw quadratic Gram map does not
commute with tensor degrees. These linear constructions retain
their precise scope in
Proposition~\ref{prop:recognition-additive-regrading}. They do
not prove that a geometric Hall bracket, pairing radical or
coefficient functor descends to the regraded source.

A finite target test window therefore starts from actual bounded
geometric charges, their normal-ordering translates and the image
identity in Definition~\ref{def:finite-target-test-window}.
A geometric exhaustion must represent every required relation
degree and its intermediate outputs. The first simple target
triples are
\[
 \gamma(\delta_1)=(1,1,0),\qquad
 \gamma(\delta_2)=(0,1,1),\qquad
 \gamma(\delta_3)=(0,-1,0).
\]
The sheaf representatives of
Proposition~\ref{prop:type-ii-compact-mukai-representatives}
realize these Gram values. Their wall semistability, reduced wall
charts and operation comparisons are further geometric data.

\subsection*{The moduli and coefficient hypotheses}

The Abramovich--Polishchuk heart and its Liu tilt retain the
polynomial charge, weak stability function, torsion pair and
central charge together. A finite subobject- and quotient-closed
collection with finitely many isomorphism classes and a strict
integer length has the noetherian and HN properties of
Propositions~\ref{prop:retained-window-noetherian} and
\ref{prop:retained-window-hn-filtration}. The maximal-phase,
maximal-length construction bounds the number of HN factors by
the initial length. If length is additive, their lengths sum to
that length. This finite collection need not be additive or
closed under all extensions.

Finite geometric stages require additional bounded class data,
finite-type loci and the object and flag maps actually used.
Semistability need not be closed under every specialization.
The ambient CY3 moduli stack has a shifted symplectic structure;
its scalar directions prevent an automatic two-term reduction on
all substacks. The universal complex may require a scalar gerbe
or a twisted recipient, even when its weight-zero endomorphism
complex descends. Scalar rigidification does not force finite
residual inertia. Closed strata, retained extensions, HN factors
and geometric duality must be supplied at the support and heart
where they are asserted.

These requirements constitute the finite-moduli and reduced-atlas
hypotheses. They are not consequences of a finite charge list.
The reduced critical object is denoted \(\Phi_o\): its
orientation enters through the sign-torsor construction, not an
ordinary tensor of a constructible complex with a coherent line.
The common translation quotient, proper or compact-support maps,
cosections, roots and full flag coherences belong to the same
coefficient system.

If an \(E_3\)-field refinement is sought, its formal target
charts, field complexes and operadic maps are additional input.
A BV realization requires its quantum master equation, anomaly
and framing conditions. The proposed
\(\operatorname{Sp}^{\mathrm{ch}}_{K3,E}\) must be constructed
on these complexes and preserve compact supports, factorization
and stage transitions. A current-envelope target does not define
this specialization functor.

\subsection{Formal sheaf models and elliptic support}
\label{subsec:hybrid-ran-carrier}

The spherical collection gives a local algebra of sheaf deformations.
To obtain Hall operations, this algebra must describe extensions of
several objects at once and retain their traces. We first construct
that formal comparison. We then distinguish sheaves supported over
finitely many points of $E$ from sheaves whose support dominates $E$.
Their extension diagrams can be formed together, and elliptic
translation acts on the entire diagram.

\subsubsection*{One cyclic model for the spherical collection}

Use the common-polynomial line bundles $L_1,L_2,L_3$ of
Theorem~\ref{k3:thm:common-triple}. Fix a nonzero holomorphic
form $\Omega_S$. Their Dolbeault dg category has morphisms
\[
 \mathcal D(i,j)=\mathcal A^{0,*}(S,L_i^\vee\otimes L_j),
 \qquad d=\bar\partial,
 \qquad
 \langle a,b\rangle=\int_S\Omega_S\wedge\operatorname{tr}(ba).
\]
The pairing has degree $-2$ and induces Serre duality. Set
$A(i,j)=\operatorname{Ext}_S^*(L_i,L_j)$ and
$\ell=\bigoplus_i\C1_i$. Thus $A^0=\ell$, $A$ is concentrated in
degrees $0,1,2$, and $A^2$ is dual to $\ell$.

\begin{proposition}[Cyclic formality of the collection]
\label{hall:prop:cyclic-collection}
There is a strictly unital cyclic minimal model on $A$ with only its
Yoneda multiplication $m_2$ nonzero. One cyclic $A_\infty$
quasi-equivalence $F:A\to\mathcal D$ can be chosen for the whole
collection, before choosing multiplicities or flags.
\end{proposition}

\begin{proof}
The cyclic minimal-model and strictification results of
\cite[Lemmas~4.4--4.5, Proposition~4.8 and Corollary~4.9]
{DavisonPurity2CY} apply to this simple-minded $2$-Calabi--Yau
collection. The nonnegative-degree hypothesis for strictification
holds. They give the cyclic model with strict vertex units.

For $n\geq3$, an input in $A^0=\ell$ makes $m_n$ vanish by strict
unitality. If all inputs have positive degree and one has degree two,
the output has degree at least three, since $|m_n|=2-n$.
The remaining possibility has every input of degree one and output
of degree two. Pair that output with a vertex unit. Cyclicity moves
the unit into an input, making the value zero. Perfection of the
pairing $A^2\otimes A^0\to\C$ forces the output to vanish.
\end{proof}

For multiplicity spaces $V_i$, write
$F_{\mathbf n}=\bigoplus_iL_i\otimes V_i$ and $n_i=\dim V_i$.
Matrix amplification replaces $A(i,j)$ by
$A(i,j)\otimes\operatorname{Hom}(V_i,V_j)$, composes morphisms
by matrix multiplication, and composes the cyclic trace with matrix
trace. Write $A_{\mathbf n}$ for this matrix amplification.
The group $G_{\mathbf n}=\prod_iGL(V_i)$ acts by changing bases.

Choose a quiver $Q$ with two arrows from $i$ to $j$ for each $i<j$.
Its double has dual variables
$X_a:V_{t(a)}\to V_{h(a)}$ and
$Y_a:V_{h(a)}\to V_{t(a)}$. With dual arrow bases chosen according
to the Serre pairing, the surface Maurer--Cartan equation is
\begin{equation}
\label{hall:eq:preprojective-moment}
 \mu_i(X,Y)=\sum_{h(a)=i}X_aY_a-\sum_{t(a)=i}Y_aX_a=0.
\end{equation}
For a local Artinian base with maximal ideal $\mathfrak m$, the
pointed deformation problem uses $\mathfrak m$-valued arrows and
the formal gauge group
$\exp(\mathfrak g_{\mathbf n}\otimes\mathfrak m)$, where
$\mathfrak g_{\mathbf n}=\bigoplus_i\operatorname{End}(V_i)$.
The full affine quotient by $G_{\mathbf n}$ is a separate algebraic
quiver model. Its comparison with sheaf moduli here is formal near
the specified polystable object.

\subsubsection*{The normal operator and the derived critical quotient}

Embed $S$ as the fibre over $e\in E$ and choose a normal coordinate
at $e$. The normal Koszul resolution adds an exterior generator
$\theta$ of degree one. The formal threefold algebra is
$A_{\mathbf n}\otimes\Lambda(\theta)$.
A degree-one element $x+z\theta$ has square
\[
 (x+z\theta)^2=x^2+(xz-zx)\theta.
\]
The first equation is \eqref{hall:eq:preprojective-moment}; the
second says that the normal endomorphisms commute with the arrows.
Both are critical equations of
\begin{align}
 \mathcal W_{\mathbf n}(X,Y,Z)
 &=\sum_i\operatorname{Tr}(Z_i\mu_i)\label{hall:eq:tripled-potential}\\
 &=\sum_{a\in Q_1}\operatorname{Tr}
       \bigl(Y_a(Z_{h(a)}X_a-X_aZ_{t(a)})\bigr).
 \nonumber
\end{align}
The potential is defined on
$\mathcal V_{\mathbf n}=\operatorname{Rep}(\overline Q,\mathbf n)
\oplus\mathfrak g_{\mathbf n}$. Its equations are
\[
 \mu_i=0,\qquad Z_{h(a)}X_a=X_aZ_{t(a)},\qquad
 Y_aZ_{h(a)}=Z_{t(a)}Y_a.
\]
Taking the derived critical locus on the smooth quotient
$[\mathcal V_{\mathbf n}/G_{\mathbf n}]$ retains gauge symmetry.
At a critical point its endomorphism-degree complex is
\begin{equation}
\label{hall:eq:four-term-gauge-complex}
 \mathfrak g_{\mathbf n}\xrightarrow{\rho}\mathcal V_{\mathbf n}
 \xrightarrow{\operatorname{Hess}\mathcal W_{\mathbf n}}
 \mathcal V_{\mathbf n}^*\xrightarrow{-\rho^*}
 \mathfrak g_{\mathbf n}^*.
\end{equation}
The degrees are $0,1,2,3$; the tangent complex is its shift by $[1]$.
Invariance of $\mathcal W$ gives
$(\operatorname{Hess}\mathcal W)\rho=0$ and its dual identity.
The outer terms record automorphisms and their duals. Omitting them
changes the derived critical quotient.

For $N=\sum_i n_i>0$, set
\[
 c=N^{-1}\sum_i\operatorname{Tr}Z_i,\qquad Z_i^0=Z_i-cI_{V_i}.
\]
Since $\sum_i\operatorname{Tr}\mu_i=0$, the potential is
independent of $c$. Therefore
\[
 \operatorname{Crit}_{[\mathcal V_{\mathbf n}/G_{\mathbf n}]}(\mathcal W)
 \simeq T^*[-1]\mathbb A^1_c\times
 \operatorname{Crit}_{[\mathcal V^0_{\mathbf n}/G_{\mathbf n}]}(\mathcal W^0).
\]
Quotienting translations of $c$ removes its tangent direction; the
centred critical model removes its dual obstruction direction as
well. These are different operations. On a flag with block dimensions
$N_i$, removing the total centre imposes $\sum_i N_ic_i=0$ and
leaves the relative centres.
The zero object has $N=0$ and requires its separate unit construction.

The ambient quotient has canonical line
$\det(\mathcal V_{\mathbf n}^*)\otimes\det\mathfrak g_{\mathbf n}$.
The opposite arrow spaces cancel in its determinant character, as do
the loop and gauge terms. An ordering of the hyperbolic pairs gives
a local trivialization. Matching this orientation with a prescribed
compact orientation is a further comparison, even when the formal
cyclic algebras have been identified.

\subsubsection*{Dimensional reduction and the point sector}

Write $\varphi_f=\phi_f[-1]$ for shifted vanishing cycles, so that
$\varphi_0\Q=\Q$. In the mixed Hodge version, $(-r)$ denotes a
Tate twist. The potential above is linear in the $r=\dim\mathfrak
g_{\mathbf n}$ loop coordinates. This gives a coefficient
comparison after compactly supported pushforward.

\begin{proposition}[Dimensional reduction of the tripled potential]
\label{hall:prop:dimensional-reduction}
There is a natural comparison
\[
 R\Gamma_c(\mathcal V_{\mathbf n},
                 \varphi_{\mathcal W_{\mathbf n}}\Q)
 \simeq R\Gamma_c(\mu^{-1}(0),\Q)[-2r](-r).
\]
The comparison is $G_{\mathbf n}$-equivariant, interpreted through
finite-dimensional approximations. On the centred model replace
$r$ by $r-1$ when $\sum_i n_i>0$; the zero locus on the right
remains $\mu^{-1}(0)$. The zero-dimensional representation is
left unchanged.
\end{proposition}

\begin{proof}
For $f(x,z)=\sum_{i=1}^r z_if_i(x)$, the dimensional-reduction
theorem identifies $\pi_!\varphi_f\pi^*K$ with
$\pi_!\pi^*i_*i^*K$, where $i$ includes the common zero locus
of the $f_i$ and $\pi$ is the rank-$r$ vector-bundle projection
\cite[Theorem~A.1 and Corollary~A.9]{DavisonCriticalCoHA}.
Apply it to the moment-map coordinates and $K=\Q$.
Compactly supported cohomology of the affine fibre is
$\Q[-2r](-r)$, giving the formula. On the trace-free loop space
the potential pairs with the projection of $\mu$ to its dual.
That projection has the same zero locus: $\mu$ already has
total trace zero. The fibre rank is now $r-1$.
\end{proof}

For the preprojective stack
$[\mu^{-1}(0)/G_{\mathbf n}]$, the dual compactly supported
cohomology is pure of Tate type \cite[Theorem~A]{Davison2017}.
This controls its mixed Hodge weights; it does not identify a
compact reduced coefficient system with the preprojective one.

This is a pushforward comparison. It does not identify the
vanishing-cycle sheaf at each critical point with a constant
sheaf, or construct reduced compact Hall operations.

For point sheaves the cyclic model is also explicit. Fix
$x\in X$ and a volume-normalized tangent basis $e_1,e_2,e_3$,
with $e_3$ in the elliptic direction. The cyclic analytic minimal
model of $\mathcal O_x$ is $\Lambda^\bullet T_xX$ with exterior
multiplication and the volume pairing
\cite[Theorem~4.23 and Proposition~5.18]{VanGarderenCyclic}.
Matrix amplification for $\mathcal O_x^{\oplus n}$ gives
\[
 a=X_1e_1+X_2e_2+Ze_3,\qquad
 \frac13\operatorname{vol}_x\operatorname{Tr}(a^3)
       =\operatorname{Tr}(Z[X_1,X_2]).
\]
Its three critical equations are the pairwise commutators of
$X_1,X_2,Z$. In the formal neighbourhood of the zero triple they
describe the deformations of $\mathcal O_x^{\oplus n}$ modulo
$GL_n$. Over a fixed closed point, support at $x$ requires the
coordinate operators to be jointly nilpotent. The full affine
commuting-triple variety also contains other supports.
Removing $\operatorname{Tr}Z/n$ centres this local model in the
elliptic direction. The gauge terms and relative centres remain
as in \eqref{hall:eq:four-term-gauge-complex}.

\subsubsection*{Finite flags use the same comparison}

We use the suspended bar convention: the bar differential has degree
one and $F$ has degree zero. Its Taylor coefficients below act on
the suspended morphism spaces. Matrix amplification of these
coefficients is the path sum
\[
 (F_n^{\mathrm{mat}}(a_n,\ldots,a_1))_{i_ni_0}
 =\sum_{i_1,\ldots,i_{n-1}}
 F_n((a_n)_{i_ni_{n-1}},\ldots,(a_1)_{i_1i_0}).
\]
An ordered twisted flag is a finite upper-triangular matrix of
objects with a strictly upper-triangular degree-one twisting element
$\delta$ satisfying the Maurer--Cartan equation. Its successive
quotients are the diagonal objects.
Write $\overline\delta=s\delta$, of degree zero in the suspended
space, and regard the $a_i$ below as suspended morphisms.

\begin{theorem}[Cyclic comparison on finite formal flags]
\label{hall:thm:formal-flags}
The chosen $F$ extends to finite twisted flags by
\begin{align*}
 F_*(\overline\delta)
 &=\sum_{n\geq1}F_n^{\mathrm{mat}}
                  (\overline\delta,\ldots,\overline\delta),\\
 F_k^{\mathrm{tw}}(a_k,\ldots,a_1)
 &=\sum_{j_0,\ldots,j_k\geq0}
 F_{k+j_0+\cdots+j_k}^{\mathrm{mat}}
 (\overline\delta_k^{j_k},a_k,\ldots,a_1,\overline\delta_0^{j_0}).
\end{align*}
A power here denotes repeated inputs. These formulas give a cyclic
quasi-equivalence, equivariant under changes of multiplicity bases
and compatible with every regrouping of a finite flag. They also
apply with nilpotent Artinian parameters and in the compatible
inverse-limit formal completion.
\end{theorem}

\begin{proof}
A strictly upper-triangular insertion advances along the finite
flag, so the sums terminate between any two displayed morphisms.
With Artinian parameters, use flag distance and the maximal-ideal
filtration together; each fixed precision still has finitely many
terms. The identity $b_{\mathcal D}F=Fb_A$, evaluated on repeated
$\overline\delta$ and then on distinguished morphism inputs, gives the
Maurer--Cartan equation and all twisted functor equations.

On the associated graded of these finite filtrations, the linear
map is a direct sum of the original quasi-isomorphisms $F_1$.
It is therefore a quasi-isomorphism on each morphism complex.
Lift a flag inductively by lifting its next degree-one extension
class and taking its cone. For a small Artinian extension, the
obstruction lies in degree two; injectivity of $H^2(F_1)$ removes
it, surjectivity of $H^1(F_1)$ matches the extension class, and a
degree-zero gauge correction matches representatives.

The trace is the original cyclic trace followed by matrix trace.
Rotating a closed matrix path preserves it with the Koszul sign.
Regrouping a flag partitions the same paths into consecutive
subpaths. Thus the comparisons preserve the cyclic class and agree
under further subdivision. The formal completion is the inverse
limit of these compatible finite-precision formulas.
\end{proof}

This theorem constructs the formal comparison on all finite flags.
An analytic realization requires convergence. A comparison of
constructible critical coefficients on compact moduli further
requires algebraic or analytic charts, their orientation descent and
the coefficient morphisms on extension correspondences. Those
requirements remain even though the formal category comparison is
available.

\subsubsection*{Critical coefficients on spectral configurations}

The polynomial quiver model also gives coefficients on its actual
spectral configuration spaces. Work in rational constructible
derived categories with smooth descent on quotient stacks, allowing
the unbounded complexes produced by classifying stacks. For this
construction keep the common centre. Put
\[
 \mathcal Y_{\mathbf n}=[\mathcal V_{\mathbf n}/G_{\mathbf n}],
 \qquad d_{\mathbf n}=\dim\mathcal V_{\mathbf n}-\dim G_{\mathbf n}
       =2\sum_{a\in Q_1}n_{t(a)}n_{h(a)},
 \qquad
 \Phi_{\mathbf n}=\varphi_{\mathcal W_{\mathbf n}}
                         \Q_{\mathcal Y_{\mathbf n}}[d_{\mathbf n}].
\]
This normalization uses the local quotient orientation already
specified. In mixed Hodge coefficients include the twist
$(d_{\mathbf n}/2)$, which is integral.

\begin{proposition}[Spectral factorization]
\label{hall:prop:spectral-factorization}
On the open where the spectrum of $Z$ splits into two ordered
disjoint parts, the critical stack and its coefficient are the
products for those two parts. These comparisons are compatible
with every refinement into finitely many disjoint spectral blocks.
\end{proposition}

\begin{proof}
Coprime factors of the characteristic polynomial give polynomial
idempotents in $Z$ by Bezout's identity. They split the vertex
spaces in families on the resultant open. The equations commuting
the arrows with $Z$ make every arrow preserve these subspaces.
The cross-block normal variables have paired potentials
\[
 \operatorname{Tr}(Y\mathcal L(X)),\qquad
 \mathcal L(X)=Z_\beta X-XZ_\alpha.
\]
The eigenvalues of $\mathcal L$ are differences of eigenvalues
in the two blocks. Its determinant is invertible on this open,
so the polarized change $Y'=\mathcal L^*(Y)$ turns the potential
into $\operatorname{Tr}(Y'X)$.

There are
$r(\alpha,\beta)=\sum_{a\in Q_1}
(\alpha_{t(a)}\beta_{h(a)}+\beta_{t(a)}\alpha_{h(a)})$
hyperbolic pairs. Their vanishing-cycle factor is
$\Q[-2r](-r)$ by dimensional reduction. The difference
$d_{\alpha+\beta}-d_\alpha-d_\beta=2r$ cancels this shift,
and the Hodge normalization cancels its Tate twist.
Thom--Sebastiani and stabilization give the coefficient
comparison \cite[Theorems~2.13 and~5.4]{BBDJS2015}.

The chosen polarization fixes its orientation generator: dimensional
reduction identifies it with the complex orientation of the
$X$-space. The change in $Y$ fixes that space. Under refinement,
the same hyperbolic pairs are regrouped; K\"unneth multiplies their
even-degree orientation generators. The resulting normalized
generators have degree zero, so all reassociations agree.
These constructions are equivariant and descend to the quotient.
\end{proof}

Let $\mathfrak M_{\mathbf n}$ be the classical critical stack and
let $\pi_{\mathbf n}$ record the characteristic divisors of its
normal operators in
$\prod_i\operatorname{Sym}^{n_i}\mathbb A^1$. Define
$C_{\mathbf n}=R\pi_{\mathbf n!}\Phi_{\mathbf n}$.
For a word $\boldsymbol\gamma=(\gamma_1,\ldots,\gamma_r)$ of
nonzero dimension vectors, let $\theta_{\boldsymbol\gamma}$
place multiplicity $(\gamma_j)_i$ at the $j$th point in colour $i$.
The word coefficient is
\[
 C_{\boldsymbol\gamma}
 =\theta_{\boldsymbol\gamma}^!C_{\sum_j\gamma_j}
 \quad\text{on }(\mathbb A^1)^r.
\]
Exceptional pullback retains the collision coefficients and their
shifts. The empty coefficient is $\Q$. For a consecutive cut
$A|B$, denote by $j_{A|B}$ the open inclusion where every point
in $A$ differs from every point in $B$. The cut tensor has
$A|B$ component
$Rj_{A|B,*}j_{A|B}^*(C_A\boxtimes C_B)$.

\begin{theorem}[The spectral cut coalgebra]
\label{hall:thm:spectral-cut}
The word coefficients have coassociative maps
\[
 C_{\boldsymbol\gamma}\longrightarrow
 Rj_{A|B,*}j_{A|B}^*(C_A\boxtimes C_B)
\]
that are equivalences after restriction to the cut open. Their
empty coefficient gives the counit and coaugmentation. With this
cut tensor, their cobar is the unital dg algebra
\[
 \Omega_{\mathrm{cut}}C
 =\widehat T_{\mathrm{cut}}(\overline C[-1]),\qquad
 d q_c=-q_{dc}+\sum(-1)^{|c_1|+1}q_{c_1}q_{c_2},
 \quad\overline\Delta c=\sum c_1\otimes c_2.
\]
Here $|q_c|=|c|+1$, and completion is the product over total
dimension vectors.
\end{theorem}

\begin{proof}
Spectral factorization, base change for $R\pi_!$ and K\"unneth
identify the restricted coefficient with the external product.
The unit of the adjunction $j^*\dashv Rj_*$ gives the displayed
map. Two successive cuts use the same three-block factorization
and the same open direct image, proving coassociativity.
Empty cuts give the unit comparisons.

Extend the displayed differential as a derivation. Its linear
and quadratic parts anticommute because $\Delta$ is a chain map;
its two successive splitting terms cancel by coassociativity.
Hence $d^2=0$. At a fixed total dimension $\mathbf n$, a word
of nonempty factors has length at most $\sum_i n_i$. Thus the
completion has only finitely many word lengths in each dimension,
even when its stack cohomology is unbounded.
\end{proof}

This is an associative construction on charged ordered
configurations. Its tensor uses consecutive cuts and their
separation opens. A Hall multiplication instead uses extensions.
Their comparison must preserve these operations on their stated
coefficient categories.

\subsubsection*{The support left by failed cuts}

For a word of length $r$, let $U_i$ be its prefix cut open after
position $i$, and put
$Z_r=(\mathbb A^1)^r\setminus\bigcup_{i=1}^{r-1}U_i$.
Write $R\underline\Gamma_{Z_r}$ for the coefficient functor
of sections supported in this closed set. Expanding the cut
cobar by the subsets of the $r-1$ cuts gives the augmented
\v Cech complex of these opens, shifted by $[-1]$. Therefore
\begin{equation}
\label{hall:eq:cut-support}
 (\overline{\Omega_{\mathrm{cut}}C})_{\boldsymbol\gamma}
 \simeq R\underline\Gamma_{Z_r}(C_{\boldsymbol\gamma})[-1].
\end{equation}
Indeed the term with $k$ cuts has the $k+1$ generator shifts,
and the cobar splitting signs are the alternating restriction
signs. The augmented \v Cech complex is the fibre of restriction
to the union of the cut opens. For two positions this says
\[
 R\underline\Gamma_{Z_2}(C)[-1]
 =\operatorname{Cone}(C\longrightarrow Rj_*j^*C)[-2],
\]
using $\operatorname{Cone}(f)^n=D^n\oplus C^{n+1}$ for
$f:C\to D$.

The first supports already distinguish cuts from the full
diagonal. If $\Delta_{ij}$ denotes equality of positions $i,j$,
then
\[
 Z_2=\Delta_{12},\qquad Z_3=\Delta_{13},\qquad
 Z_4=\Delta_{14}\cup(\Delta_{13}\cap\Delta_{24}).
\]
To verify the last formula, assume positions one and four differ.
Failure of the first and third cuts forces an equality joining
one to an interior position and another joining four to an
interior position. The two interiors cannot be equal. Pairing
one with two and three with four leaves the middle cut open;
pairing one with three and two with four closes all three cuts.
The remaining possibilities force one to equal four. This proves
the formula. In particular, the crossed-pair component must be
retained even away from the total diagonal.

The coefficient on a collision is also larger than its restriction
off that collision. For two one-dimensional colours with $r$
paired arrows, the centred potential is
$f=sQ(x,y)$, where $Q=\sum_{i=1}^r x_i y_i$.
The effective relative group $\mathbb G_m$ acts with weights
$1,-1,0$ on $x,y,s$; the full quotient retains a separate
diagonal scalar stabilizer. Its critical ideal is
\[
 (Q,sx_1,\ldots,sx_r,sy_1,\ldots,sy_r)
       =(x,y)\cap(s,Q).
\]
For the intersection identity, reduce $sa+Qb$ modulo $(x,y)$;
since $s$ is a nonzerodivisor there, membership forces
$a\in(x,y)$. For $r=2$, the invariant matrix $(x_i y_j)$ on
the collision component has trace zero and rank at most one.
Writing it as $\left(\begin{smallmatrix}u&v\\w&-u\end{smallmatrix}\right)$
identifies its affine quotient with $u^2+vw=0$.
On $x\ne0$ the quotient is $T^*\mathbb P^1$.

\begin{proposition}[The collision stalk]
\label{hall:prop:collision-stalk}
For $r\geq2$, the stalk of $\varphi_f\Q[2r+1]$ at the origin
has rank one in degrees $1-2r,-1,0$ and vanishes otherwise.
At a nonzero point of the spectral axis only degree $-1$
remains. At a smooth nonzero point of the collision quadric
only degree $1-2r$ remains.
\end{proposition}

\begin{proof}
The map $(s,x,y)\mapsto(s,x,sy)$ identifies $f^{-1}(1)$ with
$\C^*\times Q^{-1}(1)$. Projection of $Q^{-1}(1)$ to $x\ne0$
is an affine bundle with fibre $\C^{r-1}$. Thus $f^{-1}(1)$
has the homotopy type of $S^1\times S^{2r-1}$.
Homogeneity rescales a small Milnor fibre into a large ball
sublevel of $f^{-1}(1)$. The squared norm is a proper
semialgebraic function on this smooth affine fibre; its critical
set has finitely many connected components, on each of which
its value is constant. Above those values the sublevels are
deformation retracts of the whole fibre. Thus the reduced
Milnor cohomology has rank one in degrees $1,2r-1,2r$.
The shift gives
$H^j(\varphi_f\Q[2r+1])_0=\widetilde H^{j+2r}(F_0,\Q)$.
Off the origin, the transverse functions are respectively the
nondegenerate quadratic form $Q$ and the node $sQ$ with $Q$
as a local coordinate. Their Milnor spheres give the two
remaining assertions.
\end{proof}

These are atlas coefficients; the equivariant quotient also
retains its stabilizer directions. The gauge complex has the
relative infinitesimal action of $\mathbb G_m$ in addition to
the displayed critical
equations.

\subsubsection*{Supported comparison equations}

Fix a dg model of the supported coefficient category and two
uncurved $A_\infty$ algebras $A,B$ in it. Put $sA=A[1]$ and
$sB=B[1]$, and write $b_A,b_B$ for their degree-one bar
coderivations. A comparison has Taylor components $F_n$ of
degree zero satisfying $b_BF=Fb_A$. With the components below
arity $n$ fixed, set
\[
 \mathcal H_n=\operatorname{Hom}_{\mathrm{supp}}
             ((sA)^{\otimes n},sB),\qquad
 R_n=(b_BF_{<n}-F_{<n}b_A)_n.
\]
The tensor and Hom here retain the specified support and diagram
conditions. The Hom differential is
$\partial H=b_{B,1}H-(-1)^{|H|}Hb_{A,1}$, where the source
has its tensor differential.

\begin{proposition}[The next comparison component]
\label{hall:prop:supported-residual}
If all lower arity equations hold, then $\partial R_n=0$.
The next component exists exactly when
$[R_n]=0\in H^1(\mathcal H_n)$, and it solves
$\partial F_n=-R_n$. If
$\partial h+h\partial=1-P$ is a contraction onto closed
cohomology representatives and $PR_n=0$, one solution is
$F_n=-hR_n$. Solutions modulo boundaries, with the lower
components fixed, form a torsor under $H^0(\mathcal H_n)$.
\end{proposition}

\begin{proof}
For $D=b_BF-Fb_A$, the square-zero identities give
$b_BD+Db_A=0$. Its arity-$n$ part is $\partial R_n=0$
because all lower defects vanish. Adding $F_n$ changes the
defect by $\partial F_n$, giving the stated obstruction and
solution. Differences of solutions are degree-zero cycles;
homotopies fixing the lower components add boundaries.
\end{proof}

For the compact comparison, these equations must hold on the
whole collision coefficients. Already the complex
$\C[s]e\to\C[s]$, with $|e|=-1$ and $de=s$, has
$H^0=\C[s]/(s)$ although it becomes contractible after
inverting $s$. A primitive containing $s^{-1}$ therefore does
not solve the equation on the collision. At four inputs,
\eqref{hall:eq:cut-support} also retains the crossed-pair support.

\subsubsection*{Local and horizontal support}

For a coherent sheaf on $X$, call its support \emph{local over $E$}
when its image under $p_E:X\to E$ is finite, and \emph{wrapped}
when that image is all of $E$. Properness and irreducibility of $E$
give these alternatives for a nonempty support. The fibre sheaves
$i_*L_i$ above are local. In a short exact sequence, the middle
support is the union of the two endpoint supports. Hence an
extension is local precisely when both endpoints are local.

\begin{proposition}[Effective elliptic degree]
\label{prop:positive-elliptic-degree-projects}
\label{lem:wrapped-projection-obstruction}
Let $Z=\sum_i m_i[C_i]$ be an effective one-cycle on $X$, with
$m_i>0$. Write $(p_E)_*Z=b[E]$. Then $b>0$ if and only if some
$C_i$ dominates $E$, in which case $p_E(|Z|)=E$.
For a coherent sheaf of support dimension at most one, its
fundamental one-cycle has $b=0$ precisely when its support is local
over $E$.
\end{proposition}

\begin{proof}
Each proper irreducible curve maps either to a point, contributing
zero, or onto $E$, contributing its positive degree. Positivity of
the multiplicities prevents cancellation. Zero-dimensional support
adds only finitely many image points.
\end{proof}

\begin{corollary}
\label{cor:finite-support-ranE-misses-s}
A sheaf with positive effective elliptic degree is not supported
over any finite configuration in $E$. Finite-support sheaf moduli
alone therefore do not contain that geometric sector.
\end{corollary}

This is a statement about support of sheaves. A factorization
algebra on $E$ can carry other labels or act on additional modules;
the proposition does not prohibit such constructions. It also does
not identify $b$ with an exponent of the Igusa product.

\begin{definition}[Geometric degree at a bounded stage]
\label{def:geometric-elliptic-degree-map}
Choose a bounded family of coherent sheaves, a finite set of charges
$\Gamma_R$, and the extension stacks to be used at that stage.
On the sector of support dimension at most one set
$b_R^{\mathrm{geom}}(A)=\deg(p_E)_*\operatorname{cyc}_1(A)$.
On a larger coherent-sheaf sector one may instead fix an ample $H$
on $S$ and an integer $n$ for which every generic-fibre sheaf
$A|_{S_{\C(E)}}(nH)$ in the chosen family is generated by sections
and has no higher cohomology. Set
\[
 b_{R,n}^{\mathrm{geom}}(A)
 =\operatorname{rk}Rp_{E*}(A\otimes p_S^*\mathcal O_S(nH)).
\]
The regularity condition and a common choice of $n$ are part of this
larger-domain construction. We write $b_R^{\mathrm{geom}}$ for the
chosen degree when its domain is fixed.
\end{definition}

The existence of sufficiently large twists is the relative Serre
and global-generation statement; a common bound must hold on the
chosen families \cite[Tags~0B5S and~01PR]{StacksProject}. The degree
is additive on their Grothendieck classes. On the larger domain it
is the generic-fibre Euler characteristic, which by the regularity
hypothesis equals the dimension of global sections. It is
nonnegative, and it is zero precisely when that generic-fibre sheaf
is zero. Since the sheaf is generated by its sections, this is
precisely the local-support condition. On the one-dimensional
support sector the generic fibre has finite length, so the two
definitions agree and do not depend on $n$.

\begin{definition}[Additive support data]
\label{def:geometric-elliptic-degree-additivity-datum}
The charges of every permitted short exact sequence must satisfy
$\gamma_C=\gamma_A+\gamma_B$. A degree on these charges is
additive support data when it represents the chosen geometric
degree of each object and satisfies
\[
 b_R^{\mathrm{geom}}(\gamma_A+\gamma_B)
 =b_R^{\mathrm{geom}}(\gamma_A)+b_R^{\mathrm{geom}}(\gamma_B),
 \qquad b_R^{\mathrm{geom}}(0)=0.
\]
Only operations whose endpoint, intermediate and total charges are
present are included at a bounded stage.
\end{definition}

\begin{proposition}[Support and addition]
\label{prop:geometric-elliptic-degree-additivity-criterion}
Both constructions in Definition~\ref{def:geometric-elliptic-degree-map}
are additive on the stated coherent-sheaf domains. Consequently the
support types obey
$\mathrm L\mathrm L=\mathrm L$, whereas a product involving a
wrapped input has type $\mathrm W$.
\end{proposition}

\begin{proof}
Fundamental cycles are additive for short exact sequences of
sheaves of dimension at most one. For the second construction,
derived pushforward is additive in the Grothendieck group and rank
is additive. The support characterization just proved gives the
type rule. The direct support-union argument gives the same rule
without choosing a numerical degree.
\end{proof}

The regularity hypothesis matters. For a line bundle $L_i$ and
the polarization $H$ of the common-polynomial triple,
$\chi(S,L_i(nH))=4n^2-2$. Thus the horizontal sheaf
$L_i\boxtimes\mathcal O_E$ has that generic Euler characteristic,
which is negative at $n=0$. It becomes a nonnegative rank only on
the acyclic, generated range. A virtual perfect-complex class can
also have cancellations, so effectiveness must not be inferred
from its numerical degree alone.

\subsubsection*{The configuration and family spaces}

For a finite set $I$, a labelled configuration is a tuple
$\mathbf e\in E^I$. Let $Z_I\subset E\times E^I$ be the union of
the graphs of its coordinates. A closed local incidence stack at
fixed thickness $N$ parametrizes a flat family $A$ with
\[
 \mathcal I_{S\times Z_I}^{N}A=0.
\]
We denote it by $\mathfrak M^{\mathrm{loc,cl}}_{\gamma,R,I,N}$.
Keeping $N$ retains infinitesimal support at collisions. Set-theoretic
support in $Z_I$ does not impose thickness one. A bounded choice of
$N$, or the specified filtered system in $N$, is included whenever
these incidence stacks are used.

Thickness also enters extension closure. If the same ideal kills
the two endpoints to powers $N_A$ and $N_B$, it kills the middle
sheaf to power $N_A+N_B$: its $N_A$-th power sends the middle
sheaf into the subobject, and the next $N_B$ powers kill that
subobject. Thus collision and extension maps must use compatible
thickness bounds, rather than assume that one fixed bound is
closed under all extensions.

\begin{definition}[Wrapped families and scalar rigidification]
\label{def:e-equivariant-wrapped-prequotient-stack}
Let $\mathcal M^{\mathrm{wr}}_{\eta,R}$ be a chosen algebraic
stack of wrapped sheaves. When its scalar subgroup can be
rigidified, write
\[
 \rho_{\eta,R}:\mathcal M^{\mathrm{wr}}_{\eta,R}
 \longrightarrow\mathcal M^{\mathrm{wr,rig}}_{\eta,R}.
\]
The map removes the central scalar automorphisms. It is directed
from the original stack to the rigidified stack. Elliptic
translation has not been divided out. A coefficient or universal
object on the rigidification requires descent along $\rho$;
the universal sheaf on the original stack is not assumed to descend.
\end{definition}

\begin{lemma}[Translation on families]
\label{lem:e-action-law-wrapped-prequotient}
Support translation acts by
$a\cdot A=(1_S\times\tau_a)_*A$, where $\tau_a(e)=e+a$,
and gives a coherent $E$-action
on the sheaf and flag stacks. It descends through a scalar
rigidification compatible with this central subgroup.
\end{lemma}

\begin{proof}
The equality $\tau_a\tau_b=\tau_{a+b}$ gives the composition
isomorphism on sheaves and all their morphisms. Its associativity
is the associativity of composition of translations. Scalar
multiplication commutes with this action, so the universal property
of rigidification gives the descended action.
\end{proof}

\begin{proposition}[Finite type after rigidification]
\label{prop:wrapped-prequotient-finite-type}
Suppose the chosen wrapped stack is finite type and its scalar
rigidification is an algebraic $\mathbb G_m$-gerbe. Then the
rigidified stack is finite type. No universal sheaf on it is implied.
In particular, the finite-type premise holds for a finite open
cover by specified finite-type semistable substacks. Finiteness of the
numerical index set and boundedness of each substack are separate
premises of this construction.
\end{proposition}

\begin{proof}
The gerbe is smooth and surjective, and finite type over the base
can be checked on such a cover. Descent of a particular sheaf is a
separate condition on its scalar action and descent isomorphisms.
The finite disjoint union of the open substacks is a smooth
surjective finite-type cover, proving the additional assertion.
\end{proof}

\begin{definition}[A carrier for local and wrapped families]
\label{def:hybrid-ran-prestack}
\label{def:hybrid-ran-base}
Choose the charge and moduli domains above. A presentation of the
hybrid carrier consists of a finite labelled local configuration
$(I,\alpha,\mathbf e)$ and a finite list of wrapped families
$(W_j)_{j\in J}$ with charges $\eta_j$. Form the prestack colimit
of
\[
 E^I\times\prod_{j\in J}\mathcal M^{\mathrm{wr}}_{\eta_j,R}
\]
over relabellings and local collision maps. Under a surjection
$I\twoheadrightarrow I'$, the diagonal repeats coordinates and
the new local colour is the sum of the colours in each fibre.
Only defined charge sums occur. Denote this prestack by
$\operatorname{Ran}^{\mathrm{hyb}}_{R,\mathrm{pre}}(E)$.
Compatible scalar rigidifications can be used in this diagram when
the required family and coefficient descent is supplied.
\end{definition}

The indexing category contains all specified finite presentations;
it is not generally a finite category. A topology $\tau$ gives its
stackification $a_\tau\operatorname{Ran}^{\mathrm{hyb}}_{R,\mathrm{pre}}(E)$.
An operation or coefficient object on that stackification must
satisfy descent. Neither taking the colimit nor stackifying defines
a Hall coefficient system.

\begin{definition}[Local and wrapped subcarriers]
\label{def:hybrid-local-stratum-prestack}
\label{def:hybrid-wrapped-stratum-prestack}
The local subcarrier has $J=\varnothing$; the purely wrapped
subcarrier has $I=\varnothing$. Mixed presentations have both.
When the geometric degree is defined on the chosen domain, its
charge types are
\begin{equation}
\label{eq:hybrid-loc-wr-windows}
 \Gamma_R^{\mathrm{loc}}=\{\gamma:b_R^{\mathrm{geom}}(\gamma)=0\},
 \qquad
 \Gamma_R^{\mathrm{wr}}=\{\gamma:b_R^{\mathrm{geom}}(\gamma)>0\}.
\end{equation}
The local incidence stacks and the critical sheaves placed over
these carriers are additional objects, not part of their points.
\end{definition}

A wrapped family is retained as a family, with its automorphisms.
Replacing it by its degree or by a single point of $E$ would discard
that information. The ordered extension diagrams below retain still
more information: which object is a subobject and which is a quotient.

\subsubsection*{Ordered extensions and proper targets}

We use quotient-first input order. Thus
$\mathfrak E_{\alpha,\beta}$ parametrizes short exact sequences
\[
 0\longrightarrow B_\beta\longrightarrow C_{\alpha+\beta}
 \longrightarrow A_\alpha\longrightarrow0,
\]
with endpoint map $p(C)=(A,B)$ and middle-object map $q(C)=C$.
All sheaves, maps and configuration parameters belong to the chosen
domains. The charge set at height $R$ is only a window: a nonzero
positive charge cannot lie in a finite set closed under all of its
positive multiples. Operations at this stage are therefore partial,
unless a separate quotient construction makes them total.

\begin{definition}[The four binary correspondence types]
\label{def:local-local-correspondence-datum}
\label{def:mixed-local-wrapped-correspondence-datum}
\label{def:wrapped-wrapped-correspondence-datum}
Restrict the ordered extension stack to the indicated support types:
\[
 \begin{array}{c|c|c}
 \text{type}&\text{endpoints }(A,B)&\text{middle object}\\ \hline
 \mathrm{LL}&(A_{\mathrm L},B_{\mathrm L})&C_{\mathrm L}\\
 \mathrm{LW}&(A_{\mathrm L},B_{\mathrm W})&C_{\mathrm W}\\
 \mathrm{WL}&(A_{\mathrm W},B_{\mathrm L})&C_{\mathrm W}\\
 \mathrm{WW}&(A_{\mathrm W},B_{\mathrm W})&C_{\mathrm W}.
 \end{array}
\]
Local endpoints retain their closed configuration presentations and
thicknesses. Wrapped endpoints retain their sheaf families. These
are correspondences before elliptic quotient. Scalar rigidification
is applied to a compatible whole diagram when such a diagram and
its coefficient descent have been specified.
\end{definition}

\begin{proposition}[A proper coherent-sheaf domain]
\label{prop:local-local-target-properness}
Let $C$ be a flat family of coherent sheaves on the projective
threefold $X$, over a base $T$. Fix the Hilbert polynomial of the
quotient $A$. If all additional endpoint and support conditions
cut out a closed subfunctor of $\operatorname{Quot}_T(C)$, then
the corresponding middle-object map $q$ is proper over $T$.
\end{proposition}

\begin{proof}
The relative Quot scheme with that polynomial is projective over
$T$ \cite[Theorem~2.2.4]{HuybrechtsLehn}. Its universal quotient
is flat, and its kernel is flat because
$C$ is flat. The required sequences therefore form the stated
closed subfunctor of a projective $T$-scheme, which is proper.
Fixed annihilator conditions defining the local incidence spaces
are imposed inside this Quot problem.
\end{proof}

\begin{proposition}[Mixed and wrapped target maps]
\label{prop:mixed-local-wrapped-target-admissibility}
\label{prop:wrapped-wrapped-target-admissibility}
The same properness criterion holds for the mixed and wrapped
coherent-sheaf correspondences, and for their finite flag versions.
If instead the chosen conditions define only an open substack, the
restriction of $q$ need not be proper. A product with ordinary
critical cohomology is then available by this route only when the
specific coefficient complex has proper support over the target,
or another appropriate integration morphism has been constructed.
\end{proposition}

\begin{proof}
Use the same relative Quot scheme, or its projective iterated flag
version, and impose the closed conditions. An open subspace of a
proper scheme need not be proper: $\mathbb A^1\subset\mathbb P^1$
is an example. Thus mere openness of a stability condition gives
no properness conclusion. The support alternative is exactly the
condition identifying the relevant exceptional and ordinary
pushforwards.
\end{proof}

In particular, ordinary Quot properness does not establish this
criterion for all perfect complexes in a tilted heart. Such a heart
needs its actual extension moduli and the corresponding properness
or integration theorem. A bound on the number of numerical charges
also does not prove that their moduli are finite type.

\subsubsection*{Flags and their eight support words}

A three-step ordered flag is written in descending order,
\[
 C=F^0\supset F^1\supset F^2\supset F^3=0,
 \qquad A_i=F^{i-1}/F^i.
\]
This convention matches quotient-first multiplication.
Every intermediate sheaf and its configuration belongs to the
specified domain. For a partial charge window, both regroupings
are compared only when all their intermediate charges are present.

\begin{definition}[Binary support words]
\label{def:eight-local-wrapped-binary-words}
The support word of this flag is
$(\epsilon_1,\epsilon_2,\epsilon_3)$, with
$\epsilon_i=\mathrm L$ or $\mathrm W$ according to $A_i$.
The eight possibilities are
\[
 \mathrm{LLL},\quad\mathrm{LLW},\quad\mathrm{LWL},\quad
 \mathrm{WLL},\quad\mathrm{LWW},\quad\mathrm{WLW},\quad
 \mathrm{WWL},\quad\mathrm{WWW}.
\]
Each block has type $\mathrm L$ exactly when all its factors do.
The word specifies filtration order, not physical time order.
\end{definition}

\begin{proposition}[The two extension compositions]
\label{prop:eight-two-step-flag-stacks}
For each of the eight words, the two iterated extension
correspondences are naturally equivalent to the stack of flags
with those successive quotients. These equivalences commute with
base change in the parameter scheme.
\end{proposition}

\begin{proof}
For the first composition, take a sequence with quotient $C_{12}$
and subobject $A_3$, together with a sequence expressing $C_{12}$
as an extension of $A_1$ by $A_2$. The inverse image of $A_2$
defines $F^1$, and $F^2=A_3$. For the other composition, take
$F^1$ as the extension of $A_2$ by $A_3$ and then extend $A_1$
by $F^1$. These constructions are inverse on families and their
isomorphisms. Kernels, inverse images and the indicated quotients
commute with base change for these flat exact sequences. The
support-union rule gives precisely the displayed word types.
\end{proof}

This is a geometric comparison of correspondences. To deduce an
equality of operations, the two coefficient morphisms on that flag
stack must agree. We state the needed coefficient framework before
making that deduction.

\subsubsection*{Coefficient maps and six operations}

Fix a characteristic-zero coefficient field $k$; the rational
constructions above use $k=\Q$.

\begin{definition}[An admissible coefficient theory]
\label{def:compact-support-exceptional-pushforward-model}
Choose constructible derived categories $D_c(Y)$ on all spaces and
stacks in the correspondence diagram. Require the six operations
$f^*,Rf_*,Rf_!,f^!$, tensor product and internal Hom on the maps
being used, with coherent composition, the needed Cartesian
base-change isomorphisms and the projection formula. Each chosen
coefficient and its image must remain in the stated categories.
On quotient stacks this includes the full smooth or equivariant
descent diagrams. Boundedness of cohomology or perfection is an
additional condition whenever a finite dual or ordinary tensor
product requires it.
For a compactification $f=\overline f j$, with $j$ open and
$\overline f$ proper, the exceptional image is
$Rf_!K=R\overline f_*j_!K$. If $f$ is proper, take $j=1$
and $\overline f=f$. Independence of this presentation is
proved in Proposition~\ref{prop:finite-hall-product-qbang-compactification-independence}.
\end{definition}

Write $\mathcal M_\gamma$ for the chosen object stack and
choose $K_\gamma\in D_c(\mathcal M_\gamma)$.
For ordinary critical cohomology, a binary operation uses an
integer source shift $\kappa(\alpha,\beta)$ and a coefficient morphism
\begin{equation}
\label{hall:eq:binary-coefficient-map}
 u_{\alpha,\beta}:
 p^*(K_\alpha\boxtimes K_\beta)[\kappa(\alpha,\beta)]
 \longrightarrow q^!K_{\alpha+\beta},
\end{equation}
together with properness of the source coefficient's support
over the target of $q$. This identifies its ordinary and
exceptional pushforwards. Pullback on ordinary cohomology,
then $Rq_!u$ and the trace, gives the product.
A Thom--Sebastiani isomorphism on a
product of critical charts contributes to $u$ only after its
comparison with the extension diagram has been constructed.
Appendix~\ref{app:hall-product} gives a theorem that constructs
these morphisms on its unreduced moduli domain.
More explicitly, a comparison
$p^*(K_\alpha\boxtimes K_\beta)[\kappa]\simeq K_{\mathfrak E}$
can be used when it is supplied together with
$K_{\mathfrak E}\to q^!K_{\alpha+\beta}$. Their composite is
$u_{\alpha,\beta}$. The first isomorphism alone does not supply
the second arrow.

\begin{proposition}[Base change and projection formula]
\label{prop:mixed-correspondence-base-change}
\label{prop:mixed-correspondence-projection-formula}
In the chosen admissible coefficient theory, a Cartesian square
with horizontal maps $f,f'$ and vertical maps $q,q'$ gives
\[
 f^*Rq_!K\simeq Rq'_!f'^*K,
 \qquad
 Rq_!(K\otimes q^*L)\simeq Rq_!K\otimes L.
\]
Consequently the pull--push functors of composable extension
correspondences identify with the functor of their fibre-product
correspondence. The coefficient morphisms themselves must be
compatible with these identifications.
\end{proposition}

\begin{proof}
Apply the base-change and projection-formula isomorphisms in the
definition, then the composition isomorphism for exceptional
pushforward. This compares the functors on their common domain;
it makes no assertion that two independently chosen coefficient
morphisms are equal.
\end{proof}

\begin{proposition}[Critical coefficient transport]
\label{prop:mixed-correspondence-thom-sebastiani}
Suppose oriented critical charts and their coefficient morphisms
are compatible on the extension and flag diagrams. Their
Thom--Sebastiani and hyperbolic comparison maps then define the
morphisms \eqref{hall:eq:binary-coefficient-map} compatibly with
base change and projection formula. This implication holds for
local, mixed and wrapped correspondences on that domain.
\end{proposition}

\begin{proof}
The chartwise morphisms are objects of the descent diagram for
morphisms of constructible complexes. Their specified overlap
homotopies glue them. Naturality of the six operations transports
the base-change and projection-formula comparisons through this
descent. The critical Hall theorem supplies one such coherent
system under its own hypotheses; an arbitrary orientation line
alone does not supply the morphisms.
\end{proof}

\begin{proposition}[Associativity for the eight words]
\label{prop:lll-associativity}
\label{prop:llw-associativity}
\label{prop:lwl-associativity}
\label{prop:wll-associativity}
\label{prop:lww-associativity}
\label{prop:wlw-associativity}
\label{prop:wwl-associativity}
\label{prop:www-associativity}
Assume that the two composites of the maps $u$ agree on each
three-step flag stack, including their shifts and orientation
identifications. Then every defined triple product satisfies
\[
 m_{\alpha+\beta,\gamma}(m_{\alpha,\beta}\otimes1)
 =m_{\alpha,\beta+\gamma}(1\otimes m_{\beta,\gamma})
\]
with the tensor associativity constraint determined by $\kappa$.
\end{proposition}

\begin{proof}
Proposition~\ref{prop:eight-two-step-flag-stacks} identifies the
two geometric compositions with one flag stack. Base change and
projection formula identify the induced pull--push functors there.
The assumed equality of coefficient morphisms makes their traces
equal. This proves the displayed equation for all eight words.
\end{proof}

\begin{proposition}[Four-input coherence]
\label{prop:four-input-pentagon-coherence}
The geometric associators obtained by forgetting steps of a
four-step flag satisfy the pentagon. If the coefficient
comparisons are coherent on that flag diagram, their associators
satisfy the same pentagon.
Writing $a_{\alpha,\beta,\gamma}$ for reassociation from left
to right, the identity is
\[
 a_{\alpha,\beta,\gamma+\delta}
 a_{\alpha+\beta,\gamma,\delta}
 =(1\otimes a_{\beta,\gamma,\delta})
   a_{\alpha,\beta+\gamma,\delta}
   (a_{\alpha,\beta,\gamma}\otimes1).
\]
\end{proposition}

\begin{proof}
Both paths around the pentagon identify a filtration with the
same successive quotients and forget the same intermediate steps.
The inverse-image and quotient constructions in the preceding
proof agree under repeated subdivision. Applying a coherent
coefficient functor preserves this equality.
\end{proof}

A pentagon for categorical associators controls reassociation in
that category. It does not by itself provide all higher chain
homotopies of an $A_\infty$ operation or all collision coherences
of a factorization object.

\begin{definition}[The coloured flag diagram]
\label{def:higher-coloured-hybrid-tree-category}
An ordered profile consists of finitely many local or wrapped
inputs with their configuration and family data. A planar rooted
tree records an iterated extension of that profile; each edge
carries the sum of the charges below it and the corresponding
support type. Vertices carry the indicated extension or flag
correspondence. Grafting gives composition, and contracting an
internal edge regroups a flag.
A coherent coefficient system on this diagram specifies the
operations and comparison homotopies on all finite trees, together
with units, local relabellings, configuration refinements and
overlap descent. Every requested coherence is on its actual
configuration space, including collisions.
The notation $\mathfrak o^{\mathrm{col}}_R=0$ abbreviates
these tree, unit, symmetry, refinement, descent and overlap
conditions. Each condition belongs to its own diagram; the
notation does not assert that they form one cohomology class.
\end{definition}

\begin{proposition}[From coherent flags to operations]
\label{prop:higher-coloured-residual-vanishing-criterion}
A coherent coefficient system on the coloured flag diagram
induces the corresponding associative or homotopy-coherent
operations, with its specified units and local symmetries.
The conclusion is restricted to the charge and support domain of
that system.
\end{proposition}

\begin{proof}
Compose the coefficient morphisms and their pull--push functors
along each tree. Edge contraction gives the comparison between
composites. The prescribed higher compatibilities are exactly the
relations among these comparisons. The six-functor composition
constraints transport them to the resulting operations.
\end{proof}

\subsubsection*{Units, disjoint support and insertion order}

\begin{definition}[The zero-object correspondence]
\label{def:hybrid-unit-object-vacuum-correspondences}
Before elliptic quotient, the distinguished zero object has
moduli stack $\mathrm{pt}$, empty support, and charge zero.
Its two unit correspondences are
\[
 0\longrightarrow A\xrightarrow{1}A\longrightarrow0\longrightarrow0,
 \qquad
 0\longrightarrow0\longrightarrow A\xrightarrow{1}A\longrightarrow0.
\]
Choose the coefficient at zero and the comparison morphisms so
that these correspondences induce the identity. Scalar
rigidification of nonzero objects does not rigidify the zero
object by a nonexistent scalar subgroup of its automorphisms.
\end{definition}

\begin{proposition}[Local and wrapped units]
\label{prop:local-unit-compatibility}
\label{prop:wrapped-unit-compatibility}
With these unit coefficient identifications, both zero-endpoint
operations are the identity on every local or wrapped coefficient
object.
\end{proposition}

\begin{proof}
Each displayed sequence is uniquely determined by $A$, including
its morphisms. Both legs identify with the identity map on its
moduli stack. The coefficient comparison and the trace are then
the stipulated identity maps.
\end{proof}

\begin{proposition}[Local symmetry and configuration descent]
\label{prop:local-configuration-symmetric-descent}
Sheaves with disjoint supports over $E$ have vanishing mutual Ext
groups. Their extension correspondence is the direct-sum
correspondence. If its oriented coefficient comparison is the
external-product comparison, the local operation on disjoint
configurations is symmetric with the coefficient category's
Koszul signs. A coefficient object on the unlabelled carrier
requires compatible equivariant descent under relabellings and
under the specified collision refinements.
For a colour-preserving permutation, this is equivariant descent
under
$\mathfrak S_\alpha=\{\sigma\in\operatorname{Aut}(I):
\alpha_{\sigma(i)}=\alpha_i\}$.
The corresponding quotient is
$[\mathfrak M^{\mathrm{loc,cl}}_{\alpha,R,I,N}/\mathfrak S_\alpha]$;
its coefficient includes the full equivariant identifications.
\end{proposition}

\begin{proof}
All local Ext sheaves between disjoint closed supports vanish,
so their global Ext groups vanish. The extensions split uniquely
relative to their endpoints. Permuting the summands is compatible
with their product charts and external tensor products, giving
the stated symmetry. On labelled charts the descent is the full
equivariant diagram; collision maps require its additional
refinement comparisons.
Explicitly, a relabelling $\sigma:I\to I'$ sends
$(\alpha_i,e_i)$ to
$(\alpha_{\sigma^{-1}(i')},e_{\sigma^{-1}(i')})$.
It identifies the graph unions $Z_I,Z_{I'}$, preserves the
annihilator condition at thickness $N$, and leaves the sheaf
itself unchanged. Its induced coefficient isomorphisms must
satisfy the permutation composition law before quotient descent.
\end{proof}

In particular, summing independently over every $E^I$ would count
the same collision configurations repeatedly. A coefficient system
on the carrier uses the actual colimit and its compatible maps,
or an equivalent descent construction, rather than an unrelated
direct sum over all finite labels.
On a permitted collision, the one-label presentation and the
two-coincident-label presentation are identified by descent.
A direct sum over their label sets keeps two independent copies.

\begin{proposition}[Ordered wrapped insertions]
\label{prop:wrapped-insertion-order-conventions}
For overlapping supports, exchanging two inputs changes which
one is the subobject. The corresponding extension stacks need
not be identified. Planar order must therefore be retained unless
a specified symmetry or braiding compares the two operations.
\end{proposition}

\begin{proof}
The exchanged correspondence parametrizes
$0\to A\to C'\to B\to0$, whereas the original parametrizes
$0\to B\to C\to A\to0$. They involve opposite Ext groups and
have no natural identification from the definition of an
extension. The disjoint-support splitting argument does not
apply when those groups are nonzero.
\end{proof}

\begin{definition}[Two-sided mixed insertions]
\label{def:two-sided-mixed-correspondence-datum}
A two-sided mixed correspondence is the flag stack with ordered
successive quotients $(A_-,W,A_+)$, where the first and last
objects have local support and $W$ is wrapped. Collections of
local inputs may replace $A_-$ and $A_+$, with their original
orders and configurations retained. Forgetting the left or right
steps gives the one-sided mixed correspondences.
\end{definition}


\subsubsection*{Elliptic quotient over a common torsor}

Let $BE$ denote the stack of $E$-torsors, and let $Y$ be a stack
with an $E$-action. A point of $[Y/E]$ over
$T$ is an $E$-torsor $P\to T$ together with an equivariant map
$P\to Y$. The map $[Y/E]\to BE$ remembers $P$.
For the unrigidified sheaf stacks put
$\mathcal Q_\gamma=[\mathcal M_\gamma/E]$ and
$\mathcal H_{\alpha,\beta}=[\mathfrak E_{\alpha,\beta}/E]$,
using simultaneous translation of the exact sequence.

\begin{proposition}[Common-torsor Hall spans]
\label{hall:prop:common-torsor}
There is a natural equivalence
\[
 \mathcal Q_\alpha\times_{BE}\mathcal Q_\beta
 \simeq[(\mathcal M_\alpha\times\mathcal M_\beta)/E_{\mathrm{diag}}].
\]
The extension diagram descends to a relative span
\[
 \mathcal Q_\alpha\times_{BE}\mathcal Q_\beta
 \longleftarrow\mathcal H_{\alpha,\beta}
 \longrightarrow\mathcal Q_{\alpha+\beta}.
\]
Its iterated compositions are identified by the same flag
constructions, and its relative zero-object unit is $BE$.
\end{proposition}

\begin{proof}
An object of the fibre product consists of two torsors, their
families, and an isomorphism between the torsors. Transport one
family across this isomorphism. Both families now live on one
torsor; this is an object of the diagonal quotient. Repeating that
torsor in the two factors gives the inverse equivalence on objects
and arrows, compatibly with base change.

An equivariant extension on the common torsor has equivariant
endpoints and middle object. Thus the extension span descends.
The inverse-image and quotient constructions for flags are
translation equivariant, so they identify its two compositions.
Regroupings of longer flags forget the same steps. The zero
object is fixed by translation and has quotient
$[\mathrm{pt}/E]=BE$. Extensions with zero give the relative
unit identities.
\end{proof}

\begin{definition}[Quotients of correspondence diagrams]
\label{def:quotient-after-correspondence-pseudofunctor}
For an $E$-equivariant correspondence
$Y_0\leftarrow Z\to Y_1$, assign the span
$[Y_0/E]\leftarrow[Z/E]\to[Y_1/E]$ over $BE$.
Composition uses relative fibre products over the middle quotient.
Denote this assignment on the chosen diagrams by
$\mathcal Q^{\mathrm{corr}}_{E,R}$.
It is a quotient of diagrams, not a cosection reduction of
obstruction complexes.
For a coefficient realization of this assignment,
$\mathfrak o^Q_R=0$ will mean compatibility with composition,
base change, the critical coefficient maps, flag refinement,
the chosen cohomology realization, operation comparison and
transitions. Its components are defined only after the maps
in question have been supplied.
\end{definition}

\begin{proposition}[Composition and Cartesian squares]
\label{prop:quotient-after-correspondence-preserves-composition}
\label{prop:quotient-after-correspondence-preserves-base-change}
For equivariant maps $Z\to Y\leftarrow W$,
\[
 [(Z\times_Y W)/E]
 \simeq[Z/E]\times_{[Y/E]}[W/E].
\]
These equivalences preserve the composition and base-change
squares of the correspondence diagrams and all their flag
refinements.
\end{proposition}

\begin{proof}
A point on the right gives a common torsor after transporting
across its specified isomorphism. The equality of the two maps
to $Y$ is then exactly an equivariant point of $Z\times_Y W$.
This construction and its inverse commute with pullback and with
forgetting a stage of a longer diagram. The resulting
associativity identifications are therefore coherent.
\end{proof}

\begin{proposition}[What independent quotients forget]
\label{prop:quotient-after-correspondence-excludes-quotient-first}
Replacing $\mathcal Q_\alpha\times_{BE}\mathcal Q_\beta$ by
$\mathcal Q_\alpha\times\mathcal Q_\beta$ discards the
identification of their torsors. It is not generally an equivalent
input stack. A construction directly on quotient stacks is valid
when it retains this relative input and its coefficient maps.
\end{proposition}

\begin{proof}
Take two copies of $E$ with the translation action. Their diagonal
quotient is $E$, represented by the difference of the two
coordinates. Each independent quotient $[E/E]$ is a point, so
the product of independent quotients is a point. The difference
coordinate is lost. The common-torsor fibre product retains it
by Proposition~\ref{hall:prop:common-torsor}.
\end{proof}

Thus the obstruction is to forgetting the common torsor, not to
working with quotient stacks. The relative construction already
gives the correct quotient correspondence. Obtaining a reduced
critical operation on it still requires reduced coefficients and
their comparison maps.

\subsubsection*{Coefficient descent and its variance}

For an $E$-stack $Y$, a descended coefficient consists of
$\overline K_Y\in D_c([Y/E])$ and its coherent pullback
identification with $K_Y$ on $Y$. Equivalently it is an object
of the full descent diagram on the action groupoid
$E^\bullet\times Y$. This includes the higher compatibility
maps; an isomorphism on $E\times Y$ without its cocycle data is
insufficient. If critical coefficients are normalized perversely,
the smooth-atlas dimension shifts are included in the stated
identification.

\begin{proposition}[Descent of the critical comparisons]
\label{prop:quotient-after-correspondence-preserves-thom-sebastiani}
Suppose all coefficient objects, binary morphisms and their flag
homotopies form a compatible equivariant descent diagram.
Then they descend to the quotient correspondence diagram, with
its relative input over $BE$. Their base-change, projection and
Thom--Sebastiani comparisons descend with them.
\end{proposition}

\begin{proof}
The constructible category on the quotient is the totalization
of its equivariant descent categories. The supplied morphisms and
homotopies define a morphism in this totalization. The functorial
comparison diagrams are equivariant by hypothesis, so they define
the corresponding diagrams after totalization. This uses the
full descent data; forgetting to the underlying complex on $Y$
is not a fully faithful test for all equivariant derived morphisms.
\end{proof}

The translation quotient is not automatically within the
critical Hall theorem's moduli hypotheses. For example, its
relative zero has stabilizer $E$, which is not affine, whereas
Theorem~\ref{thm:hall-product} assumes affine stabilizers.
Moreover, removing an orbit and its dual obstruction direction
requires the reduced critical construction, not just the quotient
of the classical stack.

Fix a Verdier duality functor $\mathbb D_Y$ in the chosen
coefficient theory. A useful dual compact-support convention is
\[
 \mathbb B(Y;K):=R\operatorname{Hom}(R\Gamma_c(Y,K),k).
\]
When duality applies it is $R\Gamma(Y,\mathbb D_YK)$.
For $K=k_Y$, this is the usual dualizing-complex description of
Borel--Moore chains, with its indicated cohomological indexing.
Crucially, it is contravariant in $K$.

\begin{proposition}[Dual coefficients on the quotient]
\label{prop:quotient-after-correspondence-borel-moore-chain-functor}
When the indicated functors exist, the descended coefficient gives
$\mathbb B([Y/E];\overline K_Y)$. A morphism $f:K\to L$ induces
\[
 \mathbb B(Y;L)\longrightarrow\mathbb B(Y;K),
\]
and duality reverses composition. An operation obtained by dualizing a
compact-support operation must therefore reverse its coefficient
arrows. Identification of dual tensor products additionally
requires the stated finiteness or completed-tensor hypothesis.
\end{proposition}

\begin{proof}
Apply $R\operatorname{Hom}(-,k)$ to the morphism
$R\Gamma_c(Y,K)\to R\Gamma_c(Y,L)$.
For $K\xrightarrow fL\xrightarrow gM$, the resulting map is
$f^\vee g^\vee=(gf)^\vee$. The tensor-duality identification is
valid for perfect complexes; outside that domain its completion
or replacement must be specified. Descent gives the quotient
coefficient object before this duality is taken.
\end{proof}

This does not produce a natural map from nonequivariant chains on
$Y$ to chains on $[Y/E]$, nor an equivalence between them. Their
coefficient and stabilizer information differs. In particular,
finite type of a stack does not imply finite-dimensional
cohomology: already $H^*(B\mathbb G_m,\Q)=\Q[c]$ is unbounded.

\begin{proposition}[Comparison of operations]
\label{prop:quotient-after-correspondence-theta-mu-comparison}
\label{prop:quotient-after-correspondence-theta-mu-associativity}
\label{prop:quotient-after-correspondence-theta-mu-coproduct}
Suppose a coefficient realization $Q_{E,R}$ of the equivariant
correspondence diagram is supplied, together with natural
comparisons
\[
 \theta_{\mu,R}:Q_{E,R}\mu_R
       \xrightarrow{\ \sim\ }\overline\mu_RQ_{E,R}^{\otimes n}
\]
for its $n$-input operations. The tensor on the quotient side is
the specified relative tensor. If these comparisons commute with
flag refinement, they preserve associativity. If they also
commute with the separate collision diagrams, they preserve the
specified coproduct and its compatibilities.
\end{proposition}

\begin{proof}
Use the comparison to replace each operation in either composite.
Naturality and the assumed flag compatibility identify the two
replacements with the comparison for the common longer flag.
The same argument on the collision diagram proves the coproduct
assertion. A coproduct comparison is used only after that
coproduct has been defined.
\end{proof}

\begin{proposition}[Primitives and transitions under comparison]
\label{prop:quotient-after-correspondence-theta-mu-primitives}
\label{prop:quotient-after-correspondence-hn-transition-compatibility}
An operation-preserving comparison of coaugmented counital
coalgebras maps primitive elements to primitive elements. To
identify their primitive objects, an equivalence or a suitable
kernel-preservation property is required. Comparisons at two
bounded stages commute with transition only when the coefficient
maps and their natural transformations satisfy that transition
square.
\end{proposition}

\begin{proof}
Apply the comparison to
$\Delta x=1\otimes x+x\otimes1$ and use preservation of the
coaugmentation and coproduct. This proves preservation of
primitive elements. Equality with the whole target kernel needs
more than an inclusion, hence the stated additional hypothesis.
The transition assertion follows by composing the supplied
commuting squares. None is implied by the existence of the maps
on isolated stages.
\end{proof}

\subsubsection*{Anchors retain less than families}

An anchor is auxiliary data for recording a family's position or
spectral invariants. It is not needed to define the common-torsor
quotient. For a perfect family $A$, normalize its determinant
pushforward to degree zero:
\[
 \lambda^{\det}(A)=\det Rp_{E*}A\otimes\mathcal O_E(-\chi(A)0_E)
 \quad\text{in }\operatorname{Pic}^0(E)\simeq E.
\]
Proposition~\ref{prop:determinant-anchor-translation-weight} gives
$\lambda^{\det}(a\cdot A)=\lambda^{\det}(A)+\chi(A)a$.
For $\chi\ne0$, base change along $[\chi]:E\to E$ gives a
weight-one anchor and a finite residual $E[|\chi|]$-quotient,
as constructed in Proposition~\ref{k3:prop:anchor-cover}.
For $\chi=0$, this determinant is translation invariant.

\begin{definition}[Anchor data on a correspondence]
\label{def:source-target-anchor-memory-datum}
Given chosen anchor maps on the object stacks, retain their
pullbacks from every endpoint, intermediate object and middle
object of an extension or flag. If a proposed anchor comparison
identifies or composes these values, specify that morphism and
its compatibility with flag refinement before passing to a
quotient. The original families remain part of the correspondence.
\end{definition}

For semistable degree-zero vector bundles on $E$, the
Jordan--H\"older divisor
$\Lambda(V)=\sum_i[L_i]\in\operatorname{Sym}^r\operatorname{Pic}^0(E)$
has Abel sum $\det V$. It distinguishes
$\mathcal O_E^{\oplus2}$ from $L\oplus L^{-1}$ for nontrivial
$L$. But translations preserve each degree-zero line's
isomorphism class, hence preserve $\Lambda(V)$.
This spectral invariant does not repair orbit separation on a
free translation orbit.

\begin{definition}[Requirements on a chosen anchor comparison]
\label{def:finite-anchor-residual}
An anchor comparison specifies its actual target, its definition
on families, its translation weight, and any relation to the
determinant. To replace a family parameter by that target one
must additionally identify the resulting kernel equivalence
relation. Compatibility with extensions and transitions is imposed
as a commutative diagram. Divisor union applies to extensions of
semistable degree-zero elliptic bundles, where their
Jordan--H\"older factors concatenate; it is not imposed on arbitrary
perfect pushforwards.
Write $\mathfrak o^\lambda_R=0$ for existence of these maps,
their stated translation weight, their prescribed kernel
equivalence relation, extension compatibility, and transition
compatibility. These are conditions on the chosen maps, with
the zero-weight obstruction below still in force.
\end{definition}

\begin{proposition}[The zero-weight obstruction]
\label{prop:first-window-anchor-residual-vanishing-criterion}
On a nontrivial translation orbit, a translation-invariant
spectral datum cannot determine an $E$-valued anchor of weight
one. In particular this applies to the degree-zero
Jordan--H\"older divisor. On the nonzero-$\chi$ locus the
determinant cover instead gives the weight-one construction,
with its finite stabilizer descent still to be specified.
\end{proposition}

\begin{proof}
A function of an invariant datum has the same value on $A$ and
$a\cdot A$. Weight one requires those values to differ by $a$.
For $a\ne0$ these statements contradict one another. The
nonzero-weight assertion is the explicit isogeny construction
already cited. This rules out the stated invariant replacement,
not every possible anchor construction.
\end{proof}

\subsubsection*{Geometric degree and automorphic degree}

The normal-ordered charge map
$\overline\Pi_X:\widehat\Gamma\to\Lambda_{II}^{2,1}$ is additive on
its specified charge carrier. Write
$\overline\Pi_X(\gamma)=(n,\ell,m)$ and
$m^{\mathrm{Bch}}=\operatorname{pr}_3\overline\Pi_X$.
The integer $m^{\mathrm{Bch}}$ is a root or trace grading. The
geometric degree $b_R^{\mathrm{geom}}$ measures support over $E$.

\begin{definition}[A comparison of degree maps]
\label{def:geometric-borcherds-degree-comparison-datum}
A comparison specifies a domain of geometric charges and the
maps from that domain to the two grading groups. Any equality
between the two degrees is stated on that domain and checked
against addition on its permitted extensions. An affine shift
in a scalar generating series is recorded separately from an
additive charge map.
\end{definition}

\begin{proposition}[An affine shift is not additive]
\label{prop:geometric-borcherds-degree-separation}
If $b$ and $m$ are additive on two charges and their sum, then
$b=m+1$ cannot hold on all three. It may still describe the
absolute indexing of selected coefficients in a generating
series.
\end{proposition}

\begin{proof}
If the equality holds on $\alpha$ and $\beta$, then
$b(\alpha+\beta)=m(\alpha)+m(\beta)+2
=m(\alpha+\beta)+2$, rather than $m(\alpha+\beta)+1$.
\end{proof}

For example, a scalar series written in $s^{d-1}$ has an exponent
$m=d-1$. The shift accounts for its prefactor. It cannot be used
unchanged as an additive degree law for Hall extensions.

\begin{proposition}[The three simple wall labels]
\label{prop:middle-type-ii-wall-wrapped-mixed}
\label{prop:type-ii-walls-hybrid-operation-recovery}
The simple type-II wall images are
\[
 (1,1,0),\qquad(0,1,1),\qquad(0,-1,0).
\]
Suppose three actual sheaf classes have these images under
$\overline\Pi_X$, and suppose their geometric elliptic degrees
are respectively $1,2,1$. Then all three classes have wrapped
support. Local insertions against them use mixed correspondences,
and insertions on both sides use the ordered two-sided flag.
\end{proposition}

\begin{proof}
The wall images are Proposition~\ref{prop:type-ii-wall-images}.
The assumed positive geometric degrees give wrapped support by
Proposition~\ref{prop:positive-elliptic-degree-projects}, or by
the specified larger-domain degree construction. The definitions
of the mixed correspondences then give the stated profiles.
The degrees $1,2,1$ agree with $d=m+1$ on these three selected
labels; the previous proposition prevents extending that affine
rule to all their sums as an additive identification.
\end{proof}

\begin{proposition}[Wall labels after quotient]
\label{prop:type-ii-walls-survive-e-quotient}
\label{prop:no-type-i-wall-enters-hybrid-carrier}
Translation preserves the numerical charges, so the common-torsor
quotient preserves the three selected type-II labels and their
mixed profiles. It introduces no further divisor wall. The
type-I reflections act by the chamber permutations described in
Lemma~\ref{lem:bkm-chamber-s3}. Their reflecting vectors are
$f_{-2}-f_2$, $f_2-f_3$ and $f_{-2}-f_3$; they do not become type-II
divisor roots under this quotient.
\end{proposition}

\begin{proof}
Translation is connected to the identity and preserves numerical
classes. Quotienting the same object and flag stacks consequently
retains their labels. The two root types have the different
ambient divisibilities and divisor orders stated in
Remark~\ref{rem:bkm-divisibility-types}. An elliptic quotient
changes neither the reference lattice nor these divisor orders.
This conclusion concerns the selected wall classes; it neither
constructs their sheaf representatives nor exhausts every
possible geometric charge.
\end{proof}


\subsubsection*{Protected integration is a further map}

The coefficient object is not determined by its Euler
characteristic. A scalar invariant requires a specified functional
on that object, and compatibility with Hall multiplication requires
an equation for the functional on extension correspondences.

\begin{definition}[Protected integration at a bounded stage]
\label{def:finite-protected-integration-datum}
Let $\mathcal H_R=\bigoplus_{\gamma\in\Gamma_R}\mathcal H_{\gamma,R}$
be a chosen reduced coefficient realization, with the operations
already constructed on its charge domain. Choose a target complex
$T_R=\bigoplus_\gamma(T_R)_\gamma$ and degree-zero chain maps
\[
 I_{\gamma,R}^{\mathrm{prot}}:
 \mathcal H_{\gamma,R}\longrightarrow(T_R)_\gamma.
\]
Their sum is $I_R^{\mathrm{prot}}$. A scalar-labelled target has
$(T_R)_\gamma=k e_{\omega_R(\gamma)}$ in degree zero; the label
map $\omega_R$ and its relation to $\overline\Pi_X$ must be
specified. A finite charge window carries only those target
products that are defined on the window, unless a separate
algebra quotient has been proved to exist.
\end{definition}

\begin{proposition}[The integration equations]
\label{prop:protected-integration-definition}
\label{prop:protected-integration-product-compatibility}
\label{prop:protected-integration-coproduct-compatibility}
The finite sum above is a chain map. It preserves the specified
multiplication precisely when, on every permitted pair,
\[
 I_R^{\mathrm{prot}}m_R
 =m_{T,R}(I_R^{\mathrm{prot}}\otimes I_R^{\mathrm{prot}}).
\]
If source and target also have specified coproducts, preservation
of them is the additional equation
\[
 \Delta_{T,R}I_R^{\mathrm{prot}}
 =(I_R^{\mathrm{prot}}\otimes I_R^{\mathrm{prot}})\Delta_R.
\]
The relevant shifts, tensor products and units are included in
these equations. Neither equation follows from the existence of
the individual functionals.
\end{proposition}

\begin{proof}
A finite direct sum of chain maps is a chain map. The displayed
identities are the algebra and coalgebra morphism conditions on
each homogeneous component. Summing them gives the corresponding
condition on every element in the allowed domain. Conversely,
restricting an operation-preserving map to a homogeneous component
gives those identities.
\end{proof}

\begin{proposition}[Primitive compatibility and scalar targets]
\label{prop:protected-integration-primitive-compatibility}
A coaugmentation- and coproduct-preserving map sends primitives
to primitives. If the target is a monoid algebra with its
group-like coproduct
$\Delta e_t=e_t\otimes e_t$ and unit $e_1$, its primitive
space is zero. Such a map therefore kills every source primitive.
\end{proposition}

\begin{proof}
Apply the map to the primitive equation. For the second claim,
write a primitive as $x=\sum_t c_te_t$. For $t\ne1$, the
coefficient of $e_t\otimes e_t$ in
$\Delta x=e_1\otimes x+x\otimes e_1$ gives $c_t=0$.
The coefficient of $e_1\otimes e_1$ then gives $c_1=2c_1$,
hence $c_1=0$ as well.
\end{proof}

This is a restriction on that target coproduct. A numerical
integration map may still be useful as an algebra map or trace,
without preserving a primitive Lie algebra. To compare operator
primitives one needs a target with the appropriate primitive
structure and an actual intertwining map.

Another scalar coalgebra can retain a primitive. Let $e_n$ have
charge $n\geq0$, with
\[
 \Delta e_n=\sum_{a+b=n}e_a\otimes e_b,\qquad
 \varepsilon(e_n)=\delta_{n0},\qquad
 e_a e_b=\binom{a+b}{a}e_{a+b}.
\]
The map $e_n\mapsto x^n/n!$ identifies this Hopf algebra with
the additive polynomial Hopf algebra, so its primitive space is
$k e_1$. The split components of its coproduct have coefficient
one, but its product has binomial coefficients. Using instead
the monoid product $e_ae_b=e_{a+b}$ breaks bialgebra compatibility:
the coefficient of $e_1\otimes e_1$ is one in $\Delta e_2$
and two in $(\Delta e_1)^2$. Thus a scalar coproduct and a
scalar product must be compared together.

If source and target primitive projections $\pi_P,\pi_T$ are
chosen, their compatibility with integration is the additional
equation $I_R^{\mathrm{prot}}\pi_P=\pi_T I_R^{\mathrm{prot}}$.
Preservation of the primitive subspace alone does not give this
equation for arbitrary projections.

\begin{proposition}[Homogeneous integration]
\label{prop:protected-integration-gram-degree}
\label{prop:protected-integration-pi-x-separation}
Suppose the target has monomials $q^nr^\ell s^m$ graded by
$\overline\Pi_X(\gamma)=(n,\ell,m)$. Then an integration map
preserves this grading exactly when it maps each source component
to the corresponding monomial component. The geometric degree
$b_R^{\mathrm{geom}}$ does not determine $m$ without the
comparison of Definition~\ref{def:geometric-borcherds-degree-comparison-datum}.
\end{proposition}

\begin{proof}
The grading condition is the statement that each homogeneous
component has the displayed image. Additivity of
$\overline\Pi_X$ makes target multiplication homogeneous.
The geometric degree has a different definition and can be
substituted only after the asserted relation between the maps is
proved on this domain.
\end{proof}

When the numerical charge index $\Gamma_R$ is not already the
normal-ordered extension, choose its lift
$\widehat\omega_R:\Gamma_R\to\widehat\Gamma$ explicitly.
For a target degree map $g_R$, the homogeneous comparison is
\[
 g_R\omega_R=\overline\Pi_X\widehat\omega_R.
\]
All three maps respect the permitted additions. The preceding
formula uses $\widehat\omega_R=1$ when the source was already
indexed by extended charges.

For a transition $\rho_{R'R}:\mathcal H_R\to\mathcal H_{R'}$
and $t_{R'R}:T_R\to T_{R'}$, the integration compatibility is
$t_{R'R}I_R^{\mathrm{prot}}=I_{R'}^{\mathrm{prot}}\rho_{R'R}$.
It must hold before forming a limiting map. A finite-type bound
alone gives neither this equation nor exactness of the inverse
limit.

\subsubsection*{The compact factorization problem}

The K3 stability conditions and their projective moduli
\cite{BridgelandK3Stability,BayerMacriProjectivity} provide the
surface input. A choice of moduli on $S\times E$ must additionally
satisfy the product-stability and support conditions of
Appendix~\ref{app:stability-support}.

\begin{definition}[A local and wrapped factorization datum]
\label{def:hybrid-wrapped-factorization-datum}
A compact hybrid datum on the chosen charge domain consists of
local incidence stacks and wrapped family stacks, their extension
and collision diagrams, and an admissible coefficient theory on
those diagrams. Its coefficient objects are equipped with the
required critical and reduced orientation structures. Its maps
include the binary comparisons, all finite flag coherences,
local configuration descent, and the separate collision
comparisons. Elliptic quotient is taken over a common torsor,
with coherent coefficient descent and relative unit $BE$.
Optional anchors and protected integration maps carry the
specific comparison diagrams just defined.

At bounded stages all charge sums and intermediate objects must
belong to the permitted domain. A family of such data has
specified transition maps preserving its objects, operations,
units and coherences. The completion, when needed, is formed
only with a stated topology and the corresponding inverse-limit
or derived inverse-limit construction.
We abbreviate this collection of existence and compatibility
requirements by $\mathfrak O_{\mathrm{hyb},R}=0$. The notation
includes the anchor and integration conditions only when those
maps are requested. It is not an assertion that a coefficient
object exists because its prospective defects have been named.
\end{definition}

\begin{proposition}[Scope of the constructed data]
\label{prop:hybrid-datum-scope}
The formal cyclic comparison and its finite-flag extension
identify the specified formal sheaf deformations with the
corresponding completions of the quiver models. The common-torsor
construction gives the relative
elliptic quotient of the sheaf extension diagrams. The critical
Hall theorem supplies coefficient operations on its unreduced
moduli domain. A compact reduced hybrid factorization object
follows when these constructions are joined by the coefficient,
orientation and collision comparisons in the preceding definition.
\end{proposition}

\begin{proof}
The first assertion is
Proposition~\ref{hall:prop:cyclic-collection} and
Theorem~\ref{hall:thm:formal-flags}; the second is
Proposition~\ref{hall:prop:common-torsor}; the third is
Theorem~\ref{thm:hall-product}. With the additional comparisons,
all prescribed operations and coherences are maps on one
coefficient diagram. Proposition~\ref{prop:higher-coloured-residual-vanishing-criterion}
then gives the claimed factorization operations on that domain.
\end{proof}

\begin{corollary}[The compact comparison remains a coefficient problem]
\label{cor:no-quotient-first-hybrid-repair}
A degree label, a determinant anchor, a product of independent
quotients or the scalar Igusa product does not construct the
compact reduced coefficient comparisons. Working relatively over
$BE$ resolves the common-torsor part of the quotient problem;
it leaves the reduced critical maps and their collision
compatibilities to be constructed.
\end{corollary}

The distinction is visible already at a collision. A formal
comparison of deformation algebras acts on the completed
neighbourhood, whereas a constructible coefficient comparison
must act on the actual extension and collision spaces with the
required supports. A homotopy defined only after inverting a
separation parameter does not solve an equation on its zero
locus. These supported maps, together with their reduced
orientations, are the next inputs needed for the compact Hall
structure.

A holomorphic factorization theory on $X$ would provide another
possible source. Its pushforward along $p_E$ and its comparison
with the Hall coefficients require specified functors preserving
the relevant supports, tensor products, units and collision
operations. The right-$D$-module chiral construction of
\cite{BD} and the smooth factorization framework of
\cite{CostelloGwilliam2} use different coefficient categories;
an equivalence between their realizations is an additional map.
Neither the local deformation model nor a common scalar trace
supplies that map.


\subsection{Reduced complexes and coherent orientations}
\label{subsec:cosection-twisted-orientation}

An orientation is a square root together with its squaring map.  Its
compatibility with an extension is another isomorphism, and compatibility
with a flag is an equation between such isomorphisms.  These distinctions
matter on \(X=S\times E\): removing the elliptic translation direction,
changing a formal coordinate on \(E\), and descending through a stabilizer
all act on the deformation complex.  We first construct the complexes and
local comparisons.  We then formulate the global orientation they must
preserve.

\subsubsection*{The deformation complex and the reduced recipient}

Fix a smooth projective K3 surface \(S\), an elliptic curve \(E\), and
nonzero forms \(\sigma_S\in H^0(S,\Omega_S^2)\) and
\(\omega_E\in H^0(E,\Omega_E^1)\).  Their product trivializes \(K_X\).
K3 stability conditions and projective moduli of stable objects
\cite{BridgelandK3Stability,BayerMacriProjectivity} provide the surface
input.  On \(X\), we additionally fix the product-stability and support
conditions of Appendix~\ref{app:stability-support}.
For a perfect family \(\mathcal F\) on \(X\times B\), put
\[
 K_{\mathcal F}=Rp_*R\mathcal Hom(\mathcal F,\mathcal F),
 \qquad p:X\times B\longrightarrow B.
\]
We work on bases for which this complex is perfect and its formation
commutes with the base changes in use.  Proper smooth projection and a
perfect family give this property.  At a point represented by \(F\), the
cohomology of \(K_F\) is \(\operatorname{Ext}^*(F,F)\).

\begin{proposition}[Restriction of the shifted symplectic structure]
\label{prop:ptvv-finite-substack-symplectic}
The derived stack of perfect complexes on \(X\) carries the
\((-1)\)-shifted symplectic structure associated with the chosen volume
form.  A derived open substack inherits it.  On a different derived
substack, the pulled-back closed two-form is symplectic only if its
induced map from tangent to shifted cotangent complexes is an equivalence.
Finiteness of a stability window does not imply this condition.
At \(F\), the tangent complex of the unrigidified stack is \(K_F[1]\);
in particular,
\[
 H^{-1}(T_F)=\operatorname{Hom}(F,F).
\]
For nonzero simple \(F\), its cotangent complex therefore has a nonzero
\(H^1\), and is not of amplitude \([-1,0]\).
\end{proposition}
\begin{proof}
The shifted form on the derived moduli stack is the trace construction of
\cite{PTVV2013}.  An open immersion identifies tangent and cotangent
complexes, so it preserves nondegeneracy.  For a general morphism, pulling
back a closed form preserves closedness but gives no such identification.
The tangent formula follows from the endomorphism deformation complex.
Duality gives \(H^1(L_F)=\operatorname{Hom}(F,F)^\vee\).  The identity of a
nonzero simple object spans a copy of \(\mathbb C\) in this group.
\end{proof}

\begin{proposition}[Critical truncation and orientation coefficients]
\label{prop:joyce-dcritical-truncation}
Let \(\mathfrak U\) be a derived scheme locally of finite presentation
carrying a
\((-1)\)-shifted symplectic form, and let \(U=t_0\mathfrak U\).
Then \(U\) has its associated d-critical structure, with canonical line
\[
 K_{U,s}\simeq\det(L_{\mathfrak U})|_{U^{\mathrm{red}}}.
\]
A square-root pair for this line defines the oriented vanishing-cycle
sheaf on \(U\).  On a critical chart \(i:R\hookrightarrow V\), with
function \(f:V\to\mathbb A^1\), write \(\phi_f\) for the
vanishing-cycle functor shifted by \([-1]\).  Its expression is
\[
 \mathcal P_{U,s}|_R
 \simeq i^*\phi_f\mathbb Q_V[\dim V]\otimes_{\mu_2}Q_{R,V,f,i},
\]
with the convention for shifted vanishing cycles fixed as in
\cite[Theorem~6.9]{BBDJS2015}.  Here \(Q_{R,V,f,i}\) is the principal
\(\mu_2\)-bundle of square-root identifications with \(i^*K_V\).
\end{proposition}
\begin{proof}
The d-critical structure and canonical-line comparison are
\cite[Theorem~6.6]{BBDJS2015}.  The oriented coefficient sheaf and the
comparison on embeddings of charts are Theorem~6.9 of that paper.
The tensor in the displayed formula uses the sign representation of
\(\mu_2\).  It is a tensor of constructible coefficients; it is not a
tensor with the coherent orientation line.  On an Artin stack one must
also supply compatible smooth-atlas descent.  A reduced or rigidified
moduli problem must first satisfy the hypotheses on its own derived
recipient.
\end{proof}

Choose the translation vector \(v_E\) dual to \(\omega_E\).  Contraction
of the Atiyah class gives its infinitesimal action
\[
 u_F=\iota_{v_E}\operatorname{At}(F)\in\operatorname{Ext}^1(F,F).
\]
Let \(\operatorname{Tr}_X:\operatorname{Ext}^3(F,F)\to\mathbb C\)
be the Serre trace with the chosen volume form.

\begin{proposition}[Semiregularity and the translation direction]
\label{prop:k3-semiregularity-cosection}
\label{prop:k3-cosection-surjectivity-criterion}
The functional
\[
 \sigma_F:\operatorname{Ext}^2(F,F)\longrightarrow\mathbb C,
 \qquad \xi\longmapsto\operatorname{Tr}_X(\xi\circ u_F)
\]
is dual to the infinitesimal translation class.  It is nonzero, hence
surjective, precisely when \(u_F\ne0\).  With compatible relative Atiyah
and trace maps it is induced by a morphism
\[
 K_{\mathcal F}\longrightarrow\mathcal O_B[-2].
\]
\end{proposition}
\begin{proof}
Serre duality gives a perfect pairing between
\(\operatorname{Ext}^2(F,F)\) and \(\operatorname{Ext}^1(F,F)\).
Evaluation against \(u_F\) proves the pointwise assertion.  The relative
Atiyah class, contraction, composition and trace give the displayed
morphism in families.  Its target has cohomology in degree two, as the
functional requires.  Equivalently, the Atiyah--trace expression can be
paired with the pulled-back K3 two-form; the volume convention fixes the
sign of that expression.  A statement of surjectivity on a family needs
nonvanishing at every relevant point, not only at its generic point.
\end{proof}

\begin{definition}[A cosection fibre]
\label{def:rhom-red-self-ext-cone}
Let \(K_0\) be a specified perfect deformation complex with a morphism
\(\sigma:K_0\to\mathcal O_B[-2]\).  Its cosection fibre is
\[
 K^{\mathrm{fib}}=\operatorname{fib}(\sigma),\qquad
 K^{\mathrm{fib}}\longrightarrow K_0
 \xrightarrow{\sigma}\mathcal O_B[-2].
\]
One possible \(K_0\) is the fibre of the trace map
\(K_{\mathcal F}\to Rp_*\mathcal O_{X\times B}\); using it requires
restricting the cosection to that fibre.  We use the notation
\(K^{\mathrm{red}}\) only after specifying which scalar, orbit and dual
obstruction directions the reduced deformation problem removes.
\end{definition}

\begin{proposition}[Perfectness and scalar descent]
\label{prop:rhomred-perfectness}
The fibre of a morphism of perfect complexes is perfect.  The
endomorphism complex of a family on a scalar \(\mathbb G_m\)-gerbe has
weight zero and descends to the rigidified base, whenever the stated
perfect-complex descent applies.  The trace, relative Atiyah class along \(X\), contraction and
cosection maps descend with it.  These statements do not assert that the
universal family, of scalar weight one, descends.
\end{proposition}
\begin{proof}
Perfect complexes form a stable category.  Scalar conjugation acts
trivially on \(R\mathcal Hom(\mathcal F,\mathcal F)\), since the two
scalar weights cancel.  Locally trivialize the gerbe and apply descent
to this weight-zero complex and its natural maps.  The same cancellation
proves compatibility on overlaps.  It does not cancel the weight of
\(\mathcal F\) itself.
\end{proof}

The trace-free fibre need not remove scalar automorphisms.  Indeed,
\(\operatorname{tr}(\mathrm{id}_F)=\operatorname{rk}(F)\); for a
rank-zero sheaf, the identity survives this trace.  Likewise, deleting
one obstruction quotient is not by itself a symplectic reduction.
At a point where \(u_F\ne0\), Serre duality induces a perfect
pairing between \(\operatorname{Ext}^1(F,F)/\mathbb C u_F\) and
\(\ker\sigma_F\subset\operatorname{Ext}^2(F,F)\).  This is the
linear reduction of the orbit direction and its dual obstruction.  Constructing the corresponding family,
its derived critical charts and their gluing remains necessary.  A
perfect obstruction theory, a cosection-localized cycle
\cite{KiemLi2013}, and such an oriented critical family are different
constructions.

There are positive reduced geometric constructions in more specific
moduli problems.  For stable pairs on \(S\times E\) with nonzero
K3 curve class, the quotient and reduced virtual classes of
\cite[Theorems~1--2]{OberdieckReducedPT} satisfy the comparison stated
in Theorem~\ref{thm:reduced-virtual-class-K3xE}.  That result retains
its precise stable-pair domain.  Extending it to the reduced critical
coefficient and complete-flag structures above is a further problem.

\subsubsection*{Square-root pairs and their obstruction}

For a perfect complex \(K\), use the determinant convention
\[
 \det K=\bigotimes_j(\det K^j)^{(-1)^j}.
\]
Thus \(\det(K[1])=(\det K)^{-1}\) and
\(\det(K^\vee)=(\det K)^{-1}\).  Determinants of exact triangles carry
the usual functorial identifications
\cite{KnudsenMumfordDet,DeligneDetCohomology}.

\begin{definition}[The reduced determinant line]
\label{def:rhomred-determinant-line}
A specified perfect reduced complex determines
\(\mathscr D^{\mathrm{red}}=\det K^{\mathrm{red}}\).
An orientation is a pair
\[
 (o,\rho),\qquad \rho:o^{\otimes2}\xrightarrow{\sim}
 \mathscr D^{\mathrm{red}}.
\]
When this orientation is used for critical coefficients, a comparison
between \(\mathscr D^{\mathrm{red}}|_{B^{\mathrm{red}}}\) and the
canonical d-critical line is part of the datum.  A shift or dualization
of the deformation complex changes that comparison.
\end{definition}

\begin{proposition}[Kummer obstruction and choices]
\label{prop:rhomred-square-root-criterion}
\label{prop:alpha-red-vanishing-criterion}
\label{prop:alpha-red-nullhomotopy-criterion}
Over a complex scheme or stack on its relevant descent site, let
\(\delta:\operatorname{Pic}(B)\to H^2(B,\mu_2)\) be the Kummer
boundary.  A line \(\mathscr D\) has a square-root pair if and only if
\(\alpha=\delta(\mathscr D)\) vanishes.  If it vanishes, the
isomorphism classes of square-root pairs form a torsor under
\(H^1(B,\mu_2)\); the automorphisms of a pair form \(H^0(B,\mu_2)\).
The underlying root lines, after forgetting their squaring maps, form
a torsor under \(\operatorname{Pic}(B)[2]\).
\end{proposition}
\begin{proof}
Apply cohomology to
\[
 1\longrightarrow\mu_2\longrightarrow\mathbb G_m
 \xrightarrow{z\mapsto z^2}\mathbb G_m\longrightarrow1;
\]
see \cite[Tags 03PK, 03PL]{StacksProject}.  Local roots of transition
functions have a \(\mu_2\)-valued defect on triple overlaps.  This is
\(\delta(\mathscr D)\), and a trivializing cochain glues a root pair.
The ratio of two pairs is a line with a trivialization of its square,
which is a \(\mu_2\)-torsor.  A root-preserving automorphism is a
section of \(\mu_2\).  Forgetting the trivialization gives the last
assertion.  In general \(H^1(B,\mu_2)\) is larger than
\(\operatorname{Pic}(B)[2]\), because units modulo squares also occur.
\end{proof}

\begin{proposition}[Integral and rational characteristic classes]
\label{prop:rhomred-w2-formula}
The image of the Kummer class of a complex line in analytic cohomology
is its integral first Chern class reduced modulo two.  It is also the
second Stiefel--Whitney class of the underlying oriented real plane.
A rational Grothendieck--Riemann--Roch formula alone does not determine
this class.
\end{proposition}
\begin{proof}
The first assertion follows by comparing the square map on transition
functions with the integral exponential sequence.  For a complex line,
\(w_2=c_1\bmod2\), as in \cite{MilnorStasheffCharacteristicClasses}.
The map from integral to rational cohomology kills torsion, so its value
cannot recover arbitrary mod-two information.  For example, the sign
line on \(B\mu_2\) has nonzero Kummer obstruction although its rational
first Chern class vanishes.  Its character has no square root character
of \(\mu_2\).
\end{proof}

For a proper smooth morphism \(p:Y\to B\) of smooth schemes and a
perfect complex \(\mathcal K\), the rational calculation retains its
usual form:
\[
 c_1(\det Rp_*\mathcal K)_{\mathbb Q}
 =\left[p_*\bigl(\operatorname{ch}(\mathcal K)
             \operatorname{td}(T_p)\bigr)\right]_1.
\]
It is the degree-one part of Grothendieck--Riemann--Roch
\cite[Tag 02UO]{StacksProject}.  The Deligne--Mumford stack version
uses the appropriate representation-valued theory
\cite{ToenRRRoot}.  Neither rational calculation removes the need
for the integral Kummer class.

A cosection is a morphism of complexes on \(B\), not a map of bases
along which a line can be pulled back.  If a reduction constructs an
additional perfect complex \(C_{\mathrm{cs}}\), one may instead define
\[
 \mathscr D_{\mathrm{tw}}
 =\mathscr D^{\mathrm{red}}\otimes\det C_{\mathrm{cs}},\qquad
 \alpha_{\mathrm{red}}=\delta(\mathscr D_{\mathrm{tw}}).
\]
This notation requires the construction of \(C_{\mathrm{cs}}\).
The cokernel of a cosection has a determinant only when supplied with
the requisite perfect model.  The root criterion applies to the actual
line \(\mathscr D_{\mathrm{tw}}\), including its torsion information.

\subsubsection*{Formal coordinates and the full determinant unit}

Consider sheaves near one fibre of \(S\times E\), annihilated by a
locally uniform power of a chosen normal coordinate \(t\).  Such a
sheaf is equivalently a coherent sheaf \(F\) on \(S\) with a nilpotent
endomorphism \(N_F\): multiplication by \(t\) is \(N_F\).
Restriction and extension of this module structure preserve exact
sequences and finite flags.  With strict endomorphisms on chosen
resolutions, put
\[
 C(F,G)=R\operatorname{Hom}_S(F,G),\qquad
 \mathcal K(F,G)=C(F,G)\otimes\Lambda(\theta),\quad |\theta|=1.
\]
For a homogeneous coefficient \(a\), the differential and composition are
\begin{align*}
 D(a\theta^\epsilon)
 &=d_Ca\,\theta^\epsilon
   +(-1)^{|a|}(N_Ga-aN_F)\theta^{\epsilon+1},\\
 (a\theta^\epsilon)(b\theta^\delta)
 &=(-1)^{\epsilon|b|}ab\,\theta^{\epsilon+\delta},
 \qquad \theta^2=0.
\end{align*}
The relative resolution
\[
 0\longrightarrow\mathcal O_S[t]\otimes F
 \xrightarrow{t-N_F}\mathcal O_S[t]\otimes F
 \longrightarrow F\longrightarrow0
\]
combined with a resolution on \(S\) computes this Hom complex.  The commutator with \(N\) commutes with \(d_C\) and is a
derivation, so these formulas define a dg category.  This construction
uses the actual endomorphisms; it does not require formality of
\(C(F,G)\).

Let \(f(t)=\sum_{r\ge1}c_rt^r\), with \(c_1\) invertible.  For
composable maps \(a_i:F_{i-1}\to F_i\), define
\[
 D_n^f(a_n,\ldots,a_1)=
 \sum_{r\ge n}c_r
 \sum_{j_0+\cdots+j_n=r-n}
 N_{F_n}^{j_n}a_nN_{F_{n-1}}^{j_{n-1}}
 \cdots a_1N_{F_0}^{j_0}.
\]
At fixed arity every sum is finite under the chosen nilpotence bounds.
The coefficient maps need not commute with \(N\).  The first divided
difference satisfies
\[
 D_1^f(N_Ga-aN_F)=f(N_G)a-af(N_F).
\]

\begin{theorem}[Coordinate functors at every arity]
\label{thm:orientation-coordinate-functor}
Use suspension \(s\) of degree \(-1\), with bar operations
\[
 b_1(sa)=-s(Da),\qquad
 b_2(sa\otimes sb)=(-1)^{|a|}s(ab).
\]
There is a strictly unital \(A_\infty\) functor \(\mathcal F_f\)
from the category with endomorphisms \(N_F\) to that with
endomorphisms \(f(N_F)\).  Its first component is the identity on
unmarked coefficients and \(D_1^f\) on the \(\theta\) coefficient.
For \(n\ge2\), its suspended components are
\[
 (\mathcal F_f)_n
 \bigl(s(a_n\theta),\ldots,s(a_1\theta)\bigr)
 =s\bigl(D_n^f(a_n,\ldots,a_1)\theta\bigr),
\]
and vanish if any input is unmarked.  For another invertible coordinate
\(g\), these coalgebra maps satisfy
\[
 \mathcal F_g\mathcal F_f=\mathcal F_{g\circ f},\qquad
 \mathcal F_{f^{-1}}\mathcal F_f=\mathrm{id}.
\]
\end{theorem}
\begin{proof}
On marked inputs the suspended degree is the coefficient degree.
Consequently the internal differential obeys exactly the graded
Leibniz rule for \(D_n^f\); the closed endomorphisms \(N_F\) have
degree zero.  Products of two marked inputs vanish.  In the remaining
higher identity, exactly one input is unmarked.  Replacing that input
by its commutator with \(N\) makes the inner sum telescope over the
adjacent exponent.  For the monomial \(f(t)=t^r\), its endpoint terms
are the two adjacent compositions, with the tensor signs dictated by
\(b_1,b_2\).  Linearity and finite nilpotence give the identity for
\(f\).  With two or more unmarked inputs only the ordinary binary
composition identity can survive.  The unit is unmarked, so all higher
components containing it vanish.

For composition, place the \(N_{F_i}\) on the diagonal of an upper
triangular block matrix and the \(a_i\) on its first off-diagonal.
The relevant \(n\)-step block of \(f\) of this matrix is \(D_n^f\).
Applying \(g\) and summing over intervening blocks is precisely the
partition formula for composition of the suspended functors.  The
identity \(g(f(M))=(g\circ f)(M)\) proves the claimed equality.
The same calculation with the inverse formal series proves the inverse
identity at each nilpotence bound.
\end{proof}

Already in arity two, the higher map is needed:
\[
 D_1^f(ab)-aD_1^f(b)=D_2^f(N_Ha-aN_G,b).
\]
It is this identity, with its graded counterparts, that corrects the
failure of the first component to preserve composition.

The K3 Serre trace \(\tau_S:C(F,F)\to\mathbb C[-2]\) gives a CY3
trace on the \(\theta\) coefficient.  For complementary powers of
\(\theta\), the pairing is
\[
 \langle a\theta^\epsilon,b\theta^\delta\rangle
 =(-1)^{\epsilon|b|}\tau_S(ab),\qquad \epsilon+\delta=1.
\]
If the normal volume is \(h(t)\,dt\), its coefficient after changing
to \(f(t)\) is \(h(t)/f'(t)\).  Cyclicity gives the exact trace identity
\[
 \tau_S\!\left(h(N)f'(N)^{-1}D_1^f(b)\right)
 =\tau_S(h(N)b).
\]
Indeed, for any coefficient \(k(N)\), cyclicity gives
\(\tau_S(k(N)D_1^f(b))=\tau_S(k(N)f'(N)b)\).
Take \(k(N)=h(N)f'(N)^{-1}\).  This proves transport of the
displayed trace functional.  A comparison of negative cyclic classes
or a global critical coefficient functor requires its own construction.

\begin{proposition}[The determinant and its nonreduced square root]
\label{prop:orientation-coordinate-determinant}
Let the coefficient base be a local Artin \(\mathbb C\)-algebra
\(R\) with residue field \(\mathbb C\),
and let \(C(F,G)\) have a bounded finite free model.  The determinant
of the coordinate map on \(\mathcal K(F,G)\) is
\[
 d_f(F,G)=\prod_j
 \det\!\left(D_1^f|_{C^j(F,G)}\right)^{(-1)^{j+1}}.
\]
On a self-Hom complex put \(\chi=\chi(C(F,F))\).  Suppose its
reduction has determinant \(\bar c_1^{-\chi}\), and \(\chi\) is even.
Then \(\rho_f=c_1^\chi d_f(F,F)\) lies in \(1+\mathfrak m_R\),
and the unique root with reduction \(\bar c_1^{-\chi/2}\) is
\[
 j_f(F)=c_1^{-\chi/2}
 \exp\!\left(\tfrac12\log\rho_f\right),\qquad j_f(F)^2=d_f(F,F).
\]
\end{proposition}
\begin{proof}
The unmarked part of the complex is unchanged, whereas the marked
copy of \(C^j\) is in degree \(j+1\).  This gives the determinant
exponent.  Both logarithm and exponential are finite series on the
nilpotent maximal ideal.  Squaring is an automorphism of
\(1+\mathfrak m_R\), proving existence and uniqueness.  On the
reduced K3 self-Hom complex, evenness of the Mukai lattice gives
\(\chi(F,F)\in2\mathbb Z\).  The determinant of a polynomial in the
commuting nilpotent left and right actions has constant diagonal
\(c_1\), which gives the stated reduced determinant.
\end{proof}

For example, take \(R=\mathbb C[\varepsilon]/(\varepsilon^2)\),
\(F=\mathcal O_S\), \(N=\varepsilon\,\mathrm{id}\), and
\(f(t)=t+at^2\).  The two self-Ext groups in degrees zero and two
each see \(D_1^f=1+2a\varepsilon\).  Hence
\[
 d_f(F,F)=1-4a\varepsilon,
 \qquad j_f(F)=1-2a\varepsilon.
\]
Their reduced values are both one, but the actual coordinate comparison
is not.  Nilpotent units cannot be discarded from an orientation map.

\subsubsection*{Equivariant square roots and the Borel filtration}

Let \(M\) be connected and carry an \(E\)-action and an equivariant determinant line
\(\mathscr D\).  An orientation on \([M/E]\) is an equivariant
square-root pair on \(M\).  Equivariance includes its cocycle on
\(E\times E\times M\).  We write
\(M_E=EE\times_E M\) for the analytic homotopy quotient when discussing
mod-two characteristic classes.

\begin{proposition}[The base part of an equivariant obstruction]
\label{prop:free-e-borel-class-criterion}
\label{prop:free-e-borel-vanishing-criterion}
The degree-two base term in the Borel spectral sequence is
\[
 E_\infty^{2,0}
 =H^2(BE,\mathbb F_2)/\operatorname{im}
 \bigl(d_2:H^1(M,\mathbb F_2)\to H^2(BE,\mathbb F_2)\bigr).
\]
It identifies with the image of \(H^2(BE,\mathbb F_2)\) in
\(H^2(M_E,\mathbb F_2)\).  There is generally no canonical projection
from the latter group to \(H^2(BE,\mathbb F_2)\).
If a class \(\alpha_E\) restricts to zero on \(M\), then it is the
pullback of a base class \(a\), determined modulo \(\operatorname{im}d_2\).
It vanishes precisely when that representative lies in
\(\operatorname{im}d_2\).
\end{proposition}
\begin{proof}
Topologically \(E\) is a real two-torus, so \(BE\) is
\(K(\mathbb Z^2,2)\), with \(H^1(BE,\mathbb F_2)=0\).
In total degree two, the only incoming differential to \((2,0)\) is
the stated transgression, and no outgoing differential survives.
The filtration has no \((1,1)\) term.  Hence a class vanishing on the
fibre lies in its \((2,0)\) piece.  The edge map is a pullback from the
base, not a projection onto it.  Choosing generators \(u_1,u_2\) of
\(H^2(BE,\mathbb F_2)\) writes a representative as
\(a_1u_1+a_2u_2\), but does not make that representative unique.
\end{proof}

For the translation action of \(E\) on itself, the homotopy quotient
is contractible.  Both base generators are transgressions.  This free
action gives a concrete reason that two coefficients in \(H^2(BE)\)
cannot be assigned canonically to an arbitrary equivariant class.
On a free locus admitting a geometric quotient, the equivariant root
problem is the ordinary root problem on that quotient.

\begin{proposition}[Equivariant nulltrivializations]
\label{prop:free-e-borel-nulltrivialization-criterion}
The vanishing of \(\delta(\mathscr D)\) on \([M/E]\) is equivalent
to the existence of an equivariant square-root pair.  Its choices form
a torsor under \(H^1([M/E],\mu_2)\), with automorphism group
\(H^0([M/E],\mu_2)\).  A nonequivariant root and a vanishing ordinary
class on \(M\) do not specify this equivariant nulltrivialization.
\end{proposition}
\begin{proof}
Apply the Kummer argument on the quotient stack, or apply it on the
simplicial action presentation.  The equations on double and triple
overlaps include the action and its composition.  Forgetting those
equations gives only the nonequivariant root.
\end{proof}

\subsubsection*{Finite stabilizers and character lines}

For an object \(x\) with finite stabilizer \(H\subset E\), restriction
along \(BH\to[M/E]\) is always defined.  It requires no choice of a
Borel edge splitting.  Such restrictions are necessary local tests;
they need not detect the whole equivariant obstruction.
Write \(\beta_x\in H^2(BH,\mu_2)\) for the restricted gerbe class.

\begin{proposition}[Finite root gerbes]
\label{prop:finite-stabilizer-beta-class-criterion}
\label{prop:finite-stabilizer-beta-vanishing-criterion}
\label{prop:finite-stabilizer-beta-nulltrivialization-criterion}
The root gerbe at \(x\) is neutral precisely when \(\beta_x=0\).
In that case root pairs form an \(H^1(BH,\mu_2)\)-torsor.
For the root gerbe of an actual line, write its character as
\(\chi\in\operatorname{Hom}(H,\mathbb C^*)\), additively.
Its root characters are exactly the solutions of
\[
 2\lambda=\chi.
\]
\end{proposition}
\begin{proof}
A line on \(BH\) is a one-dimensional representation.  Its tensor
square doubles its character.  The Kummer sequence therefore gives
both the obstruction and the stated torsor.  More generally, a
\(\mu_2\)-gerbe has a normalized group two-cocycle; a trivialization
is a cochain whose coboundary is that cocycle.
\end{proof}

\begin{proposition}[Stack descent and inertia]
\label{prop:finite-stabilizer-linearization-character-criterion}
\label{prop:finite-stabilizer-linearization-vanishing-criterion}
\label{prop:finite-stabilizer-gerbe-linearization-independence}
A root character \(\lambda\) need not vanish for an orientation on
\(BH\) or on the quotient stack.  Vanishing is an additional condition
for descent to the coarse space, or for a specified inertia-trivial
normalization.  Existence of a root and this additional condition are
distinct requirements.
\end{proposition}
\begin{proof}
Take \(H=\mu_2\) and the trivial determinant line.  Both the trivial
line and the sign line, with their squaring maps, are roots on
\(B\mu_2\).  The sign line does not descend to the point, because its
inertia acts nontrivially.  Conversely the determinant sign character
itself has no root character: doubling either character of \(\mu_2\)
gives zero.  Thus a gerbe obstruction precedes the choice of root
linearization, while coarse descent further restricts that choice.
\end{proof}

The finite subgroups of a complex elliptic curve are abelian of rank
at most two.  Their odd-primary part contributes no positive-degree
cohomology with \(\mathbb F_2\) coefficients.  The following formulas
include cyclic subgroups and unequal two-primary factors.

\begin{proposition}[Two-primary classes]
\label{prop:finite-stabilizer-two-primary-class}
\label{prop:finite-stabilizer-two-primary-mixed-term}
\label{prop:finite-stabilizer-odd-transfer}
Let \(H_{(2)}=C_{2^a}\times C_{2^b}\), omitting a factor when it is
trivial.  A nontrivial cyclic factor has generators \(x\) in degree
one and \(u\) in degree two, with
\[
 H^*(BC_{2^a},\mathbb F_2)=
 \begin{cases}
 \mathbb F_2[x],&a=1,\quad u=x^2,\\
 \Lambda(x)\otimes\mathbb F_2[u],&a>1.
 \end{cases}
\]
For two nontrivial factors,
\[
 H^2(BH,\mathbb F_2)
 =\mathbb F_2u_1\oplus\mathbb F_2(x_1x_2)
  \oplus\mathbb F_2u_2.
\]
If a character line has exponents \(k_1,k_2\) on these factors,
its Kummer class is
\[
 \delta(\chi)=(k_1\bmod2)u_1+(k_2\bmod2)u_2.
\]
In particular the mixed term of an arbitrary gerbe is available in
\(H^2\), but it cannot occur in the root gerbe of a character line.
When \(a=b\ge2\), the nonzero class \(x_1x_2\) restricts to zero
on every cyclic subgroup of \(H_{(2)}\).  Thus cyclic restrictions
do not determine an arbitrary degree-two gerbe class.
\end{proposition}
\begin{proof}
The periodic resolution of a cyclic group computes the stated ring:
the square of the degree-one generator is nonzero for order two and
zero for higher powers of two.  The K\"unneth formula gives the three
degree-two summands.  A character of the product is a tensor product
of characters pulled back from the two factors.  The Kummer boundary
is additive and commutes with pullback.  On a cyclic factor, doubling
characters has cokernel \(\mathbb F_2\), detected by the parity of
its exponent; this defines \(u\).  There is therefore no mixed term.
Restriction and transfer for the odd-primary subgroup show that its
positive-degree \(\mathbb F_2\)-cohomology vanishes.
To check the cyclic restrictions when \(a=b\ge2\), let
\(C\subset H_{(2)}\) be cyclic.  If \(C\) has order two, then
\(C\subset2H_{(2)}\), so both mod-two characters \(x_1,x_2\)
restrict to zero.  If \(|C|\ge4\), its degree-one cohomology is
spanned by a class \(\xi\) with \(\xi^2=0\).  Both restricted
characters are multiples of \(\xi\), so their product is zero.
The trivial subgroup also has zero degree-two cohomology.  The
K\"unneth summand above shows that \(x_1x_2\) itself is nonzero.
\end{proof}

\begin{proposition}[The Klein four calculation]
\label{prop:finite-stabilizer-klein-four-class}
\label{prop:finite-stabilizer-klein-four-b20}
\label{prop:finite-stabilizer-klein-four-b11}
\label{prop:finite-stabilizer-klein-four-b02}
\label{prop:finite-stabilizer-klein-four-zero-restrictions}
For \(H=C_2\times C_2\), every degree-two class has a unique expression
\[
 \beta=b_{20}x_1^2+b_{11}x_1x_2+b_{02}x_2^2.
\]
Its restrictions to the first factor, second factor and diagonal have
coefficients \(b_{20}\), \(b_{02}\) and
\(b_{20}+b_{11}+b_{02}\), respectively.  These three restrictions
detect \(\beta\).  For a line-root gerbe one already has
\(b_{11}=0\).
\end{proposition}
\begin{proof}
Pull back the polynomial ring \(\mathbb F_2[x_1,x_2]\) along the
three inclusions.  The diagonal sends both generators to the same
class \(x\).  The resulting three linear equations recover all three
coefficients.  The final assertion is the character-line calculation
above.  For higher cyclic orders, the degree-one restriction to the
order-two subgroup may vanish; the same three-test argument cannot
be transferred without computing those restriction maps.
\end{proof}

\begin{proposition}[Scope of stabilizer detection]
\label{prop:finite-stabilizer-orientation-detection}
An equivariant root restricts to a root at every stabilizer.  The
converse requires a descent statement for the global gerbe, including
its variation on the moduli space and the action overlaps.  It does
not follow from the pointwise classes \(\beta_x\) alone.
\end{proposition}
\begin{proof}
Restriction preserves square-root pairs.  Pointwise stabilizer
restriction forgets the remaining overlap maps.  Already for a free
action all stabilizers are trivial, whereas the quotient may carry a
line without a square root.  Thus stabilizer tests cannot in general
recover its Kummer class.
\end{proof}

\subsubsection*{Cross determinants and complete flags}

For a perfect pair on a proper CY3, put
\[
 L(F,G)=\det R\operatorname{Hom}_X(F,G),\qquad
 \mathscr D(F)=\det R\operatorname{Hom}_X(F,F).
\]
Serre duality identifies
\(R\operatorname{Hom}(G,F)\simeq
R\operatorname{Hom}(F,G)^\vee[-3]\).  Both the dual and the odd shift
invert the determinant, so their composite preserves it.

\begin{proposition}[The cross-determinant square]
\label{prop:orientation-direct-sum-multiplicativity}
\label{prop:orientation-extension-multiplicativity}
For a direct sum, and for an exact triangle represented by a strict
filtered perfect model with subquotients \(F,G\), determinant
functoriality gives
\[
 \mathscr D(E)\simeq
 \mathscr D(F)\otimes\mathscr D(G)\otimes L(F,G)^{\otimes2}.
\]
For a strict finite flag with subquotients \(F_1,\ldots,F_r\),
\[
 \mathscr D(E)\simeq
 \bigotimes_i\mathscr D(F_i)\otimes
 \bigotimes_{i<j}L(F_i,F_j)^{\otimes2}.
\]
These determinant identifications are compatible with regrouping the
flag, with the determinant signs fixed by its order.
\end{proposition}
\begin{proof}
The direct-sum endomorphism complex has its four Hom blocks.  The two
cross blocks have equal determinants by CY3 duality.  For a filtered
model, filter the endomorphism complex by the difference of the flag
indices.  Its associated graded consists of all the Hom blocks.
Multiplicativity of determinant on this filtration gives the same
formula without splitting the extension.  Refining or regrouping the
filtration gives the canonical determinant comparison for the same
filtered complex.  Pair the off-diagonal blocks using Serre duality
in their fixed order.
\end{proof}

A trace-free or reduced construction changes this formula.  For
specified reduced determinant lines define the excess line
\[
 A(F,G)=\mathscr D^{\mathrm{red}}(E)\otimes
 \bigl(\mathscr D^{\mathrm{red}}(F)\otimes
 \mathscr D^{\mathrm{red}}(G)\otimes L(F,G)^2\bigr)^{-1}.
\]
One must compute \(A(F,G)\) and supply its root if a square-root
multiplication uses it.  Reducing the two endpoints independently
need not give the reduction of their extension.

For a concrete trace calculation take
\[
 H=\mathbb C[u]/(u^2)\otimes\Lambda(\eta),\quad |u|=2,
 \quad|\eta|=1,\quad\tau(u\eta)=1.
\]
In the trace-free block of sizes \(n,m>0\), the common scalar
complement is
\[
 t\longmapsto
 \operatorname{diag}(tI_n/n,-tI_m/m).
\]
Its trace pairing is \((n+m)\tau(tt')/(nm)\).  Deleting each block
trace separately would delete this nonzero paired summand.  For a
flag of block sizes \(n_i\), the same complement is
\(\{(t_i):\sum_i t_i=0\}\), with pairing
\(\sum_i\tau(t_it_i')/n_i\).  Regrouping a flag retains this one
common trace condition.

This calculation also has nonsplit families.  Over \(\mathbb C[t]\),
take the upper triangular trace-free matrices spanned by
\(h=\operatorname{diag}(1,-1)\) and \(e=e_{12}\), tensored with
\(H\), and differential \([t\eta e,-]\).  Then
\[
 dh=-2t\eta e,\qquad de=0.
\]
The cohomology contains \(\mathbb C[t]/(t)\) classes represented by
\(\eta e\) and \(u\eta e\), in degrees one and three.  Replacing
the family by a split fibre loses this torsion and its derived
specialization.  Determinants of the actual filtered complex remain
well defined throughout the family.

\begin{definition}[Strong reduced orientation]
\label{def:O1-strong-reduced-orientation}
On a specified reduced moduli problem, a \emph{strong reduced orientation}
consists of a root pair \((o,\rho)\), compatible equivariant descent, and
extension isomorphisms
\[
 \gamma_{F,G}:o_E\xrightarrow{\sim}
 o_F\otimes o_G\otimes L(F,G)\otimes a(F,G),
 \qquad a(F,G)^2\simeq A(F,G).
\]
Their squares are the prescribed determinant comparisons.  On every
three-step flag, the two composites of \(\gamma\), including the
cross and excess lines, agree.  They agree with the chosen unit on
the zero object and commute with the stated transitions and base
changes.  These are equations of line isomorphisms over the full
bases, including their nilpotents.
\end{definition}

\begin{proposition}[Coordinate roots on strict flags]
\label{prop:orientation-coordinate-flag}
For a strict flag in the K3 formal endomorphism category, with the
Serre-compatible determinant comparisons above, the coordinate
factors satisfy
\[
 d_f(E,E)=d_f(F,F)d_f(G,G)d_f(F,G)^2,
 \qquad
 j_f(E)=j_f(F)j_f(G)d_f(F,G).
\]
The second formula uses the specified Artin roots and extends to
all finite flags by multiplication over their ordered pairs.
\end{proposition}
\begin{proof}
The upper triangular endomorphism acts on the filtered Hom complex;
its associated graded has the four Hom blocks.  Determinants are
products on that filtration.  The divided-difference operators on
opposite cross blocks are adjoint for the K3 Serre pairing, so their
determinants agree.  This proves the first formula.  The Euler
characteristics obey
\(\chi(E,E)=\chi(F,F)+\chi(G,G)+2\chi(F,G)\).
Consequently both sides of the root formula have the same square
and the same reduced value \(\bar c_1^{-\chi(E,E)/2}\).  Uniqueness of
the root in \(1+\mathfrak m_R\) proves equality on the full Artin
base.  Refinement of the same filtration proves the all-flag formula.
\end{proof}

Even squares and units do not force the three-step equation.  Starting
with compatible maps indexed by positive masses, multiply
\(\gamma_{2,1}\) by \(-1\) and leave the other maps unchanged,
including the unit maps.  The two composites at masses \((1,1,1)\)
then differ by \(-1\), while all squares are unchanged.
There are nevertheless two useful positive corrections.  A normalized
mass-only sign three-cocycle \(c(a,b,c)\) on \(\mathbb N\) is the
coboundary of
\[
 \beta(a,b)=\prod_{j=0}^{b-1}c(a,j,1).
\]
The cocycle equation telescopes to \(\delta\beta=c\), so rescaling
by \(\beta\) removes this mass-only associator.  For a normalized
constant sign two-cocycle \(\alpha\) on a finite group \(G\) of odd order,
\(b(g)=\prod_{h\in G}\alpha(g,h)\) similarly satisfies
\(\delta b=\alpha^{|G|}=\alpha\).  Neither calculation trivializes
a sign cocycle varying over the moduli space.  For example,
\(\alpha((a,b),(c,d))=(-1)^{ad}\) on \(C_2^2\) is not a
coboundary: it is nonsymmetric, whereas a coboundary on an abelian
group with trivial coefficients is symmetric.

\subsubsection*{Oriented coefficients and Hall operations}

On an oriented critical moduli scheme, the pointwise Euler
characteristic of the vanishing-cycle sheaf is the Behrend function
\(\nu\); the sign local system has rank one and does not change this
Euler characteristic.  This is the numerical comparison of
\cite[Remark~6.14]{BBDJS2015}, underlying the weighted-Euler theories
of \cite{BEHREND,JoyceSong}.  Passing to \(\nu\) forgets monodromy
and extension morphisms.  Those data are still needed for a product
on critical cohomology.


\begin{proposition}[Thom--Sebastiani and quotient descent]
\label{prop:orientation-thom-sebastiani-compatibility}
\label{prop:orientation-e-quotient-descent}
Compatible strong orientation data identify the sign torsors in
Thom--Sebastiani comparisons and descend the oriented coefficient
sheaves to the quotient stack.  These comparisons use the actual
critical charts and their embeddings.  They do not follow from an
isomorphism between determinant classes in a Picard group alone.
\end{proposition}
\begin{proof}
On a critical chart the coefficient is twisted by the principal
bundle of square-root identifications.  An orientation isomorphism
identifies these bundles; equality on triple overlaps gives equality
of their descent maps.  Apply the Thom--Sebastiani isomorphism to the
vanishing-cycle factor and tensor its sign-bundle comparison.
Equivariance supplies the same equations on the action presentation.
This constructs descent to the stack, where stabilizer characters
are retained.  Coarse descent would impose the additional
inertia-trivial condition discussed above.
\end{proof}

\begin{proposition}[Operations require their coefficient morphisms]
\label{prop:orientation-descent-hall-product}
\label{prop:orientation-descent-hall-coproduct}
\label{prop:finite-dcritical-hall-product}
For an \(E\)-equivariant moduli space \(M\), put
\(\mathcal M=[M/E]\).  Suppose an extension correspondence
\[
 \mathcal M\times_{BE}\mathcal M\xleftarrow{p}\mathcal E
 \xrightarrow{q}\mathcal M
\]
has a specified critical coefficient comparison, with its shifts,
support conditions and functors, giving
\(q_!p^*(\mathcal P\boxtimes_{BE}\mathcal P)\to\mathcal P\).
Fix a cohomology theory in which the specified pull--push and
K\"unneth maps exist, so that this comparison induces the Hall product.
Assume the two resulting comparisons agree on three-step flags and
satisfy the relevant base-change identities.  Then the induced Hall
product is associative.  A collision coproduct requires a separately
specified correspondence and coefficient morphism, with its own
coassociativity and support identities.  The orientation and the
product alone do not supply it.
\end{proposition}
\begin{proof}
Both parenthesizations of the product pull back to the complete
three-step flag correspondence.  Proper or supported base change,
projection and the assumed coefficient equality identify their
pushforwards.  The same argument proves coassociativity only after
the collision coefficient morphism and its complete-flag equality
have been supplied.  Inverting Thom--Sebastiani on a local product
chart does not construct this global morphism on collision support.
\end{proof}

The tensor is over \(BE\) because both endpoints of an extension use
one elliptic torsor.  Quotient the whole correspondence with its
diagonal action before separating its maps.  The zero object retains
its map to \(BE\); it is not replaced by a point before checking the
unit.  The common-torsor equivalence of
Section~\ref{subsec:hall-correspondences} fixes this convention independently
of any reduction of the deformation complex.

\begin{proposition}[Primitive projections after descent]
\label{prop:orientation-descent-primitive-projection}
Let a descended coefficient theory already carry the specified
connected conilpotent cocommutative Hopf structure over a
characteristic-zero field.  The convolution logarithm of the identity
is then its projection to the kernel of the reduced coproduct on
the domain proved in
Section~\ref{subsec:primitive-recognition}.  If transition functors
preserve the Hopf maps and their finite convolution sums, they
preserve this projection.  Without these hypotheses, a chosen
splitting of the primitive kernel is additional data.
\end{proposition}
\begin{proof}
The formula for the projection is a finite sum on every conilpotence
stage and is natural for the given Hopf maps.  Apply the transition
functor term by term.  An orientation comparison by itself asserts
none of the coalgebra, connectedness or exactness hypotheses needed
for that formula.
\end{proof}

\subsubsection*{Weyl transport as a root-groupoid problem}

Let \(W=W^{(2)}(\Lambda_{II}^{2,1})\).  Its three simple
reflections have the presentation described in
Lemma~\ref{lem:bkm-chamber-s3}.  To transport an orientation, one
first needs geometric maps or correspondences between the relevant
strata; the lattice reflection alone gives no pullback of a line.
For clarity, write \(g^*\) for the specified pullback functor, with
its composition identifications.

\begin{definition}[The groupoid of defined transports]
\label{def:finite-partial-action-groupoid-GR}
Choose strata and specified invertible geometric transports between
them.  Let \(\mathcal G\) be their groupoid, with all defined
composites and their relations.  A finite set of strata or generating
arrows need not make \(\mathcal G\) finite: there may be infinitely
many distinct loops.  A finite cochain calculation applies only after
a finite model for the cochains in the required degrees has been
specified.
\end{definition}

\begin{definition}[Root transport and its sign cocycle]
\label{def:O1-plus-weyl-transport}
\label{def:orientation-projective-cocycle-coR}
Suppose each stratum \(U\) has a root pair \((o_U,\rho_U)\), and
each arrow \(g:U\to V\) has a determinant comparison
\(d_g:\mathscr D_U\to g^*\mathscr D_V\).
A lift is an isomorphism
\[
 \tau_g:o_U\longrightarrow g^*o_V,
 \qquad (g^*\rho_V)\tau_g^{\otimes2}=d_g\rho_U.
\]
When the determinant comparisons are coherent, two composites of such
lifts differ by a section of \(\mu_2\).  In particular define
\[
 c(g,h)=\tau_{hg}^{-1}\,(g^*\tau_h)\,\tau_g
 \in H^0(U,\mu_2).
\]
This is a two-cocycle on \(\mathcal G\) with the pullback coefficient
system \(U\mapsto H^0(U,\mu_2)\).  Strong Weyl transport, denoted
\((O1)^+\), is a coherent lift preserving the full root, extension,
unit and equivariant descent data.
\end{definition}

\begin{proposition}[Existence, torsor defects and strictification]
\label{prop:weyl-wall-transport-tau-construction}
\label{prop:weyl-wall-transport-tau-square}
\label{prop:orientation-cocycle-class-vanishing-criterion}
\label{prop:orientation-cocycle-cochain-trivialization-criterion}
\label{prop:weyl-orientation-torsor-defect-computation}
\label{prop:weyl-orientation-torsor-defect-vanishing}
Given \(d_g\), the obstruction to identifying the two root pairs is
the class of their ratio in \(H^1(U,\mu_2)\).  If it vanishes, lifts
form a torsor under \(H^0(U,\mu_2)\).  Once lifts exist, they can be
made coherent precisely when \([c]=0\).  Multiplying lifts by a
one-cochain \(b\) changes \(c\) by \(\delta b\).
For an involution \(s\), coherent normalized transport satisfies
\[
 (s^*\tau_s)\tau_s=\mathrm{id}.
\]
The determinant square equation alone does not impose this identity.
\end{proposition}
\begin{proof}
Use \(d_g\) to regard \(g^*o_V\) and \(o_U\) as roots of one line.
Their ratio is a \(\mu_2\)-torsor.  A trivialization is the required
lift, and its automorphisms give the stated freedom.  Associativity
of three lifted arrows proves the cocycle equation.  Substitution of
\(b(g)\tau_g\) gives the coboundary formula, so strictification is
equivalent to \([c]=0\).  The pair of arrows \((s,s)\) gives the
involutive relation.  Before strictification their composite may
instead be the nontrivial automorphism of the root.
\end{proof}

For constant coefficients \(\mathbb F_2\), a specified finite
cochain model makes the equation \(\delta b=c\) a linear system.
It is soluble if and only if adjoining \(c\) to the matrix of
\(\delta:C^1\to C^2\) does not increase its rank.  One must use
cochains for the actual groupoid.  A transitive action groupoid is
equivalent to the stabilizer group of one object.  Its cohomology is
therefore stabilizer cohomology, not automatically \(H^*(W)\).
Only an action on a single fixed object gives \(BW\).

\begin{proposition}[Naturality of the obstruction data]
\label{prop:weyl-alpha-red-preservation-criterion}
\label{prop:weyl-alpha-e-free-preservation-criterion}
\label{prop:weyl-finite-stabilizer-beta-preservation-criterion}
\label{prop:weyl-finite-stabilizer-lambda-zero-preservation-criterion}
A determinant comparison transports the Kummer class by pullback.
A root lift transports its chosen nulltrivialization.  At a
stabilizer isomorphism \(\varphi:H_U\to H_V\), the equations are
\[
 \beta_U=\varphi^*\beta_V,
 \qquad \lambda_U=\varphi^*\lambda_V.
\]
If inertia-trivial normalization is imposed, it is preserved by the
same equation with \(\lambda_U=\lambda_V=0\).  The base-filtration
part of a Borel class is preserved as a filtration class, without
choosing representatives in \(H^2(BE)\).
\end{proposition}
\begin{proof}
Kummer boundaries, restrictions and character pullbacks commute with
the geometric pullback.  A root isomorphism gives the same comparison
before passing to cohomology, hence transports the nulltrivialization
as well.  In chosen cohomology bases, if \(A_\varphi\) represents
\(\varphi^*\), the defect is \(\beta_U-A_\varphi\beta_V\), and
similarly in degree one.  This fixes the variance of the finite
matrix equation.  Naturality of the spectral sequence preserves its
filtration and quotient terms; it does not create a splitting.
\end{proof}

\begin{proposition}[Compatibility of strictifications]
\label{prop:orientation-cocycle-cochain-transition-compatibility}
Suppose transition pullbacks \(r\) identify two systems of
projective lifts with cocycles \(c_R,c_{R'}\).  Cochains killing
these cocycles give compatible strictifications when their
restrictions differ by the specified change of trivialization.  In
constant additive coefficients this condition is
\[
 r(b_R)-b_{R'}=\delta a_{RR'},
\]
with compatible zero-cochains \(a_{RR'}\) on successive transitions.
\end{proposition}
\begin{proof}
Conjugating the root lift at each object by its sign \(a\) changes
the arrow cochain by \(\delta a\).  Thus the displayed equation
is exactly a natural isomorphism of the two strict transports.
Composition of these natural isomorphisms is the additional
compatibility equation on a triple of transition levels.
\end{proof}

\subsubsection*{Reflection characters and divisor monodromy}

\begin{definition}[A character extracted from root transport]
\label{def:orientation-character-epsilon-o}
To obtain a scalar character from \((O1)^+\), choose compatible
reference identifications of the transported lines, or a fixed
one-dimensional fibre representation.  Suppose the residual signs
\(\sigma_i\in\{\pm1\}\) of the three simple lifts are constant on
the relevant components and compatible with transitions.  They
define
\[
 \epsilon_o(s_{i_1}\cdots s_{i_k})=\prod_{j=1}^k\sigma_{i_j}.
\]
For the present type-II group, the only Coxeter relations are
\(s_i^2=1\); hence this is a character once the residual signs have
been constructed.  An action on moving lines, without the reference
identifications, is not itself such a scalar character.
\end{definition}

\begin{proposition}[A simple Pfaffian wall]
\label{prop:orientation-character-simple-wall-sign-criterion}
Suppose a local Pfaffian section is \(p=\upsilon u e\), where
\(e\) is a line frame and \(\upsilon\) is invertible.  Fix a reference
comparison of source and target frames.  Let the base reflection
send \(u\mapsto-u\) and preserve \(\upsilon\).  Relative to that
reference comparison the pulled-back section has ratio \(-1\).
If the actual root lift satisfies
\[
 \tau_s(p)=s^*p,
\]
its residual sign relative to the reference comparison is \(-1\).
Under the component and transition compatibility in the preceding
definition, this gives \(\epsilon_o(s_i)=-1\).
\end{proposition}
\begin{proof}
In the reference frames, write \(\tau_s(e)=\lambda s^*e\).
The displayed compatibility gives
\(\lambda\upsilon u=-\upsilon u\) on the wall complement, hence
\(\lambda=-1\).  Squaring loses the sign.  Preserving a reference
frame alone would not identify its action with this section-compatible
root lift.  The calculation uses the stipulated reflection, not
parallel transport along a real path through a removed wall.
\end{proof}

\begin{proposition}[The determinant and automorphic characters]
\label{prop:orientation-character-determinant-criterion}
\label{prop:orientation-character-maass-criterion}
If the source character exists and takes value \(-1\) on all three
simple type-II reflections, then
\[
 \epsilon_o=\det_W
 =\widetilde\nu_{\Delta_5}|_W.
\]
\end{proposition}
\begin{proof}
The determinant of a real reflection is \(-1\), so equality on the
generators gives the first equality.  The second is the homogeneous
orthogonal character calculation in
Lemma~\ref{lem:bkm-maass-restriction}.  The homogeneous character is
used here: its value on an orthogonal reflection is not the bare
symplectic Maass character of a chosen lift.  This automorphic
identity does not construct the source character.
\end{proof}

\begin{proposition}[Chamber symmetries]
\label{prop:s3-chamber-maass-triviality}
\label{prop:orientation-character-s3-extension-decision}
The homogeneous character \(\widetilde\nu_{\Delta_5}\) is trivial
on the type-I chamber symmetry group \(\Sigma\simeq S_3\).
A character \(\epsilon_o\) taking value \(-1\) on all three simple
reflections is \(\Sigma\)-invariant.  Its abstract extensions to
\(W\rtimes\Sigma\) are exactly
\[
 \widetilde\epsilon_\chi(w,\sigma)=\epsilon_o(w)\chi(\sigma),
 \qquad \chi\in\{1,\operatorname{sgn}\}.
\]
The automorphically compatible extension has \(\chi=1\).
\end{proposition}
\begin{proof}
The type-I transpositions have homogeneous character \(+1\), by the
calculation in Chapter~\ref{ch:view1-automorphic}; they generate
\(\Sigma\).  Chamber symmetry permutes the three equal generator
values of \(\epsilon_o\), proving invariance.  The displayed formula
then respects the semidirect product.  Two extensions differ by a
character of \(S_3\), whose only sign characters are the two listed.
\end{proof}

\begin{proposition}[Geometric extension to chamber symmetries]
\label{prop:semidirect-hall-lift-construction-criterion}
An abstract extension is realized on an oriented Hall theory only
after supplying chamber transports of its actual correspondences,
coefficient maps and root pairs.  These must preserve the specified
product, coproduct, unit, pairing and transitions, and satisfy the
\(S_3\) relations and
\[
 \theta_\sigma\tau_i\theta_\sigma^{-1}
 =\tau_{\sigma(i)}
\]
with the appropriate pullbacks.  Such coherent transports define
an action of \(W\rtimes S_3\).
\end{proposition}
\begin{proof}
The relations present the semidirect product.  An action on the
geometric objects and all their comparison maps satisfying these
relations factors through that presentation.  The character criterion
for a line does not assert the existence of these correspondences.
\end{proof}

\begin{definition}[A local divisor monodromy datum]
\label{def:local-divisor-monodromy-datum}
Let \(\mathcal U\) be a complex analytic family near a smooth point
of a divisor, with a specified transverse disk and punctured disk
\(\Delta^*\).  Choose its reduced critical recipient and orientation.
Suppose the relevant supported cohomology of the oriented
vanishing-cycle sheaf forms a finite-rank local system \(\mathbb V\)
on \(\Delta^*\).  A base point and the positive meridian then define
\[
 T_D:\mathbb V_{z_0}\xrightarrow{\sim}\mathbb V_{z_0}
\]
up to conjugacy.  For a rank-one sign local system its value lies in
\(\{\pm1\}\).  The required local constancy, supports and
finite-dimensionality are hypotheses on the actual family.
\end{definition}

\begin{proposition}[Divisor and Weyl monodromy have different domains]
\label{prop:imaginary-divisor-not-weyl-monodromy}
\label{prop:local-divisor-monodromy-maass-irrelevance}
An imaginary simple root contributes a multiplicity to the BKM
presentation but no simple Weyl reflection.  A meridian around a
geometric divisor labelled by that root therefore has no prescribed
value under \(\epsilon_o:W\to\{\pm1\}\).  Its local monodromy
requires the preceding analytic datum; a Maass character or a root
multiplicity alone does not supply it.
\end{proposition}
\begin{proof}
The Weyl group is generated by reflections in the real simple roots.
The meridian instead belongs to the fundamental group of a specified
punctured analytic neighbourhood.  Identifying their actions would
require a map of these geometric actions and their coefficient
systems, not merely a common root label.
\end{proof}

\subsubsection*{Parity, duality and cyclic traces}

\begin{proposition}[Local root profiles and orientation lifts]
\label{prop:three-simple-wall-chart-orientation}
\label{prop:isotropic-aij-chart-orientation}
\label{prop:delta123-chart-orientation}
\label{prop:orientation-line-delta123-parity-compatibility}
The target root profiles are \((1\mid0)\) for each real simple root,
\((10\mid0)\) for each isotropic \(a_{ij}\), and \((29\mid93)\)
for \(\delta_1+\delta_2+\delta_3\), with the parity convention of
Chapter~\ref{ch:view2-bkm}.  A compact chart realizes these profiles only
if an actual graded comparison identifies its relevant cohomology
with those spaces.  Their parity does not alone determine a lift
to a determinant square root.
\end{proposition}
\begin{proof}
The profiles are the target calculations in the BKM chapter.  A
graded linear isomorphism preserves them, so it gives the conditional
chart assertion.  For the last statement, consider one odd line.
Parity acts by \(-1\) on its Berezinian.  A scalar lift \(\lambda\)
to a square root would have to satisfy \(\lambda^2=-1\), hence
\(\lambda=\pm i\) over \(\mathbb C\).  Such a lift is not an
involutive sign action.  More generally the Berezinian parity is
\((-1)^q\) on a space of dimension \((p\mid q)\); a Pfaffian-half
comparison is additional structure.  In particular \(q=93\) cannot
be replaced by an assumed orientation sign without that comparison.
\end{proof}

\begin{proposition}[Opposite degrees under CY3 duality]
\label{prop:orientation-line-negative-root-compatibility}
If a specified comparison identifies
\(K_{-\gamma}\simeq K_\gamma^\vee[-3]\), then
\[
 \det K_{-\gamma}\simeq\det K_\gamma.
\]
An unshifted dual would instead have inverse determinant.  The
corresponding square-root identification must square to the actual
shifted determinant comparison.
\end{proposition}
\begin{proof}
Dualization and an odd cohomological shift each invert determinant.
Their composite gives the displayed equality.  The distinction is
visible for \(K_\gamma=\mathcal O(1)\) in degree zero on
\(\mathbb P^1\): the shifted dual has determinant \(\mathcal O(1)\),
whereas the unshifted dual has determinant \(\mathcal O(-1)\).
\end{proof}

\begin{proposition}[Opposite products and Chevalley transport]
\label{prop:orientation-chevalley-antiinvolution}
Suppose geometric duality defines an involution \(\vartheta\) on a
specified graded associative Hall theory, with
\[
 \vartheta(xy)=(-1)^{|x||y|}\vartheta(y)\vartheta(x),
 \qquad \vartheta^2=1,
\]
and compatible orientation, coproduct and transition comparisons.
It then induces the corresponding opposite-product involution on
its primitive Lie superalgebra.  Agreement with a target Chevalley
involution is an additional equality of these maps under the
primitive comparison.
\end{proposition}
\begin{proof}
Substitute the displayed formula into the supercommutator.  One
obtains
\(\vartheta([x,y])=-[\vartheta(x),\vartheta(y)]\).
Equivalently it is the indicated Lie anti-involution, with a minus
sign converting it to a Lie automorphism if that convention is used.
The coproduct comparison preserves primitives, and the determinant
square comparison preserves their orientation data.  None of these
maps is determined by dimensions of opposite root spaces.
\end{proof}

\begin{proposition}[Associative traces and primitive forms]
\label{prop:orientation-frobenius-pairing}
\label{prop:orientation-frobenius-trace-degree}
\label{prop:orientation-cy3-serre-signs}
Let \(A\) be a graded associative algebra or dg algebra, \(\ell\) a
line, and \(\operatorname{tr}:A\to\ell[-3]\) a closed cyclic trace.
Then
\[
 B(x,y)=\operatorname{tr}(xy),\qquad
 B(x,y)=(-1)^{|x||y|}B(y,x),\qquad
 B(xy,z)=B(x,yz).
\]
Its nonzero homogeneous values have \(|x|+|y|=3\), when \(\ell\)
is in degree zero.  On the commutator Lie superalgebra the restricted
form is invariant.  This is a different condition from Hopf
adjointness on a full bialgebra.
\end{proposition}
\begin{proof}
These are respectively the definition, graded cyclicity, associativity
and degree of the trace.  Expanding the commutators gives Lie
invariance with the same Koszul signs.  In contrast, for an ordinary
Hopf pairing and primitive elements \(x,y,z\), Hopf adjointness gives
\[
 \langle xy,z\rangle
 =\langle x,z\rangle\epsilon(y)+\epsilon(x)\langle y,z\rangle=0.
\]
Thus it also annihilates \([x,y]\) against \(z\).  An invariant Lie
form on a nonabelian primitive algebra need not do this.  A product
and a trace must be defined on \(A\), not on a primitive subspace
which is generally not closed under that product.
\end{proof}

\begin{proposition}[The induced quotient pairing]
\label{prop:orientation-cy3-pairing-nondegeneracy}
Let a bilinear pairing of finite-dimensional spaces have matrix
\(G\), with convention \(B(x,y)=x^tGy\).  Its left and right
radicals are \(\ker G^t\) and \(\ker G\).  It induces a
nondegenerate pairing on the two radical quotients.  If \(Q_L,Q_R\)
are quotient maps and \(s_L,s_R\) are sections, its matrix is
\[
 \overline G=s_L^tGs_R,\qquad Q_Ls_L=1,\quad Q_Rs_R=1.
\]
The expression is independent of the sections.  It is generally not
\(Q_LGQ_R^t\).
\end{proposition}
\begin{proof}
Changing a section adds a map into the appropriate radical, which
pairs to zero.  If a quotient class pairs to zero with every opposite
class, a representative lies in the radical and its quotient is
zero.  This proves nondegeneracy over a field.  Over a coefficient
ring one must separately require finite projective quotients and
that the resulting map to the dual is an isomorphism; a nonzero
determinant or full generic rank is insufficient.  Compatibility
with a Lie bracket, a Hopf operation or transitions further requires
the corresponding ideal, coideal or kernel identities.
\end{proof}

\subsubsection*{Compatible towers and effective orientation lines}

\begin{definition}[Orientation-preserving transitions]
\label{def:transition-orientation-preservation}
For specified transition maps \(t_{RR'}\), an orientation transition
is a root isomorphism
\[
 \theta^o_{RR'}:t_{RR'}^*o_{R'}\xrightarrow{\sim}o_R.
\]
Its square is the determinant comparison, and it preserves the
chosen equivariant, extension, unit and Weyl identifications.
Identity and composition are equations of these isomorphisms on
all successive transitions.
\end{definition}

\begin{proposition}[Formal orientation towers]
\label{prop:transition-orientation-preservation-criterion}
\label{prop:orientation-class-inverse-limit-survival}
A compatible family of orientation objects and transition
isomorphisms defines an object in the homotopy limit of their
groupoids.  Every finite projection is the specified root pair,
with its comparisons.  This formal statement does not require the
orientation line itself to have trivial Picard class.
\end{proposition}
\begin{proof}
An object of the homotopy limit is precisely such a compatible
system, including its composition coherence.  Applying any finite
projection recovers its component.  Vanishing of the Kummer
obstruction says that the determinant has a root; it does not say
that the root is a trivial line.
\end{proof}

\begin{proposition}[Classes in a tower of root groupoids]
\label{prop:orientation-no-new-limit-class}
Consider a countable tower and a compatible family of isomorphism
classes of roots.  Choices of compatible representatives and edge
isomorphisms have, when nonempty, ambiguity measured by
\(\varprojlim{}^1 A_R\), where \(A_R\) is the automorphism group
of the chosen root at level \(R\).  Here
\(A_R=H^0(B_R,\mu_2)\), with the induced transition maps.
If this automorphism tower is Mittag--Leffler, the ambiguity vanishes.
To obtain a line on a geometric limit, one additionally needs an
effective comparison between its root groupoid and the homotopy
limit.  An equivalence of those groupoids, together with the stated
vanishing, identifies its isomorphism classes with the compatible
finite classes.
\end{proposition}
\begin{proof}
Fix representatives.  Changing each adjacent transition isomorphism
is an element of \(\prod_R A_R\); changing the representative
identifications acts by
\((a_R)\mapsto(a_R-t_Ra_{R+1})\) in additive notation.  The quotient
is the usual \(\varprojlim{}^1 A_R\).  The Mittag--Leffler condition
makes this quotient zero.  This is an automorphism calculation;
there is no presumed abelian group of roots of a fixed line.
Finally an effective comparison transfers the formal objects and
isomorphisms to the geometric limit.  Without it, a formal compatible
system is all that has been constructed.
\end{proof}

\begin{definition}[Coefficient-preserving transitions]
\label{def:transition-vanishing-cycle-preservation}
A transition of oriented critical coefficients is a specified
isomorphism between the chosen pullback of \(\mathcal P_{R'}\)
and \(\mathcal P_R\), with the required shift and support convention.
It commutes with the orientation sign torsors, critical-chart
comparisons, duality and the coefficient maps defining the retained
Hall operations.  These isomorphisms satisfy the transition
composition equation.
\end{definition}

\begin{proposition}[Preservation of critical coefficients]
\label{prop:transition-vanishing-cycle-preservation-criterion}
Coefficient-preserving transitions identify the induced finite
operations and their orientation comparisons.  Passing to a chosen
completed operation further requires the limit and completed-tensor
hypotheses for that coefficient category.
\end{proposition}
\begin{proof}
Apply the specified transition isomorphisms to the defining
pull--push diagram.  Their commutation with its coefficient maps
identifies the two operations.  A limiting diagram follows only if
the relevant functors, supports and tensor products commute with
that limit.  Properness of the underlying stack map alone does not
assert this coefficient comparison.
\end{proof}

The orientation theorems for sheaf moduli and real elliptic operators
\cite{JoyceUpmeier,JoyceTanakaUpmeier} provide substantial orientation
constructions in their own settings.  Applying them to the compact
reduced problem requires the comparison with its actual reduced
complex and extension data.  The formal coordinate functors, cross
determinant squares and finite character calculations above supply
explicit parts of that problem.  The remaining global question is
the construction of one reduced oriented coefficient system on the
common elliptic quotient, compatible with complete flags, geometric
transitions and the stipulated reflection transports.


\subsection{Critical Hall operations and primitive Lie algebras}
\label{subsec:hall-correspondences}

The formal deformation model of
\S\ref{subsec:hybrid-ran-carrier} supplies actual extension
categories and polynomial critical models. Their Hall algebra
requires a coefficient theory on the extension stacks.
The global critical Hall theorem gives such a multiplication on
ordinary critical cohomology under its strong orientation and
properness hypotheses. The compact reduced $S\times E$ problem
asks for additional comparisons with the reduced quotient
coefficients. A collision coproduct and a primitive Lie form are
further structures, with separate equations.

We work over a characteristic-zero field $k$, using $k=\Q$ for
constructible coefficients before extension of scalars. All
cohomological shifts and orientation parities are part of the
coefficient objects. When an ordinary superalgebra is used below,
its multiplication is even and its tensor symmetry is
$\tau(x\otimes y)=(-1)^{|x||y|}y\otimes x$.
Passing from an index-shifted critical Hall multiplication to this
category requires the stated regrading or comparison.

\subsubsection*{The coefficient object and its multiplication}

\begin{definition}[A bounded geometric coefficient object]
\label{def:finite-geometric-hall-object}
Fix chosen moduli components $\mathcal M_{R,\gamma}$ and critical
coefficients $K_{R,\gamma}$, including any reduction and orientation
data. Choose one realization, for example
\[
 \mathcal H_{R,\gamma}=R\Gamma(\mathcal M_{R,\gamma},K_{R,\gamma}),
 \qquad
 \mathcal H_R=\bigoplus_{\gamma\in\Gamma_R}\mathcal H_{R,\gamma}.
\]
If instead compact-support or dual compact-support realization is
used, name it explicitly and use its coefficient variance. In
particular
$\mathbb B(M;K)=R\operatorname{Hom}(R\Gamma_c(M,K),k)$
is contravariant in $K$.
The set $\Gamma_R$ is finite when a finite charge window is
specified. Finite-dimensional cohomology or a perfect chain
model is an additional hypothesis, not a consequence of this
finite indexing set.
\end{definition}

For a bounded ordinary-cohomology model, let
$\mathcal E_{\alpha,\beta}$ be the ordered extension stack with
maps $p$ to its endpoints and $q$ to its middle object. The
required coefficient morphism is
\[
 u_{\alpha,\beta}:
 p^*(K_\alpha\boxtimes K_\beta)[\kappa(\alpha,\beta)]
 \longrightarrow q^!K_{\alpha+\beta}.
\]
The source shift $\kappa$ is fixed by this formula, as in
\eqref{eq:D-critical-hall-comparison}. On a translation quotient,
the endpoint stack is the fibre product over $BE$ and the
coefficient tensor is the specified relative tensor.

\begin{definition}[A geometric Hall product]
\label{def:compact-hall-product}
Put $L=p^*(K_\alpha\boxtimes K_\beta)[\kappa(\alpha,\beta)]$.
Suppose $u:L\to q^!K_{\alpha+\beta}$ is defined, the support
of $L$ is proper over the target, and K\"unneth holds in
the chosen realization. The Hall product is the composite
\[
 (\mathcal H_\alpha\otimes\mathcal H_\beta)
       [\kappa(\alpha,\beta)]
 \longrightarrow R\Gamma(\mathcal M_{\alpha+\beta},Rq_!L)
 \longrightarrow\mathcal H_{\alpha+\beta}.
\]
The first arrow is pullback followed by $Rq_*L\simeq Rq_!L$.
The second applies $Rq_!u$ and the counit $Rq_!q^!\to1$.
Only permitted charge sums are included at a bounded stage.
\end{definition}

\begin{proposition}[The pullback and integration legs]
\label{prop:finite-hall-product-pullback-existence}
\label{prop:finite-hall-product-pushforward-existence}
The product is defined under the stated coefficient, K\"unneth
and proper-support hypotheses. Existence of $Rq_!$ for a
nonproper map alone does not provide the last arrow on ordinary
cohomology. For compactly supported cohomology, pullback along
$p$ needs its corresponding support condition instead.
\end{proposition}

\begin{proof}
Ordinary cohomology admits pullback through the unit
$K\to Rp_*p^*K$. Proper support identifies $Rq_*L$ with
$Rq_!L$, after which $Rq_!u$ and the trace give the last
arrow. Without that support condition an arbitrary global
section on the extension stack need not have proper support
over the target. Dually, pulling a
compactly supported section back along a map with nonproper
fibres need not leave compact support.
\end{proof}

\begin{proposition}[Independence of an exceptional-pushforward model]
\label{prop:finite-hall-product-qbang-compactification-independence}
Suppose a map in the coefficient theory is compactifiable as
$q=\overline qj$, with $j$ open and $\overline q$ proper, and
exceptional pushforward is computed by
$Rq_!=R\overline q_*j_!$. Two such compactifications give the
same functor when the chosen six-functor theory supplies its
proper base-change and composition identifications.
\end{proposition}

\begin{proof}
Compare both compactifications with the closure of the graph of
the identity on their common open domain inside their product.
The proper comparison maps restrict to the identity on that
open. Proper base change and extension by zero identify the two
pushforwards with the one through this common refinement.
The composition constraints give the cocycle identification
for three choices. This concerns $q_!$; it does not remove a
proper-support requirement elsewhere in the Hall product.
\end{proof}

\begin{proposition}[Associativity of the coefficient product]
\label{prop:finite-hall-product-associativity-two-step-flags}
If the two composites of $u$ agree on the three-step extension
flag, including their shifts and orientation identifications,
then the Hall multiplication is associative on all permitted
triples. Coherent zero-endpoint comparisons give its unit.
\end{proposition}

\begin{proof}
The two geometric compositions are the same flag stack by
Proposition~\ref{prop:eight-two-step-flag-stacks}. Base change,
projection formula and composition identify their pull--push
functors. The coefficient equality gives equality of the
resulting maps. The two zero-endpoint sequences identify with
the identity correspondence and give the unit equations.
\end{proof}

Theorem~\ref{thm:hall-product} constructs these morphisms and
their coherence for the specified unreduced $3$-Calabi--Yau
moduli with strong orientation data. This applies to the
Ginzburg categories of the tripled-quiver models. On a compact
projective threefold with trivial canonical bundle, orientation
data exists by
\cite[Theorem~4.4]{JoyceUpmeier}; the strong associativity
condition remains additional. A formal cyclic equivalence
between deformation categories does not alone identify their
constructible coefficient maps on an open algebraic moduli
stack.

\subsubsection*{An actual constructible Hall calculation}

There is a concrete algebraic result before critical
coefficients are compared. Consider the Euler-characteristic
Hall algebra of finite-dimensional nilpotent representations of
the preprojective algebra of the quiver in
\eqref{hall:eq:preprojective-moment}. Functions are constructible,
and convolution integrates over invariant subspaces.
Use the same quotient-first convention as above.
Let $e_i$ be the characteristic function of the vertex simple.
The divided power $e_i^{(b)}=e_i^{*b}/b!$ is the characteristic
function of $S_i^{\oplus b}$: a complete flag in $\C^b$ has
Euler characteristic $b!$.

\begin{theorem}[The cubic real-root Serre relation]
\label{hall:thm:constructible-serre}
For distinct vertices $i,j$,
\[
 (\operatorname{ad}e_i)^3(e_j)=0.
\]
The assignment of the three Chevalley generators to $e_i$
therefore defines a homomorphism from the positive real-root
Kac--Moody Lie algebra with Euler matrix $G$ into this
constructible Hall algebra.
\end{theorem}

\begin{proof}
Only dimension vector $(3,1)$ on $i,j$ contributes. Combine
the two forward arrows into $X:\C^3\to\C^2$ and the reverse
arrows into $Y:\C^2\to\C^3$. The preprojective relation at
$i$ gives $YX=0$, hence
$\operatorname{rk}X+\operatorname{rk}Y\leq2$.
Put $I=\operatorname{im}Y$ and $K=\ker X$.

For $e_i^{(b)}*e_j*e_i^{(3-b)}$, a descending flag has its last
$i$-subspace $U$ of dimension $3-b$. Invariance requires
$I\subset U\subset K$. If $I\not\subset K$, all these flag
spaces are empty. Otherwise, writing $r=\dim I$ and $s=\dim K$,
their Euler characteristics are
$\binom{s-r}{3-b-r}$. We have
$s-r=3-\operatorname{rk}X-\operatorname{rk}Y\geq1$.
Consequently their alternating sum is
\[
 \sum_{b=0}^3(-1)^b\binom{s-r}{3-b-r}
 =(-1)^{3-r}(1-1)^{s-r}=0.
\]
Multiplying by $3!$ converts the divided-power relation into
the cubic commutator, up to its irrelevant overall sign.
The same argument with the two vertices exchanged proves the
opposite relation.
\end{proof}

This is a relation for actual convolution operators with
constructible-function coefficients. It proves neither
injectivity nor an identification of all imaginary generators.
Transporting it to critical cohomology or to compact reduced
sheaves needs a comparison preserving convolution. In particular,
matching the Euler matrix $G$ alone is insufficient.

\subsubsection*{A spectral coproduct on critical quiver cohomology}

The quiver theory also has a coproduct constructed from direct
sum and the ambient Euler class. Let $Q'$ be a finite quiver,
$T$ a torus acting on its arrows, and $W$ an invariant potential.
This includes the tripled quiver with its canonical potential.
For a dimension vector $d$, put $G_d=\prod_iGL_{d_i}$ and
\[
 R=H^*(BT;\Q),\qquad
 M_d=[\operatorname{Rep}_d(Q')/G_d],\qquad
 C_d=H_T^*(M_d,\varphi_{\operatorname{Tr}W}\Q),\qquad
 C=\bigoplus_d C_d.
\]
Assume the relative K\"unneth maps over $BT$ are isomorphisms
for the finite products used below. The direct sum has finite
dimension-vector support. Every fixed dimension vector has
finitely many decompositions.

Write $E_{i,d}$ for the tautological bundles and $L_a$ for the
line bundle on $BT$ associated to the character of arrow $a$.
The ambient two-term Hom complex, in degrees
zero and one, is
\[
 \Theta_{d,e}=\left[
 \bigoplus_i E_{i,d}^\vee\boxtimes E_{i,e}
 \longrightarrow
 \bigoplus_{a:i\to j}E_{i,d}^\vee\boxtimes E_{j,e}\otimes L_a
 \right].
\]
Its virtual rank is the quiver Euler form
$\chi(d,e)=\sum_i d_ie_i-\sum_{a:i\to j}d_ie_j$.
This complex belongs to the smooth representation stack; it
is not the threefold self-Ext complex of the critical locus.
Set $\Psi(V,z)=z^{\operatorname{rk}V}c(V,z^{-1})$, expanding
the Chern polynomial of a virtual bundle at $z=\infty$.

Direct sum induces a critical-cohomology pullback
$\delta_{d,e}:C_{d+e}\to C_d\otimes_R C_e$.
Tensoring the vertex spaces with a line gives a coaction
$\operatorname{act}_z^*=e^{z\mathsf T}$, polynomial on each
class. Thus $\mathsf T$ is locally nilpotent, $|z|=2$ and
$|\mathsf T|=-2$.

\begin{theorem}[Critical vertex coproduct]
\label{hall:thm:critical-vertex-coproduct}
The maps
\[
 \Delta_{d,e}(z)=\Psi(\Theta_{d,e},-z)
                     e^{z\mathsf T_1}\delta_{d,e}
 :C_{d+e}\longrightarrow(C_d\otimes_R C_e)((z^{-1}))
\]
satisfy
\begin{align*}
 (\Delta_{d,e}(z)\otimes1)\Delta_{d+e,f}(w)
   &=(1\otimes\Delta_{e,f}(w))\Delta_{d,e+f}(z+w),\\
 \Delta(z)\mathsf T
   &=\partial_z\Delta(z)+(1\otimes\mathsf T)\Delta(z),\\
 (\varepsilon\otimes1)\Delta(z)&=1,
 & (1\otimes\varepsilon)\Delta(z)&=e^{z\mathsf T}.
\end{align*}
Here $\varepsilon$ is restriction to the zero representation.
Rational factors in the first identity are compared with
$z,w,z+w$ invertible and expanded in the same iterated Laurent
region. The $(d,e)$ component has ordinary cohomological
degree $2\chi(d,e)$, including the degree of $z$.
\end{theorem}

\begin{proof}
The unit of direct-sum adjunction, natural pullback for vanishing
cycles, Thom--Sebastiani and K\"unneth define $\delta$.
Associativity of direct sum makes it coassociative. Line scaling
commutes with direct sum, so $\mathsf T$ is a coderivation.
Splitting either variable of $\Theta$ gives a direct sum of
the corresponding Hom complexes. Translating its first variable
by $u$ tensors it with $L_{-u}$, and translating the second
tensors it with $L_u$.

Consequently both iterated coproducts have factor
\[
 \Psi(\Theta_{12},-z)\Psi(\Theta_{13},-z-w)
                  \Psi(\Theta_{23},-w)
\]
on $e^{(z+w)\mathsf T_1+w\mathsf T_2}\delta^{(3)}$.
This proves coassociativity. Simultaneous translation fixes
$\Theta$, and
$\partial_z\Psi(\Theta,-z)=\mathsf T_1\Psi(\Theta,-z)$,
which prove the translation identity. A zero-dimensional
factor has $\Theta=0$, giving the two counit equations.
Finally $z^{\chi-r}c_r(\Theta)$ has degree $2\chi$.
This is the construction of
\cite[Theorem~2.3.1]{JindalKaubrysLatyntsev} in ordinary
cohomological grading.
\end{proof}

For the Hall regrading $A_d=C_d[-\chi(d,d)]$, the component
degree becomes $\chi(d,e)-\chi(e,d)$. Tensor symmetries must
be transported with this regrading. With the compatible sign
twist of the Hall product, the coproduct satisfies the vertex
bialgebra identity in the meromorphically braided coefficient
category \cite[Theorems~4.5.2 and~4.5.5]{JindalKaubrysLatyntsev}.
For a symmetric quiver that category consists of modules over
the tautological cup-product ring; for a nonsymmetric quiver
it uses the separate cohomology rings in each dimension vector.
These are concrete algebraic recipients, with their prescribed
braiding and spectral completion.

The spectral parameter has a geometric role. A localized Euler
factor $y-x+\lambda$ becomes $y-x+\lambda-z$ under first-factor
line scaling. Its leading coefficient is $-1$, so it is
invertible as a Laurent series even if the unshifted class is
a zero divisor. For a single arrow $1\to2$, dimensions
$d=(1,0)$ and $e=(0,1)$, and zero arrow weight,
\[
 \Delta_{d,e}(z)(1)=\frac1{y-x-z}
       =-\sum_{r\geq0}(y-x)^r z^{-r-1}.
\]
The general localized-to-spectral comparison is
\cite[Theorem~3.3.2]{JindalKaubrysLatyntsev}.
It does not extend a localized coefficient across its
collision divisor. Nor does it identify this coproduct with
the cut coproduct of Theorem~\ref{hall:thm:spectral-cut}:
one uses line scaling on cohomology, the other uses separation
opens in spectral configuration spaces. A comparison with the
compact reduced theory must specify its own tensor, support,
coefficient maps and relation to these two constructions.

There is also a localized bialgebra on the dual compactly
supported critical cohomology of semistable quiver
representations. For a fixed slope and a stability condition
generic at that slope, with the locally closed Serre-subcategory restriction
of \cite[\S6.1]{DavisonMeinhardtCoHA}, its tensor product
inverts the prescribed cross-factor equivariant Euler classes.
The perverse associated graded has a Hopf structure without
localization in the specified symmetric category and is free commutative
there; after the stated sign twist it is free supercommutative
\cite[Theorem~C]{DavisonMeinhardtCoHA}. The PBW comparison
with the original Hall object is a comparison of graded
coefficient objects and need not be an algebra isomorphism.
Thus localization and passage to the associated graded each
change the recipient of the coproduct assertion.

\subsubsection*{A collision coproduct is separate data}

\begin{definition}[A geometric collision coproduct]
\label{def:collision-coproduct}
\label{def:compact-hall-coproduct}
For each decomposition $\gamma=\alpha+\beta$, specify a
collision span
\[
 \mathcal M_\gamma\xleftarrow{\ r\ }
 \mathcal C_{\alpha,\beta}\xrightarrow{\ s\ }
 \mathcal M_\alpha\times\mathcal M_\beta
\]
and a coefficient morphism that induces a map in the chosen
realization. For ordinary cohomology, one possible form is
\[
 v_{\alpha,\beta}:r^*K_\gamma
 \longrightarrow s^!(K_\alpha\boxtimes K_\beta)[d_{\alpha,\beta}],
\]
with the support of $r^*K_\gamma$ proper over the endpoint
stack. The resulting component is
\[
 \mathcal H_\gamma\longrightarrow
 (\mathcal H_\alpha\otimes\mathcal H_\beta)[d_{\alpha,\beta}].
\]
A coproduct on a direct sum requires finitely many nonzero
components on every element; otherwise the target and its
completion must be specified. The analogous construction on
quotients uses the relative endpoint product over $BE$.
\end{definition}

\begin{proposition}[The coproduct integration leg]
\label{prop:finite-hall-coproduct-pushforward-existence}
Under the coefficient and support hypotheses in this definition,
pullback along $r$, then $v$, then the proper trace along $s$
gives the stated coproduct component.
\end{proposition}

\begin{proof}
Apply the same ordinary-cohomology pull--push construction as
for multiplication, with the collision span and its specified
coefficient map. Proper support is used at the trace, and
K\"unneth identifies the endpoint cohomology with the chosen
tensor product.
\end{proof}

\begin{proposition}[Coassociativity and compatibility]
\label{prop:finite-hall-coproduct-coassociativity-two-step-flags}
\label{prop:finite-hall-bialgebra-compatibility}
Suppose the two successive collision diagrams are identified,
and their coefficient composites agree. Then the coproduct is
coassociative. Suppose in addition that the diagram comparing
collision after extension with extension after two collisions
carries the corresponding coefficient equality. In an ordinary
super realization it gives
\[
 \Delta m=(m\otimes m)(1\otimes\tau\otimes1)
                     (\Delta\otimes\Delta).
\]
Unit and counit compatibility are additional equations.
\end{proposition}

\begin{proof}
The common triple-collision diagram gives the first equality by
base change and composition of the coefficient maps. In the
second diagram, the middle two factors must be interchanged
before the two products are formed. This is $\tau$, so its
Koszul sign is part of the displayed equality. These arguments
compare the given maps; they do not construct a missing collision
coefficient morphism from the extension product.
\end{proof}

The ordinary associative bar coalgebra has another coproduct:
cutting a word into two consecutive subwords. Its tensor,
support and differential are those of the bar construction.
Identifying it with a geometric collision coalgebra requires an
operation-preserving comparison. Neither a shared notation
$\Delta$ nor equality of Euler characteristics supplies that
comparison; see \S\ref{subsec:koszul-source}.

\subsubsection*{The complete bialgebra equations}

Assume now that $H$ is a complex of super vector spaces with
even chain maps $m,\eta,\Delta,\varepsilon$, with the ordinary
tensor product. All subsequent assertions in this subsection
have this carrier. If a geometrically constructed object is only
an algebra in a shifted or relative category, the following
criterion applies after its specified comparison to this carrier.

\begin{definition}[Bialgebra defects]
\label{def:finite-hall-bialgebra-defect-tuple}
The bialgebra defects are the differences between the two sides
of the following equations:
\begin{align*}
 m(m\otimes1)&=m(1\otimes m),&
 m(\eta\otimes1)&=1=m(1\otimes\eta),\\
 (\Delta\otimes1)\Delta&=(1\otimes\Delta)\Delta,&
 (\varepsilon\otimes1)\Delta&=1=(1\otimes\varepsilon)\Delta,\\
 \Delta m&=(m\otimes m)(1\otimes\tau\otimes1)(\Delta\otimes\Delta),&
 \Delta\eta&=\eta\otimes\eta,\\
 \varepsilon m&=\varepsilon\otimes\varepsilon,&
 \varepsilon\eta&=1_k.
\end{align*}
Call their collection $\mathfrak O_{\mathrm{Hbialg},R}$ on a
bounded realization. Vanishing means that every displayed map
identity holds on its defined domain. In finite bases these are
matrix and structure-constant equations. A differential is
required to commute with all four structure maps.
\end{definition}

The last three equations cannot simply be omitted. For example,
let $H=ke_1\oplus ke_2$ with
$e_ie_j=\delta_{ij}e_i$, unit $e_1+e_2$,
$\Delta e_i=e_i\otimes e_i$, and $\varepsilon(e_i)=1$.
Multiplication is associative and unital, comultiplication is
coassociative and counital, and $\Delta$ is multiplicative.
Nevertheless
\[
 \Delta(1)=e_1\otimes e_1+e_2\otimes e_2\ne1\otimes1,
 \quad \varepsilon(e_1e_2)=0\ne1,
 \quad \varepsilon(1)=2.
\]
This is the algebra of the groupoid with two identity objects,
not a bialgebra with the given structure maps.

\begin{definition}[Compact geometric bialgebra data]
\label{def:compact-geometric-hall-bialgebra-datum}
A compact geometric bialgebra datum consists of a coefficient
object, the extension and collision coefficient maps, zero-object
unit and counit, and the comparison to the stated tensor category.
The geometry must supply the proper-support or exceptional
integration maps it uses. Coherent flag, collision and mixed
diagrams must induce all equations of
Definition~\ref{def:finite-hall-bialgebra-defect-tuple}.
If finite matrices are used, their bases and finite-dimensional
models are part of the datum.
\end{definition}

\begin{proposition}[Geometric data and algebraic identities]
\label{prop:geometric-hall-data-give-bialgebra-defects}
Such a datum gives a bialgebra on its specified domain. Its
primitive Lie algebra is obtained from the reduced coproduct
below. The conclusion is conditional on the geometric and
coefficient comparisons, not just on formal names for the four
structure maps.
\end{proposition}

\begin{proof}
The flag and mixed correspondence comparisons give the displayed
identities by the preceding propositions. The zero-object and
counit diagrams give the unit equations, counit equations and
the last three scalar compatibilities. These are exactly the
bialgebra axioms.
\end{proof}


\subsubsection*{Primitives, parity and an antipode}

\begin{definition}[The primitive subspace]
\label{def:finite-hall-primitive-subspace}
Let $H$ be a counital coalgebra with coaugmentation $\eta$,
$\varepsilon\eta=1$ and $\Delta\eta=\eta\otimes\eta$.
Put $I=\ker\varepsilon$ and
$\pi=1-\eta\varepsilon:H\to I$. Its reduced coproduct and
primitive space are
\[
 \overline\Delta=(\pi\otimes\pi)\Delta|_I:I\longrightarrow I\otimes I,
 \qquad P(H)=\ker\overline\Delta.
\]
For $x\in I$, this is
$\overline\Delta x=\Delta x-1\otimes x-x\otimes1$.
A linear projection onto $P(H)$, if desired, is an extra choice
unless another theorem constructs one naturally.
\end{definition}

\begin{definition}[Primitive Hall bracket]
\label{def:primitive-hall-bracket}
For an ordinary super bialgebra $H$, the primitive bracket is
\[
 [x,y]=m(x,y)-(-1)^{|x||y|}m(y,x),\qquad x,y\in P(H).
\]
It is the restriction of the supercommutator in $H$.
\end{definition}

\begin{proposition}[Primitive closure and Jacobi]
\label{prop:hall-bialgebra-defects-primitive-closure}
\label{prop:finite-hall-graded-jacobi}
\label{lem:primitive-bracket-lifting}
The bracket takes values in $P(H)$ and makes it a Lie
superalgebra. It satisfies
\[
 [x,[y,z]]=[[x,y],z]+(-1)^{|x||y|}[y,[x,z]].
\]
A compatible differential restricts to its dg Lie differential.
\end{proposition}

\begin{proof}
Multiplicativity of the coproduct gives
\[
 \Delta(xy)=xy\otimes1+x\otimes y
       +(-1)^{|x||y|}y\otimes x+1\otimes xy
\]
for primitives. The mixed terms cancel in the supercommutator,
so the bracket is primitive. Expanding the three commutators
using associative multiplication gives the displayed Jacobi
identity. Compatibility of the differential with the
coaugmentation, counit, multiplication and coproduct proves
that it restricts and is a derivation.
\end{proof}

\begin{proposition}[What multiplication alone determines]
\label{prop:hall-product-only-no-primitives}
An associative superalgebra has a Lie superalgebra under its
supercommutator, but its primitive subspace is not determined
without a coproduct. On $k[t]$, the two bialgebra structures
\[
 \Delta_a t=t\otimes1+1\otimes t,\quad\varepsilon_a(t)=0,
 \qquad
 \Delta_m t=t\otimes t,\quad\varepsilon_m(t)=1
\]
have the same multiplication and different primitive spaces.
\end{proposition}

\begin{proof}
For the additive structure, $P(k[t])=kt$: the mixed terms in
$\Delta_a(t^n)$ exclude every highest degree $n\geq2$.
For the multiplicative structure, every basis element $t^n$
is group-like. Comparing diagonal coefficients as in
Proposition~\ref{prop:protected-integration-primitive-compatibility}
shows that its primitive space is zero.
\end{proof}

The ambiguity also occurs among connected positively graded
Hopf structures. Give $k[x,y]$ degrees $1,2$. One coproduct
makes both $x$ and $y$ primitive; another has
\[
 \Delta x=x\otimes1+1\otimes x,\qquad
 \Delta y=y\otimes1+1\otimes y+x\otimes x.
\]
For the second, $y-x^2/2$ is primitive. This change of polynomial
coordinates identifies it with the additive Hopf algebra on
two generators and proves coassociativity. The two primitive
planes inside the fixed polynomial algebra are
$\operatorname{span}\{x,y\}$ and
$\operatorname{span}\{x,y-x^2/2\}$.

\begin{proposition}[Charge and parity]
\label{prop:finite-hall-normal-ordered-gram-degree}
\label{prop:finite-hall-bracket-parity}
Suppose the product preserves an additive charge grading
$\widehat\Gamma$ and cohomological parity. If the coproduct and
coaugmentation preserve the grading, its primitive space is
graded, and
\[
 [P_{\widehat\alpha,\bar p},P_{\widehat\beta,\bar q}]
 \subset P_{\widehat\alpha+\widehat\beta,\bar p+\bar q}.
\]
Regrading by the additive map $\overline\Pi_X$ preserves this
law. A raw quadratic charge expression can be used only through
the normal-ordered extension on which it becomes additive.
\end{proposition}

\begin{proof}
Both terms in the bracket have the same summed charge and
parity. Primitive closure was proved above. The composite with
an additive grading map sends their common charge to
$\overline\Pi_X(\widehat\alpha)+\overline\Pi_X(\widehat\beta)$.
The construction and group law of the normal-ordered charge
extension are given in Chapter~\ref{ch:k3xe-setup}.
\end{proof}

The primitive projection is not part of this proof. In the
connected cocommutative characteristic-zero domain, the
convolution logarithm of the identity gives a natural primitive
projection; Proposition~\ref{prop:bkm-connected-hopf-recognition}
proves that result with its local-nilpotence hypotheses. It must
not be applied to a general noncocommutative Hall bialgebra
without those hypotheses.

\begin{definition}[Antipode]
\label{def:finite-hall-antipode-datum}
An antipode on a bialgebra is a chain map $S:H\to H$ satisfying
\[
 m(S\otimes1)\Delta=\eta\varepsilon
                 =m(1\otimes S)\Delta.
\]
A bialgebra admitting this convolution inverse of its identity
is a Hopf algebra. On a connected positively graded bialgebra,
the antipode is defined recursively by
\[
 S(1)=1,\qquad S(x)=-x-\sum S(x')x'',
 \qquad\overline\Delta x=\sum x'\otimes x''.
\]
Each reduced factor has lower positive degree.
\end{definition}

\begin{proposition}[Existence from conilpotence]
\label{prop:finite-hall-nonhopf-boundary}
If the augmentation projection $\pi=1-\eta\varepsilon$ is
locally nilpotent under convolution, then
\[
 S=\sum_{n\geq0}(-1)^n\pi^{\star n}
\]
is an antipode. In particular this holds for a connected
positively graded bialgebra. It preserves additive charge.
A bialgebra need not admit an antipode: the multiplicative
structure on $k[t]$ above is an example.
\end{proposition}

\begin{proof}
The convolution unit is $\eta\varepsilon$, and
$1_H=\eta\varepsilon+\pi$. Local nilpotence makes the
geometric series finite on each element. Multiplying the series
by $\eta\varepsilon+\pi$ on either side gives
$\eta\varepsilon$, proving both identities. In a connected
positive grading, iterated reduced coproducts strictly split
the total degree, so they eventually vanish. Every term has
the original total charge. For multiplicative $k[t]$, an
antipode would satisfy $S(t)t=1$, which has no polynomial
solution.
\end{proof}

An antipode therefore preserves an additive charge; it does not
create the opposite charge sector. A two-sided comparison or a
Drinfeld double has to construct that opposite sector separately.

\subsubsection*{Finite positive windows}

Let $H=\bigoplus_{n\geq0}H_n$ have additive positive grading.
The direct sum $H_{\leq R}$ carries every product with total
degree at most $R$. If $\Delta$ preserves total degree, this
window is a subcoalgebra. Declaring all higher products to be
zero instead gives the algebra quotient
$H/H_{>R}$, but that ideal need not be a coideal.

For the additive Hopf algebra $k[x]$, take $N\geq2$. Then
\[
 \Delta(x^N)=x^N\otimes1+1\otimes x^N
       +\sum_{i=1}^{N-1}\binom Ni x^i\otimes x^{N-i}.
\]
The coefficient of $x\otimes x^{N-1}$ is $N\ne0$.
Thus $(x^N)$ is not a coideal, and the algebra
$k[x]/(x^N)$ does not inherit the additive coproduct in the
ordinary tensor category. It still has the partial operations
obtained by restricting total degree. This elementary example
prevents a bounded Hall charge window from being called a finite
Hopf algebra merely by setting operations outside it to zero.

\subsubsection*{A primitive Lie form and a Hopf pairing}

A compact effective Hall support generally contains no opposite
charge. To pair $\gamma$ with $-\gamma$, choose an enlarged
carrier or independently defined opposite object and a
comparison realizing that pairing. An antipode on a positively
graded Hall algebra does not supply it.

\begin{definition}[A positive--negative primitive form]
\label{def:finite-hall-positive-negative-pairing}
Let $L=\bigoplus_{\gamma,\bar p}L_{\gamma,\bar p}$ be a Lie
superalgebra containing the required positive and negative
sectors. An even degree-zero primitive form is a bilinear map
$B:L\otimes L\to k$ pairing only opposite charges and equal
parities. It is supersymmetric when
$B(x,y)=(-1)^{|x||y|}B(y,x)$, and invariant when
\[
 B([x,y],z)=B(x,[y,z]).
\]
Its invariance defect is
$\mathfrak o_F(x,y,z)=B([x,y],z)+(-1)^{|x||y|}B(y,[x,z])$.
For an even supersymmetric form, its vanishing for all
homogeneous triples is equivalent to the displayed invariance
identity, by super skew-symmetry of the bracket.
One possible geometric construction starts with an even
supertrace $\operatorname{tr}_0$ on a two-sided graded
associative realization and sets
$B(x,y)=\operatorname{tr}_0(xy)$ on its chosen Lie subspace.
The existence and descent of that trace are part of the data.
\end{definition}

\begin{proposition}[Trace forms]
\label{prop:finite-hall-pairing-homogeneity}
\label{prop:finite-hall-pairing-supersymmetry}
\label{prop:finite-hall-pairing-invariance}
If the trace is even, vanishes in nonzero charge, and satisfies
$\operatorname{tr}_0(xy)=(-1)^{|x||y|}\operatorname{tr}_0(yx)$,
then its product form has the stated homogeneity,
supersymmetry and Lie invariance.
\end{proposition}

\begin{proof}
Charge preservation of multiplication and the degree-zero trace
give homogeneity. Its parity gives zero between unequal
parities, and its cyclic identity gives supersymmetry. Finally,
\begin{align*}
 B([x,y],z)
 &=\operatorname{tr}_0(xyz)
       -(-1)^{|x||y|}\operatorname{tr}_0(yxz)\\
 &=\operatorname{tr}_0(xyz)
       -(-1)^{|y||z|}\operatorname{tr}_0(xzy)
 =B(x,[y,z]).\qedhere
\end{align*}
\end{proof}

\begin{definition}[Pairing identities on their carriers]
\label{def:finite-hopf-pairing-identity-tuple}
For a primitive Lie form, its identities are evenness,
homogeneity, supersymmetry and Lie invariance. Its radical and
quotient pairing are then formed on $L$.
A Hopf pairing instead is a bilinear map
$\beta:H\otimes K\to k$ between full Hopf superalgebras such
that
\[
 \beta(xy,z)=\beta(x\otimes y,\Delta_Kz),\qquad
 \beta(x,yz)=\beta(\Delta_Hx,y\otimes z),
\]
with $\beta(1,z)=\varepsilon_K(z)$ and
$\beta(x,1)=\varepsilon_H(x)$.
The tensor pairing includes the super sign
\[
 \beta(x\otimes y,u\otimes v)
 =(-1)^{|y||u|}\beta(x,u)\beta(y,v).
\]
Antipode compatibility is imposed with the convention of the
chosen pairing. These equations concern the full algebra
products, not a product asserted to close on primitives.
The notation $\mathfrak O_{\mathrm{Hpair},R}=0$ records the
identities for the pairing on its specified carrier. A Lie-form
comparison uses the first set; a double uses the full Hopf
pairing equations. It does not require the same form to serve
both roles.
\end{definition}

\begin{proposition}[Hopf adjointness on primitives]
\label{prop:finite-hall-pairing-coproduct-adjointness}
If $x,y$ are primitives in $H$ and $z$ is primitive in $K$,
then $\beta(xy,z)=0$. Hence the restriction of a Hopf pairing
to primitives annihilates $[P(H),P(H)]$ in its first argument.
If that restricted pairing is nondegenerate there, then
$P(H)$ is abelian.
\end{proposition}

\begin{proof}
Insert $\Delta_Kz=z\otimes1+1\otimes z$ in the first
adjointness equation. Each term contains either
$\varepsilon_H(x)$ or $\varepsilon_H(y)$, both zero.
Subtract the two products to obtain the commutator assertion.
Nondegeneracy forces every such commutator to vanish.
\end{proof}

Thus a nondegenerate invariant form on a nonabelian primitive
Lie superalgebra cannot simultaneously be the restriction of
this Hopf pairing. The primitive-recognition problem uses the
Lie form. A Drinfeld double uses an appropriate Hopf pairing
on its full positive and negative halves. Their comparison
must retain this distinction.


\subsubsection*{Radicals and quotient forms}

\begin{proposition}[Finite quotient pairing]
\label{prop:finite-hall-quotient-pairing-nondegeneracy}
Let $G:V\otimes W\to k$ be a bilinear pairing of
finite-dimensional spaces. Quotienting $V$ by its left radical
and $W$ by its right radical gives a perfect pairing. If
$Q_V,Q_W$ are matrices of chosen quotient lifts, its matrix is
\[
 \overline G=Q_V^tGQ_W.
\]
It has rank equal to both quotient dimensions. To apply this
criterion to geometric primitives, the actual matrix $G$ and
its two radicals must first be identified.
\end{proposition}

\begin{proof}
Both radicals vanish against the opposite space, so the pairing
descends. A nonzero class in either quotient has nonzero pairing
with some vector in the opposite space by the definition of its
radical. Thus both adjoint maps are injective. Their dimensions
are both $\operatorname{rank}G$, so they are isomorphisms.
Evaluating on the chosen lifts gives the displayed matrix.
\end{proof}

\begin{definition}[Geometric pairing data]
\label{def:compact-geometric-hopf-pairing-datum}
For the compact primitive comparison, geometric pairing data
specify the positive and negative carrier, the pairing
correspondence or trace, its orientation and degree shifts, and
the cyclic three-input comparison proving Lie invariance.
The form and its radical are then computed on the resulting
primitive Lie algebra.
If a Hopf pairing for a double is also required, its two full
Hopf algebras, product--coproduct adjointness, unit--counit
normalization and antipode compatibility are additional data.
They are not imposed as substitute axioms on the primitive
Lie form.
\end{definition}

\begin{proposition}[Lie radicals and enveloping quotients]
\label{lem:hopf-radical-ideal-coideal}
Let $L$ have an even supersymmetric invariant form $B$, and let
$R=\operatorname{rad}B$. Then $R$ is a Lie ideal, and $L/R$
has the descended nondegenerate invariant form. The ideal
generated by $R$ in $U(L)$ is a Hopf ideal, and
\[
 U(L)/(U(L)R)\simeq U(L/R).
\]
This construction does not require $B$ to be a Hopf pairing.
\end{proposition}

\begin{proof}
For homogeneous $x\in L$, $r\in R$ and $z\in L$, invariance
and supersymmetry give
\[
 B([x,r],z)=-(-1)^{|x||r|}B(r,[x,z])=0.
\]
Thus $R$ is an ideal. The form descends and is nondegenerate
by the definition of the radical. In $U(L)$ each $r$ is
primitive, so $\Delta r=r\otimes1+1\otimes r$,
$\varepsilon(r)=0$ and $S(r)=-r$. The generated ideal is
therefore also a coideal and is stable under the antipode.
The universal property of the enveloping algebra identifies
its quotient with $U(L/R)$.
\end{proof}

\begin{proposition}[Radicals of a full Hopf pairing]
\label{hall:prop:full-hopf-radical}
Let $\beta:H\otimes K\to k$ be a Hopf pairing. Suppose its
induced quotient tensor pairings separate points, as they do
in finite-dimensional quotients. Then its left and right
radicals are Hopf ideals in their respective full algebras.
\end{proposition}

\begin{proof}
If $r$ is in the left radical, product adjointness gives
$\beta(hr,z)=\beta(h\otimes r,\Delta z)=0$, and similarly
for $rh$. Pairing $r$ against the unit gives
$\varepsilon_H(r)=0$. The other adjointness equation says
$\beta(\Delta r,y\otimes z)=\beta(r,yz)=0$.
Separation by the quotient tensor pairing forces the image
of $\Delta r$ in the tensor of the two quotient spaces to
vanish. Thus $\Delta r\in R\otimes H+H\otimes R$.
Antipode compatibility gives stability under $S_H$.
The argument for the right radical is identical with the
roles exchanged.
\end{proof}

\begin{proposition}[What the pairing matrix proves]
\label{prop:pairing-matrix-not-hopf-pairing-proof}
\label{prop:geometric-hopf-pairing-data-give-defects}
A finite matrix and its kernel determine a bilinear radical
and its perfect quotient pairing. Lie-ideal stability requires
invariance of the primitive form. Hopf-ideal stability in a
full algebra requires the full pairing hypotheses just stated,
or the enveloping-algebra construction from a Lie ideal.
These are distinct ways to construct a quotient.
\end{proposition}

\begin{proof}
The linear-algebra assertion is
Proposition~\ref{prop:finite-hall-quotient-pairing-nondegeneracy}.
The other assertions follow from the two preceding propositions,
whose proofs use the respective invariance and adjointness
identities. A kernel computation alone contains neither identity.
\end{proof}

The resulting quotient Lie algebra is still a geometric source
for the recognition problem. A comparison with the BKM target
must prove its generators, brackets, form and absence of extra
relations as stated in \S\ref{subsec:primitive-recognition}.
A matrix rank or a signed dimension does not give that map.

\subsubsection*{Transition maps}

Let $R'\leq R$. Write
$T_{R'R}:H_R\to H_{R'}$ for a proposed inverse transition.
It must first be defined by a morphism of coefficient diagrams
or by another specified chain construction. The direction of
a dual transition is reversed and is not silently identified
with this one.

\begin{definition}[Multiplicative transitions]
\label{def:transition-hall-product-preservation}
A multiplicative transition commutes with the differential,
unit and multiplication:
\[
 T_{R'R}m_R=m_{R'}(T_{R'R}\otimes T_{R'R}),\qquad
 T_{R'R}\eta_R=\eta_{R'}.
\]
For partial charge windows the first equation is imposed only
where both sides are defined. The transition comparisons must
also respect shifts, orientation and the relative tensor when
those are present.
\end{definition}

\begin{proposition}[Multiplication under transition]
\label{prop:transition-hall-product-preservation-criterion}
A coefficient-diagram transition preserving the extension
maps and their trace comparisons satisfies these equations.
\end{proposition}

\begin{proof}
Move the transition through each leg of the defining composite
for the product, using its supplied pullback, coefficient and
trace naturality squares. Their composition is the displayed
multiplicative square. The zero-object square gives the unit.
\end{proof}

\begin{definition}[Coalgebra transitions]
\label{def:transition-hall-coproduct-preservation}
A coalgebra transition satisfies
\[
 \Delta_{R'}T_{R'R}=(T_{R'R}\otimes T_{R'R})\Delta_R,
 \qquad\varepsilon_{R'}T_{R'R}=\varepsilon_R,
 \qquad T_{R'R}\eta_R=\eta_{R'}.
\]
The maps for successive stages compose and preserve the
specified degree and parity.
\end{definition}

\begin{proposition}[Coproduct under transition]
\label{prop:transition-hall-coproduct-preservation-criterion}
Transitions of collision coefficient diagrams preserving all
their comparison maps satisfy these equations.
\end{proposition}

\begin{proof}
Compose the naturality squares for the collision pullback,
coefficient morphism and trace. They give the coproduct square.
The empty-output and zero-object squares give the counit and
coaugmentation identities.
\end{proof}

\begin{definition}[Primitive transitions]
\label{def:transition-primitive-subspace-preservation}
A primitive transition is the restriction
$T^{\mathrm{Prim}}_{R'R}:P(H_R)\to P(H_{R'})$ of a
coaugmentation- and coproduct-preserving transition. If an
independently chosen primitive projection is used, its
commutation with $T_{R'R}$ is an additional equation.
\end{definition}

\begin{proposition}[Preservation of the primitive kernel]
\label{prop:transition-primitive-subspace-preservation-criterion}
Every coalgebra transition above preserves primitives, and
multiplicative transitions also preserve their Lie bracket.
\end{proposition}

\begin{proof}
The reduced-coproduct square commutes, so its source kernel
maps into its target kernel. Multiplicativity then preserves
the two terms of the supercommutator.
\end{proof}

\begin{definition}[Radical transitions]
\label{def:transition-radical-preservation}
Let $L_R$ and $L_{R'}$ have the actual primitive forms $B_R$
and $B_{R'}$, and let $T:L_R\to L_{R'}$ preserve the Lie
bracket. A sufficient pairing condition for radical descent is:
for every opposite homogeneous $y'\in L_{R'}$, there exists
$a(y')\in L_R$ such that
\[
 B_{R'}(Tx,y')=B_R(x,a(y'))\quad\text{for all }x\in L_R.
\]
An adjoint lift satisfying $Ta(y')=y'$ is one possible stronger
choice. In finite coordinates the condition is
$T_\gamma^tG_{R',\gamma}=G_{R,\gamma}A_{-\gamma}$.
\end{definition}

\begin{proposition}[Radical descent]
\label{prop:transition-radical-preservation-criterion}
Under this condition $T(\operatorname{rad}B_R)
\subset\operatorname{rad}B_{R'}$. Hence $T$ descends to the
quotient Lie algebras. Compatible radical-preserving
transitions give a quotient inverse system.
\end{proposition}

\begin{proof}
For $x$ in the source radical, the right side of the pairing
identity vanishes for every $y'$. Thus $Tx$ lies in the target
radical. The quotient map is well defined, and it is a Lie
map because $T$ is one. Its composition laws are inherited
from the original transitions.
\end{proof}

\subsubsection*{Coalgebras and their completions}

\begin{definition}[A geometric conilpotent coalgebra]
\label{def:conilpotent-source-coalgebra}
A geometric coalgebra $C_{X,R}$ is a coefficient object with a
specified tensor, coaugmentation, counit and geometric collision
coproduct. It is conilpotent if each element is killed by a
sufficiently long iterated reduced coproduct. A sufficient
filtration is
$F_0=k1\subset F_1\subset\cdots$, exhaustive, with
\[
 \overline\Delta F_n\subset
 \sum_{i,j>0,\ i+j\leq n}F_i\otimes F_j.
\]
If this is a chiral coalgebra on $E$ or on its Ran space, the
coefficient category, chiral tensor, supports and filtration
are part of the definition; ordinary vector-space data alone
do not define it. Transitions preserve all these structures.
\end{definition}

\begin{proposition}[Conilpotence and a completed limit]
\label{prop:conilpotent-source-coalgebra-firewall}
The displayed positive filtration implies conilpotence of
each algebraic element. An inverse limit of such coalgebras
need not be conilpotent as an ordinary coalgebra. It carries
the corresponding completed coproduct when that coproduct
is continuous for the specified inverse-limit topology.
\end{proposition}

\begin{proof}
Every positive factor in an iterated reduced coproduct consumes
at least one unit of filtration degree, so more than $n$
factors kill $F_n$. For the limit, consider the tensor
coalgebra truncated by word length and take its inverse limit.
An element with nonzero components in all lengths has
nonzero reduced coproducts of arbitrarily high length.
The compatible finite coproducts land instead in the completed
tensor product, where the induced limit map is continuous.
\end{proof}

\begin{proposition}[Exactness of finite towers]
\label{prop:sectorial-hall-truncation}
For a countable inverse tower of actual finite-dimensional
vector spaces in each fixed degree, the Mittag--Leffler
condition holds and $R^1\varprojlim$ vanishes. Thus a short
exact sequence of such towers remains exact under inverse
limit. In particular, once compatible finite primitive and
radical towers have been constructed,
\[
 \varprojlim_R(L_R/\operatorname{rad}B_R)
 \simeq(\varprojlim_RL_R)/(\varprojlim_R\operatorname{rad}B_R).
\]
The kernel of an actual commuting map of inverse systems
commutes with inverse limit without a Mittag--Leffler
hypothesis. Operations on the limits still require their
stated continuous tensor products.
\end{proposition}

\begin{proof}
The images in a fixed finite-dimensional term form a descending
chain, which stabilizes. This is Mittag--Leffler and gives
$R^1\varprojlim=0$ for the countable tower, as proved in
Proposition~\ref{prop:bkm-finite-tower-ml}. Apply the derived
inverse-limit exact sequence to the radical short exact
sequence. Kernels commute with inverse limits because both
are limits. The continuous operation statement concerns the
definition of their target topology, not exactness alone.
\end{proof}

The finiteness premise is substantial for moduli coefficients.
For unbounded or infinite-dimensional systems one must prove
an appropriate Mittag--Leffler property or retain the derived
inverse-limit correction. Fibrewise K3 boundedness alone also
does not construct a finite-type stability moduli space on
$S\times E$. On the product hearts and stability conditions of
\cite{AbramovichPolishchukTStructures,LiuProductStability},
the required input is boundedness in every selected class,
together with a support bound making the permitted charge
window finite. Threefold boundedness results such as
\cite{PiyaratneTodaModuli3folds} apply only after their stability
hypotheses have been verified. Appendix~\ref{app:stability-support}
states these product-stability and support requirements.

\subsubsection*{Integral forms and the compact comparison}

\begin{definition}[An integral-form quantum-group comparison]
\label{def:integral-form-stable-envelope-comparison-datum}
An integral-form comparison starts with a specified graded
Hopf algebra and a lattice stable under its operations. It
includes a separated complete topology, an appropriate full
Hopf pairing for the double, its cross-relations and Cartan
part, and a comparison preserving these structures with the
proposed quantum group. A stable-envelope realization also
specifies the geometric correspondence, chamber,
polarization, support and normalization conditions, and its
chamber composition maps.
For the Igusa problem, one proposed target is denoted
$Y_\hbar(\widehat{\mathfrak g}_{\Delta_5})$. Here the chosen
central extension, its quantum-group presentation and its
coproduct have to be defined independently. The desired
comparison is then
\[
 D_\hbar(H^+)^{\mathrm{int}}_{\mathrm{compl}}
 \longrightarrow Y_\hbar(\widehat{\mathfrak g}_{\Delta_5}).
\]
An isomorphism would require an inverse preserving these
operations. The notation for the target does not itself
construct a Yangian for the Borcherds Lie superalgebra.
The compact Hall--Drinfeld proposal of \cite{CYQGVolIII} is
subject to the same coefficient, pairing and completion requirements.
\end{definition}

\begin{proposition}[Scope of stable-envelope transport]
\label{prop:integral-form-stable-envelope-transport-firewall}
The stable-envelope theorem and the Schiffmann--Vasserot
comparisons construct quantum-group objects on their specified
symplectic and instanton moduli. Applying them to the compact
reduced Hall construction requires the preceding comparison
data. The formal tripled-quiver chart alone does not extend
those objects or their operators to compact moduli.
\end{proposition}

\begin{proof}
Appendix~\ref{app:stable-envelope} states the geometric
hypotheses and the actual algebras of those constructions.
A formal chart identifies a completed local deformation
problem. It does not furnish a proper symplectic resolution,
a torus action and stable envelopes on the whole compact
moduli space, or the required coefficient and operator
comparison. When those data are supplied, the induced map
can be tested on the products, coproducts and double
cross-relations stated in the definition.
\end{proof}

The normal-ordered charge group is independent of a stability
chamber, but the selected sheaf moduli and their Hall operations
can change with the chamber. A wall-crossing comparison must
preserve the actual coefficient operations and match the
primitive curve-counting chamber before its scalar series is
identified with the normalized Igusa denominator. This is
additional to a charge-lattice identification.

The original compact problem therefore has several concrete
next maps. First construct the reduced oriented coefficient
morphism on extensions and its equality on three-step flags.
Then construct the collision coproduct and the mixed
product--coproduct comparison. An operator realization of the
BKM algebra further requires the primitive Lie form and the
recognition map; a geometric Pfaffian or protected trace
requires its own family, state and modular-line comparison.
The local constructions above supply domains and equations
for these tasks without replacing them by the scalar product.
The collective notation $\mathfrak O_{\mathrm{Hall},R}=0$
means that these geometric data and their relevant bialgebra,
pairing and transition equations have actually been supplied.
Exactness of an inverse limit is imposed at the chosen
completion, not inferred from the name of a finite stage.


\subsection{Finite skew blocks and the protected Pfaffian}
\label{subsec:dirac-pfaffian}

A square root of a determinant line need not come with a section.
A skew operator supplies one: its Pfaffian.  This subsection first
constructs that finite algebraic object, then states the geometric
comparisons needed to regard it as the protected section on the
compact reduced moduli space.

\subsubsection*{The Pfaffian of a skew two-term complex}

Let \(B\) be a complex base, \(V\) a vector bundle of even rank
\(2n\), and \(D:V\to V^\vee\) a skew map.  Regard
\[
 K_D=[V\xrightarrow{D}V^\vee]
\]
as a complex in degrees \([-1,0]\).  In a frame its skew matrix has
a Pfaffian, and under a frame change \(A\) it satisfies
\[
 \operatorname{Pf}(A^tDA)=\det(A)\operatorname{Pf}(D).
\]
Thus these local functions define a section of \(\det V^\vee\).

\begin{definition}[Pfaffian line and section]
\label{def:pfaffian-line}
For this skew complex define
\[
 \mathscr L_{\mathrm{Pf}}(K_D)=\det V^\vee,
 \qquad \operatorname{pf}(D)\in
 H^0(B,\mathscr L_{\mathrm{Pf}}(K_D)).
\]
The determinant identity gives
\[
 \mathscr L_{\mathrm{Pf}}(K_D)^{\otimes2}
 \simeq\det K_D,
 \qquad \operatorname{pf}(D)^2=\det(D).
\]
A geometric Pfaffian datum specifies such local skew models and
compatible Pfaffian-line and section comparisons on their overlaps,
or another precise construction of the same line and section.
An orientation of a critical locus is a root pair; it does not by
itself specify these operators or their Pfaffian sections.
\end{definition}

The section is also \(\omega_D^n/n!\), where
\(\omega_D\in\Lambda^2V^\vee\) is the alternating form of \(D\).
In a hyperbolic basis its coefficient is the product of the skew-block
entries and its square is the determinant.  The identity holds on the
dense set of invertible skew matrices and hence, being polynomial, on
all skew matrices over \(\mathbb C\).  It therefore holds in families
by the same universal polynomial identity.

The degree convention matters.  For the unshifted CY3 self-Ext
complex \(K=R\operatorname{Hom}(F,F)\), Serre duality has the form
\(K\simeq K^\vee[-3]\).  If a specified reduced complex has only
Ext degrees one and two, shifting it by \([2]\) places those terms
in degrees \([-1,0]\), and changes the self-duality to a map to its
dual shifted by \([1]\).  A comparison with \(K_D\) must include
this shift, the skew symmetry and the determinant-square map.
One cannot impose the latter self-duality on the unshifted complex
without changing its meaning.  The shifted degree condition alone
also does not provide a skew realization or a Pfaffian section.

\subsubsection*{Hyperbolic factors and the super ratio}

Choose a positive charge monoid \(\Gamma^+\) and a height function
\(h\) which is positive on its nonzero elements and locally finite
on the active degrees.  Work in its completed monoid algebra, or
at a fixed height jet.  Then \(1-x^\gamma\) is a unit for
\(\gamma\ne0\).  Let \(P_\gamma=P_{\gamma,\bar0}\oplus
P_{\gamma,\bar1}\) have finite dimension \((p_\gamma\mid q_\gamma)\).

\begin{definition}[A finite skew block]
\label{def:first-order-dirac-block}
For one basis vector of either parity use the ordered hyperbolic
block
\[
 J_\gamma=
 \begin{pmatrix}0&1-x^\gamma\\-(1-x^\gamma)&0\end{pmatrix}.
\]
Take \(p_\gamma\) such blocks in the even part and
\(q_\gamma\) in the odd part.  Their super-Pfaffian means the ratio
of the two ordinary Pfaffians:
\[
 \operatorname{sPf}(J_{\gamma,P})
 =\frac{\operatorname{Pf}(J_\gamma^{\oplus p_\gamma})}
        {\operatorname{Pf}(J_\gamma^{\oplus q_\gamma})}
 =(1-x^\gamma)^{p_\gamma-q_\gamma}.
\]
The ordinary Pfaffian of all the blocks together has exponent
\(p_\gamma+q_\gamma\), not \(p_\gamma-q_\gamma\).
The ratio is defined on the formal unit domain just specified, or
as a meromorphic section where its denominator is allowed.
\end{definition}

The matrix is linear in the scalar \(1-x^\gamma\).  This does not
make it the symbol of the Hall bracket or an analytic Dirac
operator.  An abelian three-dimensional Lie algebra and the
three-dimensional Heisenberg algebra give identical such blocks
when assigned the same graded vector space, but have different
brackets.  In a Chevalley--Eilenberg model the bracket is a binary
operation in its differential; in an enveloping algebra it appears
in the Poisson bracket of the PBW associated graded.  Relating either
construction to a differential operator requires a separate map.

\subsubsection*{A finite geometric comparison}

\begin{definition}[Finite geometric Pfaffian data]
\label{def:finite-geometric-pfaffian-datum}
A finite geometric Pfaffian datum \((P_{\mathrm{fin}})\) consists
of a specified reduced quotient base \(\mathfrak Q_R\), its perfect
reduced complex, oriented critical coefficients and skew realization
with Pfaffian line \(\mathscr L_{\mathrm{Pf},R}\) and section
\(\operatorname{pf}_{X,R}\).  The skew realization is compatible
with the shift and square map above.

On a specified cusp chart, the datum further gives an oriented
block comparison with the finite family \(J_{\gamma,P}\), including
its determinant of frame change and every residual unit.  The scalar
normalization is fixed as
\(64x^\rho\), where \(x^\rho=q^{1/2}r^{1/2}s^{1/2}\).  Thus the
block comparison identifies the section itself, rather than only
its rank or squared determinant.  Each \(P_\gamma\) is compared as
a graded space with the target root space \(V_\gamma\).  In
particular the source is zero where the target graded space is zero;
vanishing of its signed dimension alone is not this condition.

At a simple type-II wall there is a compatible local factorization
\(\operatorname{pf}_{X,R}=\upsilon u\,
\operatorname{pf}^{\mathrm{tan}}_{X,R}\), with simple normal zero,
invertible \(\upsilon\) and a stipulated reflection lift.  For the
sign conclusion the reference comparison preserves the tangential
factor and unit, the base map sends \(u\) to \(-u\), and the actual
root lift carries the section to its pullback.

Finally the datum specifies its parameter map, Pfaffian-to-automorphic
line comparison and transition isomorphisms.  Transitions preserve
the square maps, the sections, the root comparisons and the critical
coefficient data.  An effective geometric limit is required when
these are claimed to define a line and section on that limit.
\end{definition}

This definition retains the distinction between a positive height
jet and an actual finite moduli space.  On the former, adding roots
of higher height does not change a lower jet.  It need not leave
an ordinary section on the latter unchanged.  A transition of the
moduli problem must therefore be supplied separately.  Likewise a
root with signed multiplicity zero can still have equal nonzero
even and odd parts; the full graded comparison detects them.

\begin{proposition}[The finite Pfaffian product]
\label{prop:pfaffian-product-finite-stage}
Suppose \((P_{\mathrm{fin}})\) supplies the stated block and
normalization comparison on a finite active set \(A_R\).  Then
\[
 \operatorname{pf}_{X,R}|_{\mathfrak c_\infty}
 =64q^{1/2}r^{1/2}s^{1/2}
 \prod_{\gamma\in A_R}(1-x^\gamma)^{p_\gamma-q_\gamma}.
\]
If \(P_\gamma\simeq V_\gamma\) with both parities preserved,
its exponents equal the target Borcherds exponents
\(a_\Delta(\gamma)=\operatorname{sdim}V_\gamma\).
\end{proposition}
\begin{proof}
The Pfaffian is multiplicative on the ordered orthogonal sum of
skew blocks.  Divide the odd Pfaffian into the even Pfaffian and
use the specified frame and normalization to obtain the formula.
The graded root comparison gives the exponent equality.  Merely
knowing the ranks and the leading coefficient would not prove this
formula: multiplying a section by a nonconstant unit with constant
term one preserves those data.  The actual block comparison excludes
that additional unit.
\end{proof}

Write \(J(z)=\left(\begin{smallmatrix}0&z\\-z&0\end{smallmatrix}\right)\).
Then \(J(1-x)\) and \(J(1-x)(1+x^2)\) have the same rank,
leading coefficient and simple divisor at \(x=1\) wherever
\(1+x^2\) is invertible.  Their Pfaffian sections differ.  The
calculation exhibits the need to identify every unit in the block
comparison.

\begin{definition}[The automorphic line comparison]
\label{def:pfaffian-to-automorphic-line-comparison}
Let \(\kappa:B\to\mathcal A_2\) be the specified parameter map,
on the chosen stack or cover carrying the weight-five character
line.  A Pfaffian-to-automorphic comparison is an isomorphism
\[
 \iota_{\mathrm{aut}}:\mathscr L_{\mathrm{Pf},X}
 \xrightarrow{\sim}
 \kappa^*(\mathcal L^5\otimes\nu_{\Delta_5}),
\]
compatible with the square map, the cusp frame and the stated
boundary extension.  A section identity means
\(\iota_{\mathrm{aut}}(\operatorname{pf}_X)=\kappa^*\Delta_5\).
It includes the map \(\kappa\); the two lines do not initially live
on the same base.
\end{definition}

\begin{proposition}[What scalar squaring forgets]
\label{prop:scalar-pfaffian-data-not-section}
Local scalar products, their divisors, their leading coefficients
and their squares do not determine the descent of a square-root
section.  In particular they do not construct
\(\iota_{\mathrm{aut}}\).
\end{proposition}
\begin{proof}
Let \(T_\chi\) be a nontrivial sign local system on a punctured
quotient chart.  Tensoring a local square-root section by a flat
frame of \(T_\chi\) leaves its local expression and its square
unchanged.  Around a loop with \(\chi=-1\), that frame changes
sign.  The two local sections therefore have different descent
data, even though their squares agree.  This example concerns the
flat or equivariant line with its specified descent; it makes no
claim that a given scalar expression has two distinct global
holomorphic descents.  Those descents are precisely the missing
information.
\end{proof}

\subsubsection*{Type-II crossings and local normalization}

\begin{definition}[A rank-one Pfaffian crossing]
\label{def:rank-one-type-ii-crossing}
A normal Pfaffian crossing is a rank-two skew block
\[
 D_u=\begin{pmatrix}0&u\\-u&0\end{pmatrix},
 \qquad\operatorname{Pf}(D_u)=u,
\]
with a specified reference frame comparison and base reflection
\(u\mapsto-u\).  Its tangential Pfaffian is an invariant unit in
that reference comparison.  The orientation lift is required to carry
the section to its pullback.  The phrase ``rank one'' counts one hyperbolic
Pfaffian factor; the skew matrix itself has size two.
\end{definition}

\begin{proposition}[The normal sign]
\label{prop:rank-one-crossing-o2-sign}
\label{prop:generic-local-sign-orientation-character}
A rank-one crossing changes its Pfaffian section by \(-1\) and
has determinant \(u^2\).  If all three simple walls have these
comparisons and they assemble to the coherent character
\(\epsilon_o\), then \(\epsilon_o=\det_W\) on the type-II Weyl
group.  Its automorphic target is the homogeneous orthogonal
character \(\widetilde\nu_{\Delta_5}|_W\).
\end{proposition}
\begin{proof}
The two-by-two Pfaffian is \(u\), so its pulled-back expression
is \(-u\) and its square is \(u^2\).  Section compatibility forces
the root lift to have residual sign \(-1\) in the reference frames.
The three generator values determine
the character, as proved in the preceding subsection.  A real path
crossing \(u=0\) is not a path in the wall complement; the stipulated
lift is part of the datum.
\end{proof}

The abstract sign set
\[
 \Omega_{12}=\{\pm1\}^{12}
\]
has cardinality \(2^{12}=4096\).  To interpret
this number as a count of geometric wall configurations, one needs
a bijection from a complete geometric atlas to that set, including
stabilizer weights and identifications on overlaps.  The finite
sign sum alone does not construct such an atlas.

\subsubsection*{From the finite product to the compact section}

Local finiteness in positive height makes the products compatible
as jets: only finitely many factors contribute to any fixed height.
This gives a formal product.  To obtain the holomorphic automorphic
section one also needs convergence on a cusp domain, the parameter
and line comparisons, and analytic continuation with the stated
boundary behaviour.  Under these assumptions the target product is
\(\Delta_5\), with normalization
\[
 \Delta_5=64q^{1/2}r^{1/2}s^{1/2}(1+\cdots),
 \qquad \Delta_5^2=\chi_{10},
\]
where the complete boundary factors are those of
Chapter~\ref{ch:view1-automorphic}.

If an independently defined protected trace satisfies the geometric
sewing identity
\[
 Z_X=-4096\,\iota_{\mathrm{aut}}(\operatorname{pf}_X)^{-2},
\]
then the section comparison gives
\(Z_X=-4096\,\kappa^*(\Delta_5^{-2})\).
The trace and its sewing identity remain additional geometric data.
The scalar branch sign in this formula does not determine
\(\epsilon_o\), and the finite algebraic skew blocks do not alone
construct the compact state space or its operator.


\subsection{The source bar coalgebra and the compact comparison}
\label{subsec:koszul-source}

The hybrid correspondences must first define an algebra in a specified
coefficient category. Only then does its associative bar have a
meaning. We use the unitalized Ran category
$(\mathcal R_E^+,\star,K)$ of Section~\ref{sec:bkm-ran-envelope},
with right differential-module coefficients, cross-block chiral
convolution and empty unit $K$. A realization of constructible Hall
coefficients in this category is additional geometric data.
It must carry the products, orientation transport and higher collision
coherences, including wrapped and mixed configurations.

Write $V_E=\mathsf{Den}_{\Delta_5,E}$ for the BD current envelope
and $\mathcal A_{\Delta,E}=U_\star(\iota_*V_E)$ for its associative
Ran envelope. The associative comparison in this subsection has
$\mathcal A_{\Delta,E}$ as its target. A comparison directly with
$V_E$ as a BD chiral algebra is a different statement and needs the
corresponding chiral functor. The bar--cobar equivalence in
Proposition~\ref{prop:bkm-ran-bar-envelope} does not identify these
two recipients.

\begin{lemma}[Source and target bars]
\label{lem:source-target-koszul-separation}
For an augmented associative target $T$ with $T_\varnothing=\C$,
the counit $\Omega BT\to T$ is an internal equivalence in the
specified Ran category. It constructs neither a compact-source
algebra nor a morphism from such an algebra to $T$.
\end{lemma}
\begin{proof}
The counit depends functorially on the multiplication and
augmentation of $T$. Its construction has no input consisting of
compact correspondences, a realization of their coefficients,
orientations or a source-to-target morphism. In particular,
$C_X:=BT$ would define a coalgebra from $T$, without identifying it
with a coalgebra obtained from the compact geometry.
\end{proof}

\subsubsection*{Coefficient realization and the source coalgebra}

At a stage $R$, denote the realized hybrid algebra by $A_{X,R}$.
The index may bound geometric complexity or configuration arity.
It does not by itself specify a root-height quotient. Whenever a
target stage $T_R$ is used, it must be an actual augmented
associative object with supplied transition maps. A projection to
symmetric positive and negative root degrees is not such an object,
by Proposition~\ref{prop:bkm-positive-completion}. Identification
of the limit of the $T_R$ with $\mathcal A_{\Delta,E}$ is a
separate comparison.

\begin{definition}[Chiral Koszul source datum]
\label{def:chiral-koszul-source-datum}
For fixed stability, support and stage data $(\sigma,S,R)$, a
source datum $\mathfrak K^{\mathrm{Kos}}_{X,R}$ consists of the
following objects and morphisms in $\mathcal R_E^+$.
\begin{enumerate}
\item A realization of
$\mathcal F^{\mathrm{hyb}}_{X,\sigma,S,\leq R}$ as an augmented
associative algebra $A_{X,R}$ with empty component $\C$.
Its multiplication and all coherent collision operations come from
the supplied hybrid correspondences and orientations.
\item A coaugmented counital coassociative coalgebra $C_{X,R}$
with empty component $\C$, differential $d_C$, coproduct
$\Delta_R^{\mathrm{ch}}$ and an exhaustive filtration
$F^{\mathrm{co}}_{R,\bullet}$. The reduced coproduct is locally
nilpotent on each finite configuration and retained filtration
piece. If a global finite bound on that stage is used, it is included
in the datum.
\item A morphism of coalgebras
$b_{X,R}:C_{X,R}\to BA_{X,R}$ compatible with charge descent,
coaugmentation, counit and the supplied geometric operations.
Its cone is acyclic. The geometric construction of $C_{X,R}$
and this comparison are specified on local, wrapped and mixed
inputs, with all higher-colored collision coherences.
\item Compatible transition morphisms on these objects. Their
coherent diagrams include the reduced $E$-quotient and the
normal-ordered charge homomorphism $\overline\Pi_X$.
\end{enumerate}
A source-to-target comparison additionally supplies an augmented
associative morphism $\vartheta_R:A_{X,R}\to T_R$ and proves
its cone acyclic. Define
\[
 \varepsilon^{\mathrm{src}}_{\mathrm{bar},R}
 :=\varepsilon_{A_{X,R}}\circ\Omega b_{X,R}:
 \Omega C_{X,R}\longrightarrow A_{X,R},
 \qquad
 \Theta_{\mathrm{Kos},R}
 :=\vartheta_R\circ\varepsilon^{\mathrm{src}}_{\mathrm{bar},R}.
\]
The symbols $B,\Omega$ in this datum always refer to $\star$ and $K$.
\end{definition}

The finite binary flag atlas supplies particular associativity
and coproduct diagrams. It does not replace the higher-colored,
unit, refinement and quotient coherences required by this
definition. Nor does the projection-finite local sector provide the
coefficient realization of wrapped inputs. A restriction to ordered
words records the word-cut part of $C_{X,R}$. Comparison with the
full Ran coalgebra must preserve every collision operation.

The obstruction notation can retain these separate requirements:
\[
\begin{split}
 \mathfrak O_{\mathrm{Kos},R}=
 (&\mathfrak o^C_R,\mathfrak o^{\mathrm{fil}}_R,
 \mathfrak o^{\mathrm{conil}}_R,\mathfrak o^{\Delta,\mathrm{ch}}_R,
 \{\mathfrak o^{C,\mathrm{tr}}_{R'R}\},
 H^\bullet\operatorname{Cone}b_{X,R},\mathfrak o^\Omega_R,\\
 &\mathfrak o^{\mathrm{wrKos}}_R,\mathfrak o^{\mathrm{hyb}}_R,
 \mathfrak o^\Pi_R,\mathfrak o^{\mathrm{sep}}_R,
 \mathfrak o^{\Prim}_R).
\end{split}
\]
Here the first entries mean the stated coalgebra, filtration and
transition identities. $\mathfrak o^{\mathrm{wrKos}}_R$ includes
the wrapped coefficient realization and its bar domain.
The last entries require compatible charge descent, separation of
source from target, and any specified ordinary primitive realization.
An absent object or map is an unresolved construction, rather than
a numerical value of this tuple.

\begin{proposition}[Finite-length source bar]
\label{prop:canonical-source-bar-coalgebra}
Suppose the realized augmented algebra $A_{X,R}$ is supplied and
its positive augmentation ideal $a_{X,R}$ has no nonzero retained
$\star$-word of length greater than $N_R$. Then its normalized
bar coalgebra is
\[
 C^{\mathrm{bar}}_{X,R}=BA_{X,R}
 =K\oplus\bigoplus_{1\leq p\leq N_R}(sa_{X,R})^{\star p},
\]
with the internal and multiplication differentials. The suspension
$s$ is the bar suspension. Multiplication has degree one after
suspension. The coproduct is deconcatenation:
\[
 \Delta[x_1|\cdots|x_p]
 =1\otimes[x_1|\cdots|x_p]+[x_1|\cdots|x_p]\otimes1
  +\sum_{i=1}^{p-1}[x_1|\cdots|x_i]\otimes[x_{i+1}|\cdots|x_p].
\]
It is conilpotent and supplies a source-bar model with $b_{X,R}$
the identity.

Without a uniform $N_R$, the same construction is
finite at each configuration by the positive arity bound.
\end{proposition}
\begin{proof}
The normalized associative bar construction in this tensor category
has the displayed terms. Associativity with its coherent form gives
the square-zero bar differential. Cutting words gives a counital
coassociative coproduct compatible with it. Each reduced cut lowers
the available length, so more than $N_R$ nonempty blocks cannot
occur.

On a configuration $I$, positivity gives the bound $|I|$
instead. These are coefficient-category statements. To identify
the coproduct with geometric cutting correspondences requires the
realization and coherence maps stipulated in the source datum.
\end{proof}

If all nonempty geometric colors have positive HN height, the
bound on total HN height can imply $N_R\leq R$. That is a hypothesis
on the retained geometric algebra, not a conclusion about the full
two-sided root algebra. Additivity of $\overline\Pi_X$ ensures that
a bar word, each cut and each multiplication term have the same
descended Gram degree. This proves compatibility of gradings. It
constructs neither the realization nor the comparison $\vartheta_R$.

\begin{corollary}[The remaining source-to-target map]
\label{cor:source-bar-remaining-cobar-obstruction}
For a source datum in the specified Ran category,
$\varepsilon^{\mathrm{src}}_{\mathrm{bar},R}$ is an equivalence.
Hence $\Theta_{\mathrm{Kos},R}$ is an equivalence exactly when
$\vartheta_R$ is. A bar construction in another coefficient category
requires its own counit theorem before this conclusion applies.
\end{corollary}
\begin{proof}
The coalgebra map $b_{X,R}$ is an equivalence by assumption.
Apply the associative cobar and the bar--cobar counit of
Proposition~\ref{prop:bkm-ran-bar-envelope}. The two-out-of-three
property proves the assertion. A target-internal counit supplies
no part of $\vartheta_R$.
\end{proof}

For a strict representative, the comparison equations include
\[
 d_T\vartheta_R=\vartheta_Rd_X,\quad
 \vartheta_R\mu_X=\mu_T(\vartheta_R\star\vartheta_R),\quad
 \vartheta_R\eta_X=\eta_T,\quad
 \vartheta_R\rho^X_{R'R}=\rho^T_{R'R}\vartheta_{R'}.
\]
For coherent representatives, these equations are replaced by the
specified coherent homotopies. In either case the full coefficient
cone $\operatorname{Cone}\vartheta_R$ must be acyclic.
The Lie--commutative chiral theorem of
\cite[Theorems~5.1.1--5.2.1]{FrancisGaitsgoryCKD} does not supply
this compact associative comparison. It enters the envelope
construction at precisely the scope proved in
Proposition~\ref{prop:bkm-ran-bar-envelope}.

\subsubsection*{Compatibility with the geometric structures}

\begin{definition}[Finite Koszul comparison datum]
\label{def:finite-koszul-comparison-datum}
A comparison datum $\mathfrak K^{\mathrm{cmp}}_{X,R}$ supplements
the source datum by the following maps and commuting diagrams.
\begin{enumerate}
\item[\textup{(K0)}] The source counit and $\vartheta_R$ have
acyclic cones. If finite filtrations are used to prove this, the
associated-graded cone comparisons are acyclic at every retained
coefficient and parity.
\item[\textup{(K1)}] Every specified Weyl transport is intertwined
by $\Theta_{\mathrm{Kos},R}$. A reflection may map between chambers
or require a larger finite operation domain. Both domain and
codomain are specified. In chain representatives the equation is
$\Theta\tau_i^X-\tau_i^T\Theta=dH_i^W+H_i^Wd$.
The $H_i^W$ carry the higher coherences for their defined
compositions. No finite symmetric root set is assumed invariant
under the whole Weyl group.
\item[\textup{(K2)}] On the perfect complexes defining the geometric
Pfaffian and determinant lines, a supplied line isomorphism
intertwines orientation transport. Its type-II reflection sign is
$\epsilon_o(s_{\delta_i})=\widetilde\nu_{\Delta_5}(s_{\delta_i})=-1$.
The wall and chamber isomorphisms satisfy the cocycle identities of
$(O1)^+$ and $(O2_{\mathrm{atlas}})$.
\item[\textup{(K3)}] The associative map preserves products and
units, and its bar preserves coproducts, coaugmentations and
counits. If an ordinary Hall/Hopf realization is used to extract
primitives, it is separately supplied and its operation diagrams
commute. The invariant Lie form on those primitives is intertwined
with $B_{\mathfrak g}$. It is not identified with a pairing adjoint
to the ordinary primitive Hopf coproduct. Such an identification
is excluded by Proposition~\ref{prop:bkm-hopf-pairing-obstruction}.
\item[\textup{(K4)}] The primitive comparison carries the invariant source Lie radical
onto the target radical. It induces the comparison on their quotients.
When the target is $\mathfrak g$, its invariant form has zero
radical. Kernel and quotient squares are required in the actual
coefficient realization.
\item[\textup{(K5)}] This primitive comparison is the same map
$\iota_R^{\Prim}$ used in recognition. The induced maps of ordinary
enveloping algebras preserve PBW length and the finite weighted
pieces of Proposition~\ref{prop:bkm-finite-filtered-pbw}.
\item[\textup{(K6)}] Stage transitions and all the preceding
structures form coherent diagrams. In a strict chain model,
$\Theta_R\Omega(\rho^C_{R'R})-\rho^T_{R'R}\Theta_{R'}
=dH^{\mathrm{tr}}_{R'R}+H^{\mathrm{tr}}_{R'R}d$, with the
composition homotopies also specified.
\item[\textup{(K7)}] Completion uses the homotopy inverse limit
of these diagrams. To use ordinary inverse limits, require degreewise Mittag--Leffler
chain groups for the objects and comparison cones. Alternatively,
supply another explicit comparison between ordinary and homotopy
limits. Bar and cobar
commute with the chosen completed construction only when that
additional comparison has been proved.
\end{enumerate}
The notation $\mathfrak O_{\mathrm{Kcmp},R}$ denotes the two cone
cohomologies together with the Weyl, line, operation, radical, PBW
and transition defects in these diagrams. Their vanishing includes
the supplied coherent nullhomotopies, not just equal matrix ranks.
\end{definition}

\begin{proposition}[Compatible finite comparison]
\label{prop:finite-koszul-comparison-compatibility}
Under \textup{(K0)--(K6)}, the map
$\Theta_{\mathrm{Kos},R}:\Omega C_{X,R}\to T_R$ is an augmented
associative equivalence compatible with the specified Weyl,
orientation, primitive, invariant-form and enveloping diagrams.
\end{proposition}
\begin{proof}
Clause (K0) and the source counit give the equivalence.
Each remaining compatibility is a supplied commuting square or
coherent homotopy, and composition with the counit preserves it.
The radical comparison descends to quotient spaces. The ordinary
enveloping universal property gives the PBW comparisons.
This conclusion uses the stated realizations of these structures,
not an automatic extraction of a Hall object from the target.
\end{proof}

\begin{proposition}[Agreement on selected generators]
\label{prop:primitive-restriction-not-koszul}
Agreement with a prescribed map on a selected family of root
vectors does not imply an equivalence of associative coefficient
objects. For a supplied comparison $\vartheta_R$, the remaining
condition is acyclicity of its full coefficient cone.
\end{proposition}
\begin{proof}
The failure already occurs for ordinary augmented dg algebras.
Let $B$ be such an algebra and let $M$ be a dg $B$-bimodule with
nonzero cohomology. The square-zero extension $A=B\oplus M$ has
\[
 (b,m)(b',m')=(bb',bm'+mb'),\qquad M^2=0.
\]
Its augmented projection $p:A\to B$ fixes every selected vector
in the copy of $B$, but
$\operatorname{Cone}(p)\simeq M[1]$ is not acyclic.
This example asserts neither a connected Hopf structure on $A$
nor agreement on all its primitives. Under those additional
hypotheses, Proposition~\ref{prop:bkm-connected-hopf-recognition}
gives the positive ordinary Hopf recognition theorem.
\end{proof}

To establish the analogous conclusion in a chiral coefficient
category one must inspect the local comparison cone, not merely
the ordinary example. Additional wrapped, mixed or descendant
coefficients can lie outside the selected root-vector family.
The invariant form and the ordinary primitive map likewise need
the specified realizations in (K3)--(K5).

\begin{proposition}[Passage to completion]
\label{prop:chiral-koszul-source-datum-consequence}
Suppose source and comparison data are supplied on a cofinal
sequence of stages, with the coherent transitions of (K6).
Then their homotopy inverse limits give equivalences
\[
 \operatorname*{holim}_R C_{X,R}
       \simeq\operatorname*{holim}_R BA_{X,R},\qquad
 \operatorname*{holim}_R\Omega C_{X,R}
       \simeq\operatorname*{holim}_R T_R,
\]
with all specified compatible operations.
Under (K7), the same assertions hold for the ordinary inverse
limits of the chosen chain models. Writing these limits as
$B A_X$, $\Omega C_X$ and $\mathcal A_{\Delta,E}$ additionally
requires the corresponding bar, cobar and target-limit comparisons.
\end{proposition}
\begin{proof}
A homotopy limit is a functor on the coherent diagram category and
preserves equivalences. Apply it to $b_{X,R}$ and
$\Theta_{\mathrm{Kos},R}$ and their compatibility diagrams.
The ordinary-limit assertion follows from
Proposition~\ref{prop:bkm-limit-complexes}.
For actual degreewise finite-dimensional towers, the required
Mittag--Leffler property is automatic once their transition maps
exist. In general it is not automatic, as the polynomial-tail
example there shows. Nothing in preservation of equivalences by
limits alone says that bar or cobar commutes with that limit.
\end{proof}

\subsubsection*{Twisting cochains and the boundary algebra}

A twisting cochain from $C_{X,R}$ to an augmented algebra has
unsuspended degree one. It satisfies
\[
 d\tau+\tau d+\mu(\tau\star\tau)\overline\Delta=0,
\]
with the suspension signs fixed by the bar convention.
Equivalently it is a morphism from the cobar. For a bar source it
includes the canonical twisting cochain. This equation describes
coupling data at the specified carrier and tensor. It does not
identify the bar coalgebra with the bulk algebra.

A proposed bulk center belongs to the associative boundary algebra
$A_b=R\operatorname{End}(b)$ and requires its separate map and
hypotheses, as in Remark~\ref{rem:bkm-universal-holography-scope}.
The empty coefficient $K$, the BD vacuum, the augmentation of
$A_{X,R}$ and a prepared state are not interchangeable. In
particular, obtaining a completed coefficient equivalence to
$\mathcal A_{\Delta,E}$ still leaves the state comparison and its
protected trace to be constructed. The geometric primitive
recognition in the next subsection is another operation-preserving
comparison, with its own explicit source data.


\subsection{Primitive recognition}
\label{subsec:primitive-recognition}

The compact recognition problem asks for a graded Lie superalgebra
isomorphism from geometric primitives to the denominator algebra
$\mathfrak g=\mathfrak g_{\Delta_5}$ of
Chapter~\ref{ch:view2-bkm}. Its first interior comparison must identify
$29$ even and $93$ odd directions and their brackets. The scalar
$-64$ specifies only their difference.

We first work with algebraic direct sums over root degrees. Positive
height completion is available by
Proposition~\ref{prop:bkm-positive-completion}. A completion of both
signs must be separately defined. In particular, the full product of
positive and negative root spaces is not a topological Lie algebra
with the coordinatewise topology. This distinction applies equally
to the proposed compact source and to the target.

\subsubsection*{The source and its presentation map}

Assume the hybrid geometry supplies an ordinary Hall super-bialgebra
$H_X$ in a specified coefficient realization. Its associative
product $m_X$, unit and counital coproduct $\Delta_X$ satisfy
\[
 \Delta_Xm_X=(m_X\otimes m_X)(1\otimes\tau\otimes1)
                         (\Delta_X\otimes\Delta_X),
\]
where $\tau$ is the Koszul braiding. Its actual primitives form a
Lie superalgebra under the supercommutator. If a cochain realization
is used, all differentials and the passage to protected cohomology
must preserve the operations needed for this assertion.

Denote the descended primitive algebra by $P_X^{\Pi,\mathrm{desc}}$.
A formula such as
\[
 P_X^{\Pi,\mathrm{desc}}=
 H^0\Prim_{\mathrm{prot}}(\overline\Pi_{X,*}^{\Theta}\mathcal F_X)
\]
means that the coefficient realization, primitive functor and charge
descent have been constructed and shown to commute in that formula.
It is not a definition of those maps from the charge labels alone.

Let $B_X$ be an independent even invariant Lie form on this algebra,
and set
\[
 R_X=\operatorname{Rad}B_X,\qquad
 P_X^{\Pi,\mathrm{red}}=P_X^{\Pi,\mathrm{desc}}/R_X.
\]
Invariance makes $R_X$ a Lie ideal. Its generated ideal in the
ordinary enveloping algebra is a Hopf ideal by
Proposition~\ref{prop:bkm-lie-radical-quotient}.
The invariant form is not required to be adjoint to the ordinary
primitive Hopf coproduct. That requirement would contradict the
real $\mathfrak{sl}_2$ sector.

Let $\mathfrak F$ be the free Lie superalgebra on a basis of
$\mathfrak h=L\otimes\C$ and the symbols $e_i,f_i$ with their
prescribed parities. Let $J_{\mathrm{BK}}$ be the ideal generated
by Definition~\ref{def:bkm-relations}, including the linear
identification $h_i=\alpha_i\in\mathfrak h$. Thus
$\mathfrak F/J_{\mathrm{BK}}=\mathfrak g$.
This is the Gritsenko--Nikulin presentation
\cite[Section~3 and Statement~6.4]{GN}, with its product and lattice
forms in \cite{GNPartI,GNII,BorcherdsGKM88,KAC:1}.
No further radical quotient of $\mathfrak g$ is needed: its invariant
form is already nondegenerate.

\begin{theorem}[Recognition from geometric representatives]
\label{thm:primitive-recognition}
Suppose the following data are supplied over $\C$.
\begin{enumerate}
\item[\textup{(S1)}] The source primitive algebra and its
normal-ordered descent are constructed. Their bracket is graded by
$\Gamma_{\mathrm{gram}}$, and each descended homogeneous parity
space is finite-dimensional.
\item[\textup{(S2)}] An even identification
$P^{\Pi,\mathrm{red}}_{X,0}\simeq\mathfrak h$ matches the source
degrees with $L$ through $\alpha(n,l,m)$.
\item[\textup{(S3)}] For every simple index $i$, the source supplies $e_i^X,f_i^X$
in degrees $\pm\alpha^{-1}(\alpha_i)$ and parity $p(i)$.
It also supplies $h_i^X$ corresponding to $\alpha_i$ in the Cartan. The imaginary
fibers have the prescribed $\mu_{\bar p}(a)$ generators.
\item[\textup{(S4)}] These representatives satisfy the Cartan,
Chevalley, real Serre, orthogonality and super sign relations of
Definition~\ref{def:bkm-relations} in the source bracket.
\item[\textup{(S5)}] The invariant source Lie form descends to
$P_X^{\Pi,\mathrm{red}}$, and the presentation map preserves it
and $B_{\mathfrak g}$. The radical and quotient maps are actual
maps of the supplied source objects.
\item[\textup{(S6)}] For the induced map
$\pi:\mathfrak F\to P_X^{\Pi,\mathrm{red}}$, one has
$\ker\pi=J_{\mathrm{BK}}$.
\item[\textup{(S7)}] The Cartan and simple representatives
generate $P_X^{\Pi,\mathrm{red}}$ by finite bracket words.
\item[\textup{(S8)}] The source and target ordinary enveloping
algebras carry their PBW length filtrations and the finite weighted
pieces of Proposition~\ref{prop:bkm-finite-filtered-pbw}. The
comparison uses these actual algebras, not envelopes of symmetric
root fragments.
\item[\textup{(S9)}] Whenever a topological conclusion is desired,
a completion domain and continuous operations are specified on
both sides. The comparison and its inverse extend continuously.
For positive height this uses the genuine positive Lie quotients.
For another domain it is a separate hypothesis.
\item[\textup{(S10)}] The source parity dimensions agree with
$e_\beta,o_\beta$ in every positive and negative root degree,
and vanish in non-root degrees. At $\delta_{123}$ they are
$29|93$, not a choice inferred from the signed value $-64$.
\end{enumerate}
Then $\mathfrak g\to P_X^{\Pi,\mathrm{red}}$ is a graded Lie
isomorphism preserving the invariant form. It induces a filtered
Hopf isomorphism of ordinary enveloping algebras. Under (S9), it
also identifies the specified completions. Its superdenominator
is $64^{-1}\Delta_5(2Z)$.
\end{theorem}
\begin{proof}
The representatives and (S4) factor $\pi$ through $\mathfrak g$.
Clause (S6) gives injectivity and (S7) gives surjectivity.
The degree, parity and form conditions make this the claimed
structured isomorphism.

The ordinary enveloping universal property
and PBW give the filtered Hopf assertion. Continuity of both maps
proves the completed statement in (S9). The denominator follows
from Theorem~\ref{thm:bkm-denominator-identity}.

In particular, (S10) is a consequence of the established
isomorphism, while its finite parity tests can be used in proving
injectivity or surjectivity. It is not a substitute for them.
\end{proof}

For a real pair the required Serre exponent is three.
At $a_{ij}=\delta_i+\delta_j$, the source must supply nine even
null simple representatives in addition to $[e_i^X,e_j^X]$.
For the complementary real index $k$, their exponent is five.
The relations at $\delta_{123}$ also restrict possible source
differentials: Theorem~\ref{thm:bkm-serre-nonextension} forbids the
rank-$28$ cancellation on the unchanged presentation beyond height
five. A derived replacement must therefore identify its changed
relations and its comparison map.

If the independent coefficient data of
Definition~\ref{def:finite-koszul-comparison-datum} are supplied,
their primitive comparison can be required to be the inverse of the
isomorphism above. The resulting statement concerns the associative
Ran envelope $\mathcal A_{\Delta,E}$ in its specified category.
It does not follow solely from the ordinary primitive isomorphism.

\subsubsection*{Finite root sets and geometric support}

\begin{definition}[Cofinal finite recognition data]
\label{def:cofinal-primitive-recognition-datum}
For $\nu\geq1$ set
\[
 W_\nu=\{\beta\in\mathcal R_+:\operatorname{ht}\beta\leq\nu\},
 \qquad S_{W_\nu}=W_\nu\cup(-W_\nu)\cup\{0\},
 \qquad \Gamma(W_\nu)=\alpha^{-1}(W_\nu).
\]
Finite recognition data on an exhaustion consist of the following
compatible homogeneous spaces and maps.
\begin{enumerate}
\item[\textup{(R0)}] Geometric source spaces with specified compact
supports, coefficient summands, orientation data and charge maps.
Their basis vectors are classes of these objects.
\item[\textup{(R1)}] The simple representatives and Cartan maps
in each retained degree, with the prescribed parities.
\item[\textup{(R2)}] The homogeneous bracket maps and the defining
relation identities. Each iterated identity is evaluated in a
finite enlargement retaining all intermediate degrees.
\item[\textup{(R3)}] Even bijections
$A_{\beta,\bar p}:P^{\Pi,\mathrm{red}}_{X,\alpha^{-1}(\beta),\bar p}
\to\mathfrak g_{\beta,\bar p}$ for all retained root degrees in both signs and in degree zero.
Require the source spaces to vanish in every non-root degree.
Include these zero-space tests in the finite degree exhaustion.
\item[\textup{(R4)}] Opposite-degree invariant Lie pairings, their
radicals and quotient maps. The comparison intertwines the pairings
and brackets. Where ordinary Hall products and coproducts are
used to define primitives, their bialgebra identities hold in
finite operation diagrams with the necessary output spaces.
\item[\textup{(R5)}] No extra relation in each bounded word space:
$\ker(\pi|_{F_N\mathfrak F})=J_{\mathrm{BK}}\cap F_N\mathfrak F$.
Here $F_N$ bounds total generator weight, assigning weight one to
Cartan generators and $\operatorname{ht}\alpha_i$ to $e_i,f_i$.
This finite-dimensional word bound is different from fixing only
the total root degree.
\item[\textup{(R6)}] Every retained source vector is a finite
bracket-word combination of the supplied representatives, in some
finite enlargement containing its intermediate degrees.
\item[\textup{(R7)}] The ordinary enveloping comparison preserves
PBW length and the finite weighted pieces. Its maps are the ones
induced by the primitive comparison.
\item[\textup{(R8)}] Restriction and enlargement maps agree on
common homogeneous operations. For a completed statement, actual
quotient or inverse-system maps are supplied in a completion domain
with continuous operations, as in (S9).
\end{enumerate}
The maps $A_{\beta,\bar p}$ must commute with every retained bracket
and pairing square. Root sets carry partial operations as in
Definition~\ref{def:bkm-finite-target-presentation-datum}. No
zero extension of omitted brackets is imposed.
\end{definition}

A bounded root degree alone does not make the free presentation
finite: positive and negative generators can cancel in arbitrarily
long words. The total generator-weight bound in (R5) does give
finite-dimensional word spaces. The intersection
$J_{\mathrm{BK}}\cap F_N\mathfrak F$ is the actual ideal
intersection, not merely the span of defining relators whose
written degrees happen to lie in a root set.

\begin{proposition}[Height sets and their inverse degree map]
\label{prop:height-root-windows-gamma}
The $W_\nu$ are finite, downward saturated and cofinal in
$\mathcal R_+$. Their Gram preimages are
\[
 \Gamma(W_\nu)=
 \{(b_1,b_1+b_2-b_3,b_2):\textstyle\sum_i b_i\delta_i\in W_\nu\}.
\]
They lie in the product cusp support.
\end{proposition}
\begin{proof}
There are finitely many nonnegative triples with sum at most $\nu$.
A proper positive subdegree has no greater height. Every positive
root has finite height.

Solving the three coefficients in
\eqref{eq:bkm-alpha-map} gives the displayed inverse.
If $b_2>0$ then $m>0$. If $b_2=0,b_1>0$ then $m=0,n>0$.
If $b_1=b_2=0,b_3>0$ then $n=m=0,l<0$.
\end{proof}

\begin{corollary}[Exhaustion of positive roots]
\label{prop:active-root-window-exhaustion}
Every positive root $\beta$ lies in $W_{\operatorname{ht}\beta}$,
and $\bigcup_\nu W_\nu=\mathcal R_+$.
\end{corollary}
\begin{proof}
Finiteness of height proves the assertion. A root of zero
superdimension still occurs, since roots were defined by their spaces.
\end{proof}

\begin{definition}[Geometric support exhaustion]
\label{def:source-window-compact-support-exhaustion}
A discrete geometric support exhaustion uses a set
$\mathscr S_X^{\mathrm{cpt}}$ of compact support types and a height
$h_X:\mathscr S_X^{\mathrm{cpt}}\to\N$ with finite fibers. Its stages are
$\mathscr S_{X,\leq R}^{\mathrm{cpt}}=\{s:h_X(s)\leq R\}$.
Each type carries its actual source stratum, proper support,
stabilizer and orientation data. Geometric transition maps between
the associated moduli or coefficient objects are also specified.
Inclusions of the index sets do not alone give these maps.
\end{definition}

\begin{proposition}[Support exhaustion]
\label{prop:source-window-compact-support-criterion}
Such an exhaustion contains every retained support type in a
finite stage. It does not assert that the union of the corresponding
moduli has finite dimension or finitely many points.
\end{proposition}
\begin{proof}
The type $s$ occurs at $R=h_X(s)$. Finiteness concerns the index
set of types, not the geometric objects they index.
\end{proof}

\begin{definition}[Compatibility of source and root degrees]
\label{def:target-source-window-compatibility}
For each geometric stage, specify its normal-ordered degree map
$d_R:\mathscr S_{X,\leq R}^{\mathrm{cpt}}\to\Gamma_{\mathrm{gram}}$
and an index $\nu(R)$ such that
$\alpha(d_R(s))\in S_{W_{\nu(R)}}$ whenever that support contributes
to the primitive comparison. The maps agree on shared support types
as $R$ increases, and $\nu(R)$ can be chosen nondecreasing.
\end{definition}

\begin{proposition}[Degree-compatible finite comparison]
\label{prop:target-source-window-compatibility-criterion}
These data place each retained geometric primitive space in a
specified target degree and preserve that assignment under transition.
\end{proposition}
\begin{proof}
Apply the defining equality with $\alpha$ at each support type.
Compatibility of $d_R$ makes the same equality hold in every later
stage. Negative degrees and the Cartan use $-W_\nu$ and $\{0\}$.
\end{proof}

The existence of a geometric support exhaustion and its operation
maps is part of the compact construction. The target height
exhaustion supplies no such geometric maps. In a proposed source
whose support cannot be indexed with finite fibers, another finite
coefficient construction must replace this discrete formulation.
That limitation concerns this exhaustion method, not existence of
all possible compact sources.

\subsubsection*{Recognition in finite homogeneous coordinates}

\begin{proposition}[Target operation spaces]
\label{prop:target-presentation-prerequisite}
A finite comparison with $\mathfrak g$ requires its two parity
spaces, bracket maps, opposite-degree invariant pairings and
ordinary enveloping pieces on the retained operation domains.
The signed dimensions $f_\beta$ or simple coefficients $m(\beta)$
alone do not define these maps.
\end{proposition}
\begin{proof}
A bracket square has a vector space at each input and at its
output, and a pairing square has the specified opposite-degree
spaces. A scalar dimension cannot supply their structure maps.
The parity recursion constructs dimensions for the target
presentation. Bases and operations still come from that presentation.
\end{proof}

\begin{proposition}[Finite homogeneous recognition]
\label{prop:finite-window-presentation-recognition}
Finite recognition data determine an isomorphism of the retained
homogeneous operation diagrams
\[
 A_W:P_X^{\Pi,\mathrm{red}}|_{\alpha^{-1}(S_W)}
       \xrightarrow{\sim}\mathfrak g|_{S_W}.
\]
It preserves the Cartan, simple representatives, retained brackets
and invariant pairings. Cofinal compatible data identify the full
algebraic graded Lie algebras.
\end{proposition}
\begin{proof}
Each parity block of $A_W$ is bijective by (R3), and the bracket
and pairing squares commute by (R2)--(R4). Their inverses preserve
the same maps. Every homogeneous operation occurs in some finite
enlargement, so the compatible family identifies all operations
on the direct sums. Equivalently, (R5)--(R6) prove injectivity and
surjectivity of the source presentation map. This is
Theorem~\ref{thm:bkm-finite-operation-reconstruction}. No enveloping
algebra of a symmetric root fragment is used.
\end{proof}

\begin{definition}[Finite source matrices]
\label{def:finite-source-proof-tables}
Choose homogeneous bases of the geometric coefficient and primitive
spaces. Denote the ordinary Hall product, coproduct, bracket and
invariant Lie form by matrices $M,D,B,G$, and the unit and counit
by $\eta,\epsilon$. These matrices act on their specified input
and output spaces, with enlargement when a product increases the
retained weight. Kernel inclusions, quotient maps and source-to-target
maps are denoted $K,Q,A$.

The coefficient basis is obtained from the supplied compact support,
vanishing-cycle or intersection-cohomology summand, reduced and
quotient orientations, finite-stabilizer action and protected
integration. Thom--Sebastiani and transition isomorphisms specify
how these coefficients enter the Hall operations. These geometric
objects and maps precede their coordinate matrices.

Require associativity of $M$, coassociativity of $D$, and the left
and right unit and counit laws. Require also
\[
 DM=(M\otimes M)(1\otimes\tau\otimes1)(D\otimes D),
 \qquad B=M-M\tau\quad\hbox{on primitives},
\]
and the invariant-form identities
\[
 G([x,y],z)=G(x,[y,z]),\qquad
 G(x,y)=(-1)^{|x||y|}G(y,x).
\]
The comparison squares include
\[
 A_{\beta+\gamma}B^X_{\beta,\gamma}
 =B^{\mathfrak g}_{\beta,\gamma}(A_\beta\otimes A_\gamma),
 \qquad
 G^X_\beta=A_\beta^{\mathrm t}G^{\mathfrak g}_\beta A_{-\beta}
\]
in homogeneous bases, with the super sign convention on reversed
arguments.

Radical inclusion and quotient matrices satisfy
$\operatorname{im}K=\ker G$, $QK=0$, and induce a nondegenerate
form on the quotient. The relation matrices express the defining
Cartan, Chevalley, Serre and orthogonality identities. Their actual
ideal intersections enter (R5). The bracket-word spans enter (R6).

The induced enveloping matrices use the super PBW relations and
finite weighted pieces. They need not satisfy Hopf adjointness
with the same matrix $G$. Proposition~\ref{prop:bkm-hopf-pairing-obstruction}
proves that the latter requirement is incompatible with this target.
\end{definition}

\begin{proposition}[Coordinate form of recognition]
\label{prop:finite-source-proof-tables-equivalence}
Over $\C$, the finite source spaces and maps of (R0)--(R8) are
equivalently specified by the matrices above. Their geometric
construction, exact kernel equalities and commuting operation
diagrams remain part of the specification. This equivalence is independent of the chosen
homogeneous bases.
\end{proposition}
\begin{proof}
A choice of basis turns every given linear map into a matrix.
Composition, kernel and intertwining identities are unchanged by
conjugating the source and target bases. Conversely, matrices on
the specified geometric spaces define those maps when they satisfy
the stated identities. Equality of a proposed relation subspace
with a kernel requires inclusion and equality of dimensions, not
only two unrelated ranks. Invertibility of each comparison block
gives its actual inverse. Bare matrices of the right sizes do not
identify their underlying spaces with geometric coefficients.
\end{proof}

\begin{proposition}[Cofinal recognition]
\label{prop:global-recognition-cofinal-datum}
Suppose the algebraic source and its presentation map are
constructed. Conditions (S1)--(S8) and (S10) can be checked on
cofinal finite homogeneous operation domains and bounded word spaces.
A completed recognition statement additionally requires (S9).
\end{proposition}
\begin{proof}
Every algebraic vector and relation is a finite linear combination
of words, and hence occurs at some total generator-weight bound.
Every bracket or pairing square involves finitely many degrees.
Thus the finite data exhaust the defining maps and identities.
Conversely, a global graded isomorphism and its specified geometric
source restrict to these diagrams and preserve their kernels and
parities. A topology and convergence of its operations are not
consequences of an algebraic finite-word assertion, so the completed
conclusion retains (S9).
\end{proof}

\subsubsection*{Transitions, parity and PBW}

Charge descent changes the grading before any source-to-target
comparison is made. Its formal effect on supplied algebraic data
can be stated separately.

\begin{proposition}[Additive regrading]
\label{prop:recognition-additive-regrading}
Let $q:\Gamma\to\Lambda$ be a homomorphism of charge groups and
$V=\bigoplus_{\gamma\in\Gamma}V_\gamma$ a graded super vector space.
Define
\[
 (q_*V)_{\lambda,\bar p}
   =\bigoplus_{q(\gamma)=\lambda}V_{\gamma,\bar p}.
\]
Regrading preserves both parity summands and any supplied
homogeneous filtration. For a graded Lie superalgebra $L$, it gives
\[
 q_*U(L)\simeq U(q_*L),\qquad
 q_*\operatorname{gr}_{\mathrm{PBW}}U(L)
   \simeq\operatorname{Sym}_{\mathrm{super}}(q_*L).
\]
A supplied positive/Cartan/negative direct-sum decomposition
survives as a labeled decomposition. To recover those labels from
the new charge alone, require $q$ to preserve the corresponding
positive, zero and negative supports.
\end{proposition}
\begin{proof}
Regrading retains every vector, its parity and its filtration
membership. Additivity gives $q(\gamma+\gamma')=q(\gamma)+q(\gamma')$,
so each homogeneous product, coproduct or bracket has the new
required degree. It changes neither the tensor algebra nor its
supercommutator relations. Hence it preserves the enveloping
algebra, its PBW length and its associated graded.

The three labeled summands remain direct because their underlying
subspaces have not changed. If their new charge supports overlap,
charge alone cannot distinguish them. Finite retained vector spaces
remain finite, but arbitrary charge fibers need not have finite
dimension. A completed regrading therefore requires its own
summation and continuity hypotheses.
\end{proof}

Apply this proposition to the supplied normal-ordered homomorphism
$\overline\Pi_X$. It preserves constructed parity, PBW and triangular
data. It does not supply source dimensions, representatives,
relations or the geometric maps in the recognition theorem.

An inclusion of finite root sets is an enlargement of an operation
diagram. It is not automatically a quotient morphism of Lie
algebras. In the following statements, $L_\nu$ and their transitions
are actual Lie superalgebras and morphisms, for example the
positive-height quotients. This supplies the domain needed by the
enveloping functor.

\begin{definition}[PBW transition data]
\label{def:transition-pbw-filtration-datum}
For an even Lie morphism $T^Q_{\nu+1,\nu}:L_{\nu+1}\to L_\nu$,
its PBW transition is the induced map
$U(T^Q_{\nu+1,\nu}):U(L_{\nu+1})\to U(L_\nu)$.
When weighted finite pieces are retained, assume $T^Q$ respects
that weight filtration as well. The comparison maps and transitions
are required to form commuting diagrams.
\end{definition}

\begin{proposition}[Functorial PBW filtration]
\label{prop:transition-pbw-filtration-preservation-criterion}
The induced map preserves PBW length. Its associated-graded map is
$\operatorname{Sym}_{\mathrm{super}}(T^Q)$, and these maps compose
functorially. Under the additional weight condition, they preserve
the specified finite weighted pieces.
\end{proposition}
\begin{proof}
A generator maps to a generator in $L_\nu$. A word of length at
most $d$ maps to such a word. Preservation of the Lie bracket
carries the source supercommutator relation to the target relation,
so the tensor-algebra map descends. Super PBW identifies its symbol
with the symmetric map. Both construction and composition are
functorial. The weight assertion follows by adding the weights of
generators in each word.
\end{proof}

\begin{definition}[Parity transition data]
\label{def:transition-parity-decomposition-datum}
Let $J_\nu$ be $+1$ on $V_{\nu,\overline0}$ and $-1$ on
$V_{\nu,\overline1}$. A parity-preserving transition satisfies
\[
 J_\nu T_{\nu+1,\nu}=T_{\nu+1,\nu}J_{\nu+1}.
\]
The same equation is required in each retained positive, negative
and Cartan degree, compatibly with its comparison map.
\end{definition}

\begin{proposition}[Parity blocks]
\label{prop:transition-parity-decomposition-preservation-criterion}
The equation above is equivalent to zero off-diagonal parity
blocks. It restricts the inverse system separately to both parities.
\end{proposition}
\begin{proof}
If $J_{\nu+1}x=(-1)^px$, then
$J_\nu T x=TJ_{\nu+1}x=(-1)^pTx$.
Conversely, a map preserving both eigenspaces commutes with $J$.
The restricted maps inherit identity and composition from $T$.
\end{proof}

\begin{definition}[Finite primitive towers]
\label{def:primitive-space-ml-datum}
A finite primitive tower consists of actual finite-dimensional
spaces $P_{\nu,\gamma,\bar p}$ and compatible inverse-system maps
$t_{\nu'\nu}$ preserving degree and parity. The notation does not
assert that these are projections of a symmetric root fragment as
Lie algebras.
\end{definition}

\begin{proposition}[Primitive inverse limits]
\label{prop:primitive-space-lim1-vanishing-criterion}
Every such countable tower is Mittag--Leffler and satisfies
$R^1\varprojlim_\nu P_{\nu,\gamma,\bar p}=0$.
\end{proposition}
\begin{proof}
Images in a fixed finite-dimensional term form a descending chain
and stabilize. Apply Proposition~\ref{prop:bkm-finite-tower-ml}.
No additional stabilization hypothesis is needed after these finite
spaces and their actual transition maps have been constructed.
\end{proof}

\begin{definition}[Pairing-kernel towers]
\label{def:pairing-kernel-ml-datum}
At each degree and parity, the invariant form defines
\[
 \lambda_{\nu,\gamma}:P_{\nu,\gamma}\to P_{\nu,-\gamma}^*,
 \qquad
 \rho_{\nu,\gamma}:P_{\nu,-\gamma}\to P_{\nu,\gamma}^*,
\]
with kernels $K^\ell_{\nu,\gamma}$ and $K^r_{\nu,\gamma}$.
To form a kernel tower, require that the specified transitions
carry each kernel into the corresponding kernel and compose.
For a map of inverse systems with these dual codomains, specify
the codomain transitions separately. Dualization reverses the
direction of an inverse transition.
\end{definition}

\begin{proposition}[Kernel inverse limits]
\label{prop:pairing-kernel-lim1-vanishing-criterion}
The countable towers of these finite kernels have
$R^1\varprojlim K^\ell_{\nu,\gamma}
=R^1\varprojlim K^r_{\nu,\gamma}=0$.
For an actual commuting map of inverse systems, its kernel
commutes with inverse limit without a Mittag--Leffler assumption.
\end{proposition}
\begin{proof}
Finite-dimensionality gives the first assertion. The second follows
because the kernel and the inverse limit are limits, which commute.
It makes no identification between an inverse limit of dual
spaces and a continuous dual unless the necessary variance and
topology have been supplied.
\end{proof}

\begin{definition}[Finite PBW associated-graded towers]
\label{def:pbw-associated-graded-ml-datum}
Suppose actual enveloping transitions preserve finite weighted
pieces $F_NU(L_\nu)$ and PBW length $P_dU(L_\nu)$. Their finite
associated-graded pieces are
\[
 G_{\nu,N,d,\chi,\bar p}
 =\left(\frac{F_NU(L_\nu)\cap P_dU(L_\nu)}
              {F_NU(L_\nu)\cap P_{d-1}U(L_\nu)}\right)_{\chi,\bar p}.
\]
The induced quotient maps define the corresponding inverse towers.
\end{definition}

\begin{proposition}[PBW inverse limits]
\label{prop:pbw-associated-graded-lim1-vanishing-criterion}
These towers satisfy $R^1\varprojlim G_{\nu,N,d,\chi,\bar p}=0$.
For the limit filtration defined by the finite filtered inverse
systems, taking the indicated associated-graded quotient commutes
with inverse limit.
\end{proposition}
\begin{proof}
The $G_\nu$ are finite-dimensional. The short exact sequence of
the two filtered pieces and their quotient also has
finite-dimensional kernel. Proposition~\ref{prop:bkm-finite-tower-ml}
therefore proves exactness after inverse limit and gives the
associated-graded quotient. This does not identify
$U(\varprojlim L_\nu)$ with $\varprojlim U(L_\nu)$. Such a
comparison has its own completion requirements.
\end{proof}

\begin{proposition}[Compatible limits and radical quotients]
\label{prop:strict-ml-completion-primitive-recognition}
Compatible finite-level isomorphisms $A_\nu:V_\nu\to W_\nu'$
induce an isomorphism of ordinary inverse limits without an
additional exactness hypothesis. For a compatible short exact
sequence of finite-dimensional spaces
$0\to R_\nu\to P_\nu\to Q_\nu\to0$, one has
\[
 \varprojlim Q_\nu\simeq
       (\varprojlim P_\nu)/(\varprojlim R_\nu).
\]
If the towers are given the inverse-limit linear topology,
$\varprojlim R_\nu$ is closed. These assertions preserve operations
only when those operations are defined continuously on the limits.
\end{proposition}
\begin{proof}
The inverses $A_\nu^{-1}$ are compatible and induce the inverse
map on limits. The quotient assertion follows from
$R^1\varprojlim R_\nu=0$. Its kernel is closed because it is the
intersection of the inverse images of the finite-stage kernels.
Continuous operation squares pass to limits. The obstruction in
Proposition~\ref{prop:bkm-two-sided-completion-obstruction} shows
why a vector-space limit of symmetric root fragments is not, by
these linear facts alone, a completed Lie algebra.
\end{proof}

\subsubsection*{The first enlarged degree set}

The height-three root set contains thirteen degrees. To evaluate the
first real Serre relations, doubled null degrees, complementary
strings and odd-odd products, consider the finite degree set
\[
 \begin{split}
 W_{\mathrm{rel}}^+={}&W_3
 \cup\{3\delta_i+\delta_j:i\ne j\}
 \cup\{2a_{ij}:i<j\}\\
 &\cup\{a_{ij}+s\delta_k:\{i,j,k\}=\{1,2,3\},\ 2\leq s\leq5\}
 \cup\{2\delta_{123}\}.
 \end{split}
\]
It has $13+6+3+12+1=35$ degrees, including nine zero root spaces
at the real and complementary terminal degrees. Its nonzero
proper subdegrees add exactly
$D_k=\delta_k+2a_{ij}$, giving a $38$-degree set. Thus a degree
set that includes zero codomains is different from a subset of
$\mathcal R_+$.

\begin{proposition}[Comparison spaces in the enlarged set]
\label{prop:target-admissibility-obstruction-wrel}
The full parity dimensions at
$C_{k,2}=a_{ij}+2\delta_k$, $D_k$ and $2\delta_{123}$ are
respectively
\[
 201|93,\qquad510|1023,\qquad10136|6120.
\]
Every downward-saturated root set containing $2\delta_{123}$
contains the $D_k$. Recognition there additionally requires the
source coefficient spaces, invariant pairings and the bracket and
enveloping comparison maps in those degrees.
\end{proposition}
\begin{proof}
The dimensions are Proposition~\ref{prop:bkm-height-six-parities}.
The inequalities $0<D_k<2\delta_{123}$ force their inclusion.
The operation requirements are those of
Proposition~\ref{prop:target-presentation-prerequisite}.
The signed values $108,-513,4016$ and simple coefficients
$90,54,-540$ describe different quantities and do not supply the
geometric maps.
\end{proof}

A further degree-enlargement rule illustrates why one should retain
finite operation diagrams instead of requiring a single finite
set closed under all such enlargements. For a degree $\beta$ with
$(\delta_i,\beta)<0$, define
\[
 T_i(\beta)=\beta+\bigl(1-(\delta_i,\beta)\bigr)\delta_i.
\]
On an imaginary simple generator this is its real Serre output
degree. On a composite root, this formula only specifies a test degree.
It does not assert that the corresponding power of
$\operatorname{ad}e_i$ kills the entire root space.
The distinction is detected by
Theorem~\ref{thm:bkm-serre-nonextension}.

For example, apply $T_i$ to the ten degrees
$2a_{ij}$, $C_{k,2}$, $D_k$ and $2\delta_{123}$ inside the
$38$-degree set. There are twenty-one distinct new output degrees.
Eighteen have nonzero signed coefficient. Adjoining their nonzero
proper subdegrees creates $126$ incidences across thirty distinct
new degrees, so the enlarged set has $38+21+30=89$ degrees.

These finite counts use the full Jacobi coefficients after elliptic
reduction. In particular,
\[
 \begin{array}{c|c|r}
 \beta&(n m,l)&f(nm,l)\\\hline
 (2,6,1)&(12,7)&1\\
 (8,2,2)&(16,8)&10\\
 (5,2,1)&(10,6)&108\\
 (7,2,2)&(14,7)&-513
 \end{array}
\]
and the same values hold after exchanging the first two root
coordinates. Elliptic translation reduces these coefficients to
$f(0,1)$, $f(0,0)$, $f(1,0)$ and $f(2,1)$, respectively.

A finite Fourier expansion cannot assign zero to an uncomputed
coefficient beyond its degree bound. Such a substitution would
already miscount this first enlargement.

\begin{proposition}[An unbounded affine staircase]
\label{prop:recognition-affine-staircase}
Suppose a set of nonnegative lattice degrees is closed under
$T_i$ whenever its pairing is negative. Suppose it contains all
proper subdegrees with nonzero Jacobi coefficient. If it
contains the three real simple degrees, it is infinite.
\end{proposition}
\begin{proof}
For $N\geq2$, define
\[
\begin{split}
 A_N={}&\{(0,N,N-1),(0,N,N),(0,N+1,N),\\
       &\hspace{18mm}(N,0,N-1),(N,0,N),(N+1,0,N)\},\\
 B_N={}&\{(0,N,N+1),(0,N+1,N+1),(0,N+1,N+2),\\
       &\hspace{18mm}(N,0,N+1),(N+1,0,N+1),(N+1,0,N+2)\}.
\end{split}
\]
For $\beta=(0,N+1,N)$, the pairing with $\delta_3$ is $-2$,
so $T_3(\beta)=(0,N+1,N+3)$.
Its three indicated proper subdegrees in the first half of $B_N$
have Gram $q$-degree zero and elliptic exponents $-1,0,-1$,
hence coefficients $1,10,1$. Exchanging the first two coordinates
gives the other half. Therefore $A_N$ forces $B_N$.

For $\beta=(0,N+1,N+2)$, the pairing with $\delta_2$ is $-2$,
so $T_2(\beta)=(0,N+4,N+2)$.
Its proper subdegrees include the first half of $A_{N+2}$,
again with coefficients $1,10,1$. The coordinate exchange gives
$B_N\Rightarrow A_{N+2}$. Both displayed output degrees have
zero coefficient, since their elliptic exponent has absolute
value two at $q$-degree zero.

Finally, $T_3(\delta_2)=(0,1,3)$ forces the root $(0,1,2)$.
Its $T_2$ output is $(0,4,2)$, whose nonzero proper subdegrees
include $(0,2,1),(0,2,2),(0,3,2)$. Starting with $\delta_1$
gives their coordinate exchanges, so $A_2$ is present.
Induction yields $A_{2+2k}$ and $B_{2+2k}$ for all $k\geq0$.
Their coordinates are unbounded. In particular this includes the
six-degree families at every even $N\geq14$.
\end{proof}

This is an obstruction to that finite closure rule. It does not
obstruct the positive Lie quotients or the cofinal finite operation
diagrams used above. A compact construction need only provide each
specified operation in a sufficiently large valid domain, with
compatible maps. It need not make an infinite root algebra into a
single finite algebra.

\subsubsection*{The remaining compact construction}

\begin{lemma}[The same determinant with an abelian bracket]
\label{lem:signed-dimensions-miss}
Equip the graded super vector space underlying $\mathfrak g$ with
the zero bracket. Its root superdeterminant is unchanged, but it
is not isomorphic to $\mathfrak g$ as a Lie superalgebra.
\end{lemma}
\begin{proof}
The determinant uses only the signed graded dimensions.
The new bracket is abelian, while $[e_i,f_i]=h_i\ne0$ in
$\mathfrak g$. An abelian Lie superalgebra is a valid comparison
object here. No counital Hopf structure with zero coproduct is
being asserted.
\end{proof}

\begin{lemma}[Scalar data and geometric recognition]
\label{lem:first-window-scalar-exclusions}
A scalar denominator, Pfaffian product or protected index does
not by itself determine the geometric primitive recognition datum.
Target dimensions or target PBW spaces likewise do not construct
source coefficient spaces or an intertwining map.
\end{lemma}
\begin{proof}
The preceding lemma keeps the entire signed determinant while
changing the bracket. Adding an even-odd pair also keeps it while
changing both parity dimensions. The matrices in
Definition~\ref{def:finite-source-proof-tables} require actual
source spaces, operations and geometric coefficient maps before
their comparison with the target can be formed. Those objects are
absent from the scalar formulas.
\end{proof}

The unresolved compact data include admissible correspondences,
coefficient realizations, the invariant Lie form and source
representatives. Generation, kernel equality and the preceding
coefficient comparison must also be constructed. They can be
expressed as the relation and operation defects
\[
 \mathfrak O_{\mathrm{Rec},\nu}
 =\bigl(\mathfrak o^{(R1)}_\nu,\ldots,
        \mathfrak o^{(R8)}_\nu,
        \mathfrak o^{\mathrm{tr}}_{\nu+1,\nu}\bigr),
\]
with each entry referring to the corresponding map and identity.
An unresolved geometric map cannot be replaced by a vanishing
numerical defect. The orientation and Pfaffian
data of Appendix~\ref{app:obstruction-discharges} do not alone
construct these source maps.

The target denominator is established by
\cite[Sections~3--4]{GN}, \cite[Theorem~2.1 and Section~5.1.1]{GNII}
and the Borcherds lift \cite{Borcherds95,B}. An additional Drinfeld
double, a derived center or a vacuum-evaluated algebra is a further
construction with its own tensor and pairing data. A local toric
Hall algebra does not identify the compact $K3\times E$ primitive
algebra without a map between the actual coefficient objects.
The first compact comparison must construct the source operations
on the explicit real, null and interior spaces above. These finite
comparisons must then extend to global recognition.


\subsection{Finite geometric systems and their limits}
\label{subsec:eight-finite-constructions}

The preceding constructions have different meanings of finiteness.
A bounded set of charges is combinatorial. A moduli stack of finite
type is geometric. A finite-dimensional coefficient space is a
property of a specified cohomology theory. None of these conditions
alone supplies the other two. For example,
\(H^*(B\mathbb G_m;\mathbb Q)=\mathbb Q[c_1]\) is unbounded even
though \(B\mathbb G_m\) is an Artin stack of finite type.

We organize the compact comparison by finite operation diagrams and
then by inverse systems. A finite stage must contain the input and
output spaces of the operations being compared. It need not be
closed under every extension or bracket. When a calculation leaves
the retained range, its larger output stage must be specified.
This permits finite calculations without asserting that a bounded
root fragment is a Lie algebra or that a bounded nonzero charge set
is closed under arbitrary addition.

\subsubsection*{Charges and target windows}

Fix the stability and support data of
Appendix~\ref{app:stability-support}. Let
\(F_{\sigma,S}(R)\subset\Gamma_X^{\mathrm{form}}\) be a supplied
finite set of charges with positive height bounded by \(R\), and
let \(\Gamma_R^{HN}\subset F_{\sigma,S}(R)\) be the charges of
the geometric strata to be used. The product-stability construction
has central charge
\[
 L_F(n)=Z_S\bigl(Rp_{S,*}(F\otimes p_E^*H^n)\bigr),
\]
where \(H\) has degree one on \(E\) and the K3 stability datum is
rational \cite{LiuProductStability}. Its numerical formula does not
prove boundedness of the chosen moduli problem. The relevant class
bounds, bounded families and finite-type substacks must be supplied
at their actual support and stability scope.

Let \(\Pi_X\) and \(B\) be the Gram map and normal-ordering
cocycle of Section~\ref{subsec:d6-d2-d0-dictionary-appendixC}.
The cocycle equation is
\[
 B(c,c')+B(c+c',c'')=B(c',c'')+B(c,c'+c'').
\]
It makes the central lift associative on its declared charge monoid.
Restrictions to finite sets retain only the products whose outputs
are present.

\begin{definition}[Finite target test window]
\label{def:finite-target-test-window}
For \(c\in\Gamma_R^{HN}\), set
\[
 \mathcal T_R(c)=
 \left\{\sum_{i<j}B(c_i,c_j):
   c=\sum_i c_i,\ c_i\in F_{\sigma,S}(R),\
   \sum_i h_S(c_i)\le R\right\}.
\]
Assume these translate sets are finite, and define
\[
 \begin{split}
 \widehat\Gamma_R&=
    \{(c,T):c\in\Gamma_R^{HN},\ T\in\mathcal T_R(c)\},\\
 \Gamma_R^{\mathrm{test}}&=
    \{\Pi_X(c)+T:(c,T)\in\widehat\Gamma_R\}.
 \end{split}
\]
This is a finite target-degree set. Compatible inclusions of the
charge and translate sets give inclusions of the target sets as
\(R\) increases.
\end{definition}

The image formula constructs a set of Gram degrees after its inputs
have been supplied. It does not construct sheaves representing those
degrees, nor does it make the test window exhaustive.

\begin{definition}[Exhaustion of a specified target set]
\label{def:relation-closed-target-finite-r-exhaustion}
Let \(\mathcal D_{\mathrm{rel}}\) be a specified set of degrees
needed by the target relations. A finite-stage realization of this
set assigns to each \(\gamma\in\mathcal D_{\mathrm{rel}}\) a
finite height \(R(\gamma)\), a charge
\(c_\gamma\in\Gamma_{R(\gamma)}^{HN}\), and a translate
\(T_\gamma\in\mathcal T_{R(\gamma)}(c_\gamma)\), such that
\[
 \Pi_X(c_\gamma)+T_\gamma=\gamma.
\]
The assignments are compatible with the stage inclusions. If the
target set has a finite prefix and symbolic tail families
\(\gamma_N\), the tail requires formulas
\(R_N,c_N,T_N\) valid for every allowed \(N\), with the same
image identity. All intermediate degrees of each asserted relation
must be included in an appropriate finite enlargement.
\end{definition}

\begin{proposition}[Coverage of target degrees]
\label{prop:relation-closed-target-finite-r-exhaustion-criterion}
Under this datum, every degree of \(\mathcal D_{\mathrm{rel}}\)
belongs to \(\Gamma_R^{\mathrm{test}}\) for some finite \(R\).
\end{proposition}
\begin{proof}
The pair \((c_\gamma,T_\gamma)\) belongs to
\(\widehat\Gamma_{R(\gamma)}\) and has image \(\gamma\).
The same argument applies uniformly to each symbolic family.
\end{proof}

A relation cascade, including a family of the form
\(A_N\Rightarrow B_N\Rightarrow A_{N+2}\), can describe target
degrees that must be retained. It supplies neither the geometric
charges \(c_N\) nor their normal-ordering translates. Conversely,
a finite list of represented degrees does not prove coverage of an
infinite cascade. The target calculations and the geometric
exhaustion are distinct parts of recognition.

\subsubsection*{Geometric stages and transitions}

A geometric stage consists of the chosen finite-type object stacks,
extension stacks and flag stacks, with the support and stability
conditions needed for their maps. The full CY3 deformation problem
has its shifted symplectic structure on its stated derived stack
\cite{PTVV2013}. Restricting or reducing that problem requires a
comparison preserving the relevant structure; a finite-type
substack is not automatically a two-term critical chart.
The reduced constructible coefficients and orientations are those
constructed on the actual reduced critical problem
\cite{BBDJS2015,JoyceUpmeier}.

A universal complex may live on a scalar gerbe. Its weight-zero
endomorphism complex and natural maps can descend without descent
of the weight-one universal object. Scalar rigidification likewise
does not imply finite residual inertia: direct sums can retain
positive-dimensional automorphism groups. Both facts affect the
coefficient theory and its quotient maps.

\begin{definition}[Finite geometric transition datum]
\label{def:finite-moduli-transition-geometry}
For \(R'\le R\), specify object, extension and flag morphisms
\[
 t^a_{RR'}:\mathfrak M^a_R\longrightarrow\mathfrak M^a_{R'},
 \qquad a\in\{\mathrm{obj},\mathrm{ext},\mathrm{fl}\},
\]
and closed retained substacks
\(\mathfrak Z^a_R\subset\mathfrak M^a_R\) carried into
\(\mathfrak Z^a_{R'}\). Require the restricted morphisms to be
proper; a closed immersion is one allowed special case.
For every coefficient operation, include the appropriate base-change
square and its comparison morphism. Include the maps of deformation
complexes, cosections, constructible coefficients and root pairs,
with their overlap and composition coherences.
\end{definition}

\begin{proposition}[Restricted geometric transitions]
\label{prop:finite-moduli-transition-geometry-criterion}
These data give proper transition maps on the retained substacks.
When their coefficient and base-change comparisons are equivalences,
they transport the specified coefficient operations and orientation
pairs compatibly with composition.
\end{proposition}
\begin{proof}
The retained-locus condition restricts each supplied morphism to the
chosen source and target. Properness is a hypothesis on those maps.
The coefficient assertion follows by composing their specified
pullback or pushforward comparisons with the base-change
isomorphisms. Their coherence gives equality of the composite
transports.
\end{proof}

A closed graph alone does not make a morphism a closed immersion.
Nor do inclusions of bounded HN substacks naturally give morphisms
in the opposite direction. The direction of a transition and the
variance of its coefficient functor are part of the datum. A direct
system obtained from open inclusions is not the inverse system above
without an additional construction.

\subsubsection*{The eight kinds of structure}

The following structures specify the operations that a finite
comparison must retain.
\begin{enumerate}
\item \emph{Charge and support.} The geometric classes, Gram map,
  normal-ordering cocycle and finite operation degrees are those
  above. The elliptic support degree is defined by
  \((p_E)_*\operatorname{cyc}_1(F)=b^{\mathrm{geom}}[E]\).
  It is additive on the relevant exact sequences. It is not the
  third Borcherds coordinate \(m^{\mathrm{Bch}}\) without a proved
  comparison. A positive effective elliptic degree projects onto
  all of \(E\); at degree zero the image of the curve support
  in \(E\) is a finite set.
\item \emph{Moduli and coefficients.} Specify the derived object
  and extension stacks, full flag correspondences, reduced
  deformation theory, critical coefficient system and supports.
  A proper target map can be obtained from a closed relative Quot
  locus in the sheaf recipient when its boundedness and closedness
  hypotheses hold. It does not imply properness of the source map,
  and the argument does not apply to an arbitrary tilted-heart
  extension problem without another construction.
\item \emph{Hybrid operations.} The local and wrapped sectors of
  Section~\ref{subsec:hybrid-ran-carrier} use their actual
  configuration spaces, wrapped prequotients and anchors. The
  \(LL,LW,WL,WW\) correspondences retain the source and target
  anchor data before quotient. Their three-input comparisons include
  all eight words in \(\{L,W\}^3\), and the full higher-flag
  coherences are further structure. The wrapped product of quotient
  stacks uses a common \(BE\) fibre, not two independently
  forgotten translation torsors. The quotient-after-correspondence
  pseudofunctor and protected integration maps must satisfy the
  relevant base-change, projection and Thom--Sebastiani identities.
\item \emph{Orientation and transport.} Specify root pairs on the
  actual quotient, cross-determinant factors on extensions and
  coherent flag comparisons. On a finite-stabilizer locus, a root
  obstruction and the character of a chosen root are separate
  conditions. Weyl lifts act on this same datum and preserve its
  square maps, critical coefficients and transition comparisons.
\item \emph{Hall and primitive operations.} Specify product,
  coproduct, unit, counit and the tensor in which the bialgebra
  equation is asserted. The ordinary Koszul and meromorphic
  braided tensors give different equations. A product alone gives
  neither cooperation nor a primitive projection. In a connected
  cocommutative ordinary realization over characteristic zero,
  the Eulerian projection has the hypotheses of
  Section~\ref{subsec:primitive-recognition}. The invariant Lie
  form, its radical and its generated Hopf ideal are distinct from
  any separately supplied full Hopf pairing.
\item \emph{Pfaffians.} Supply the finite skew complexes, their
  shifts and square maps, the Pfaffian lines and sections, and the
  geometric wall charts. The actual parameter and automorphic-line
  comparisons fix the target and its normalization. Positive-height
  product convergence and analytic continuation are separate from
  geometric effectivity of the family.
\item \emph{Associative coefficients.} Supply the counital,
  coaugmented source coalgebra with its conilpotence or the precise
  replacement filtration, and the associative cobar recipient of
  Section~\ref{subsec:koszul-source}. The finite comparison maps
  preserve its differential, support, tensor and operations.
  Lie--commutative chiral duality is not a substitute for this
  associative comparison.
\item \emph{Recognition.} Supply geometric simple and Cartan
  representatives, relation maps, kernel equalities, generation,
  both root parities and PBW comparisons. A finite calculation must
  retain all intermediate outputs. Compatible comparisons in every
  bounded weighted word space give the all-level algebraic
  assertion; a topological conclusion also requires continuous
  inverse maps on a specified completion.
\end{enumerate}

For wrapped sheaves, a determinant anchor can record translation
when its Euler characteristic is nonzero. At Euler characteristic
zero that information can vanish. Jordan--H\"older data may then
separate additional objects but do not automatically determine an
elliptic translation orbit. Additional anchors must be chosen and
their separating property proved. Similarly, a change of origin can
change a line's Picard class without determining an equivariant
stabilizer character. The constructions and examples of
Section~\ref{subsec:hybrid-ran-carrier} keep these distinctions.

On the formal nilpotent support of the local deformation problem,
the cyclic local model and
its full flag extensions provide genuine local operations. Their
comparison with a global reduced compact coefficient functor, its
proper correspondence maps and quotient descent remains a separate
geometric construction. The finite organization above preserves
this positive local result without treating a list of required
objects as proof of global existence.

\subsubsection*{Morphisms of finite comparison data}

\begin{definition}[Finite-stage Dirac--Igusa datum and morphism]
\label{def:finite-stage-dirac-igusa-morphism}
In a fixed coefficient setting, a stage is a tuple
\[
 \begin{split}
 \mathfrak D^{\mathrm{DI}}_{X,R}=
 (&\mathcal A_{X,R}^{E_3},F_{X,R}^{\mathrm{hyb}},
   \Gamma_{X,R},\Pi_X,\Phi_R,o_R,H_R,P_R^\Pi,\\
  &C_{X,R},\Theta_{\mathrm{Kos},R},
   \mathscr L_{\mathrm{Pf},R},\operatorname{pf}_{X,R},
   \epsilon_{o,R},\mathrm{Rec}_R),
 \end{split}
\]
whose components carry the specified structures above. In particular,
an \(E_3\)-source is actual additional operadic data when that
component is required; it is not inferred from its scalar series.
A morphism consists of the following compatible maps.
\begin{enumerate}
\item[(M1)] Maps of charge and support diagrams preserving
  \(\Pi_X\), the cocycle \(B\) and the retained products.
\item[(M2)] Geometric and coefficient transition maps as in
  Definition~\ref{def:finite-moduli-transition-geometry}, with the
  required proper support and variance.
\item[(M3)] Maps of the hybrid carriers and all specified
  coloured-flag operations, compatible with the quotient
  pseudofunctor and protected integration.
\item[(M4)] Isomorphisms of the pulled-back orientation pairs,
  compatible Weyl lifts and wall charts, preserving their square
  maps, normal coordinates and section comparisons.
\item[(M5)] Maps of the chosen Hall data preserving product,
  coproduct, units and counits in their declared tensor. Preserve
  primitives, the invariant Lie form and its radical; require a
  further Hopf-pairing comparison only where that pairing is
  independently supplied.
\item[(M6)] Isomorphisms of pulled-back Pfaffian lines carrying
  the specified sections to one another and commuting with their
  square and automorphic-line comparisons.
\item[(M7)] Coalgebra maps preserving the differential, counit,
  coaugmentation, filtration and the associative comparison
  \(\Theta_{\mathrm{Kos}}\).
\item[(M8)] Maps of the recognition diagrams preserving the
  representatives, relations, kernels, parity pieces, invariant
  forms, radicals and PBW filtrations, with all required
  intermediate operation spaces.
\end{enumerate}
Composition is componentwise, with the coherence required by the
chosen coefficient category. Fix strict models when an ordinary
category \(\mathsf C_X\) is used below. Alternatively retain an
enhanced category and all its mapping spaces. A finite-stage
isomorphism is an invertible morphism of these structured data,
not merely an isomorphism of their underlying vector spaces.
\end{definition}

The preservation results in the preceding subsections apply to
these maps at their individual scopes: orientation pairs and
critical coefficients, Hall products and coproducts, primitive
subspaces, invariant-form radicals, PBW filtrations and parity.
They do not follow from a scalar equality. In particular, maps
between nonorthogonal matrix quotients must be checked before a
radical or pairing is transported.

\subsubsection*{Pro-objects and actual inverse limits}

\begin{definition}[The Dirac--Igusa pro-object]
\label{def:dirac-igusa-pro-isomorphism}
A compatible cofiltered diagram
\(R\mapsto\mathfrak D^{\mathrm{DI}}_{X,R}\) in \(\mathsf C_X\)
defines the pro-object
\(\mathfrak D_X^{\mathrm{DI}}=\{\mathfrak D^{\mathrm{DI}}_{X,R}\}_R\).
A pro-isomorphism is an invertible morphism in
\(\operatorname{Pro}(\mathsf C_X)\).
No componentwise inverse limit is implicit in this notation.
\end{definition}

\begin{proposition}[Mapping formula and cofinal invariance]
\label{prop:dirac-igusa-pro-object-universal-property}
For two such systems \(\mathfrak E=(\mathfrak E_S)_S\) and
\(\mathfrak D=(\mathfrak D_R)_R\),
\[
 \operatorname{Hom}_{\operatorname{Pro}(\mathsf C_X)}
       (\mathfrak E,\mathfrak D)
 =\varprojlim_R\varinjlim_S
       \operatorname{Hom}_{\mathsf C_X}(\mathfrak E_S,\mathfrak D_R).
\]
For a constant source \(T\), this reduces to
\[
 \operatorname{Hom}_{\operatorname{Pro}(\mathsf C_X)}(T,\mathfrak D)
 =\varprojlim_R\operatorname{Hom}_{\mathsf C_X}(T,\mathfrak D_R).
\]
Cofinal restriction gives an isomorphic pro-object. Compatible
finite-stage isomorphisms on common cofinal systems are a sufficient
construction of a pro-isomorphism.
\end{proposition}
\begin{proof}
The first formula defines the morphisms of the pro-category. A
constant source removes the inner colimit. Cofinal restriction
preserves these mapping functors, hence the represented pro-object.
Compatible finite-stage inverses give inverse elements in the two
mapping sets, proving the last assertion. In an enhanced
pro-category, the corresponding statement uses mapping spaces and
homotopy limits and colimits.
\end{proof}

The sufficient condition is not necessary in a general pro-category.
For example, the vector-space tower
\[
 \mathbb C\xleftarrow{0}\mathbb C\xleftarrow{0}\mathbb C
       \xleftarrow{0}\cdots
\]
is pro-zero: for every target stage, the next transition kills the
identity class in the defining colimit. Every stage is nevertheless
nonzero, so no cofinal set of stages is levelwise isomorphic to the
zero tower. Replacing the zero transitions by identities gives the
constant nonzero pro-object. Thus the transition maps are decisive.

When an actual inverse limit exists, it belongs to a specified
complete category. For modules its projection topology has basic
neighbourhoods given by kernels of the stage projections. For line
or root-pair components the appropriate construction is a limit of
Picard groupoids or stacks, with its descent data. Neither existence
of a pro-object nor Mittag--Leffler vanishing for some cohomology
groups proves geometric effectivity of all its components.

\begin{definition}[Degreewise Mittag--Leffler condition]
\label{def:moduli-cohomology-ml-datum}
Choose a countable cofinal sequence and coefficient complexes
\(K_{R,a}\) on the retained object, extension, wrapped or flag
stacks. Let
\[
 C^q_{R,a}=H_c^q(\mathfrak M_{R,a},K_{R,a})
\]
when compactly supported cohomology is the selected theory, and use
the specified dualizing-complex model in a Borel--Moore theory.
The actual six-functor maps define a tower only after their variance
and identity and composition laws have been fixed.
It is Mittag--Leffler in degree \(q\) if, for every \(R_0\), the
images of \(C^q_{R,a}\to C^q_{R_0,a}\) are constant for all
sufficiently large \(R\).
\end{definition}

\begin{proposition}[Degreewise exactness of a countable limit]
\label{prop:moduli-cohomology-lim1-vanishing-criterion}
For a countable tower of modules satisfying the preceding condition,
\(R^1\varprojlim_R C^q_{R,a}=0\).
For a tower of complexes \(C_R\), there is a short exact sequence
\[
 0\longrightarrow\varprojlim{}^1 H^{q-1}(C_R)
 \longrightarrow H^q(R\varprojlim C_R)
 \longrightarrow\varprojlim H^q(C_R)\longrightarrow0.
\]
In particular, degreewise Mittag--Leffler cohomology gives the
right-hand identification.
\end{proposition}
\begin{proof}
The derived limit of a countable module tower is represented by the
two-term product complex with map
\[
 1-T:\prod_n A_n\longrightarrow\prod_n A_n,
 \qquad (1-T)(x)_n=x_n-u_nx_{n+1}.
\]
For a surjective tower this map is surjective by successive lifting.
For a Mittag--Leffler tower, its stabilized images form a
surjective subtower. The quotient is pro-zero. On a pro-zero tower,
\(\sum_{j\ge0}T^j\) is finite in each fixed output coordinate
and inverts \(1-T\). The exact sequence of these towers therefore
gives surjectivity of \(1-T\) on the original tower, proving the
vanishing of its cokernel \(\varprojlim{}^1\).
Applying the same product-cone model to a tower of complexes gives
the displayed cohomology sequence.
\end{proof}

Every tower of finite-dimensional vector spaces over a fixed field
is Mittag--Leffler: its images in a fixed stage form a descending
chain of subspaces. The analogous assertion for finite-rank
integral modules is false, as the tower
\(\mathbb Z\xleftarrow{2}\mathbb Z\xleftarrow{2}\cdots\)
shows. The images \(2^n\mathbb Z\) never stabilize.
Finite type of a stack, a proper transition map or an Euler
characteristic does not establish the finite-dimensional or
Mittag--Leffler hypotheses for a chosen coefficient theory.

\subsubsection*{Finite and limiting comparison statements}

\begin{theorem}[Formation of the finite comparison system]
\label{thm:finite-stage-dirac-igusa-pro-object}
Suppose the eight kinds of structure above are constructed, with
morphisms \textup{(M1)}--\textup{(M8)} satisfying their identity and
composition laws on a cofiltered stage system. They define a
Dirac--Igusa pro-object. The primitive, pairing-kernel, PBW and
coefficient cohomology limits agree with the cohomology of their
derived limits whenever each of those towers satisfies the
Mittag--Leffler hypotheses just stated. No geometric effectivity or
analytic continuation is implied by this conclusion.
\end{theorem}
\begin{proof}
The component maps and their coherence define the diagram in
\(\mathsf C_X\), hence its pro-object. Apply the preceding
proposition to each of the indicated coefficient towers. The last
sentence distinguishes these limits from the further geometric and
analytic realizations.
\end{proof}

\begin{theorem}[Finite Pfaffian comparison]
\label{thm:pfaffian-dirac-finite-stage}
At a finite stage carrying the geometric skew and orientation data,
the Pfaffian product and local reflection identities hold under
their specified line comparisons. A primitive recognition or
associative cobar conclusion holds only with the corresponding
finite source maps and relations. If these comparisons and their
inverses commute with every stage transition, they give the stated
comparison of pro-objects and, when defined, of their compatible
limits.
\end{theorem}
\begin{proof}
The finite skew identity is
\(\operatorname{pf}(D)^2=\det D\), and orthogonal block sums
multiply Pfaffians. Section-compatible wall lifts give the local
sign of Proposition~\ref{prop:rank-one-crossing-o2-sign}.
The primitive and coefficient conclusions are the separate
recognition and associative comparison theorems of the preceding
subsections. Compatible inverse maps induce a pro-isomorphism by
Proposition~\ref{prop:dirac-igusa-pro-object-universal-property};
applying an available limit functor preserves that equivalence.
\end{proof}

For \(K3\times E\), the finite moduli, reduced coefficients,
orientation transport and compact comparison maps remain the
geometric inputs. The local cyclic charts, formal charge formulas,
target relations and finite skew identities provide concrete parts
of the construction. A global realization requires their compatible
geometric extensions, rather than only their scalar shadows.


\subsection{The four geometric comparison conditions}
\label{subsec:four-named-obstructions-ch5}

The compact comparison needs a reduced coefficient theory, compatible
orientation pairs, Weyl transport and geometric wall charts. We denote
these requirements by \((D0)\), \((O1)\), \((O1)^+\) and \((O2)\).
They are existence and compatibility conditions; they are not four
classes in a common cohomology group. The Kummer and group-cocycle
classes below detect particular parts of them.

The tuple of Definition~\ref{def:G-retained-obstruction-datum}
keeps these inputs together with finite geometric Pfaffians:
\[
 (D0_{\mathrm{HN}},O1_{\mathrm{quot}},\mathcal T_W,
       O2_{\mathrm{atlas}},P_{\mathrm{fin}}).
\]
Primitive recognition, an associative coefficient comparison and a
protected state require additional data. In particular, constructing
this tuple would not identify the compact primitive algebra from its
superdimension alone.

\subsubsection*{Degeneration and reduced coefficients}

At each finite geometric stage \(R\), condition \((D0)\) asks for a
specified degeneration or specialization to a zero-dimensional
sector. Its domain must be stated before using flatness. On a fixed
projective target, a flat family of coherent sheaves has constant
Hilbert polynomial. A sheaf with nonzero effective curve class has a
positive linear coefficient, whereas a zero-dimensional sheaf has a
constant polynomial. Such a family cannot specialize the former to
the latter. Likewise a family of perfect complexes preserves its
Chern character in a fixed cohomological marking.

A proposed comparison changing the target, support theory or charge
therefore has to construct that change. The datum
\(D0_{\mathrm{HN}}\) specifies its geometric recipient, coefficient
functor and transition maps, together with the reduced deformation
and orientation comparisons. It must also provide specialization of
the constructible critical coefficients, Pfaffian lines and protected
integration maps on every required correspondence. The boundary
\(b_R^{\mathrm{geom}}\to0\) in a hybrid Ran carrier is not by itself
such a geometric specialization.

For extension maps \(q,p_1,p_2\), a compatible semiregularity
cosection satisfies
\[
 q^*\operatorname{CS}_{C}
   =p_1^*\operatorname{CS}_{A}+p_2^*\operatorname{CS}_{B}
\]
on the common deformation complex, with the stated reductions and
shifts. This equality is required throughout the specialization and
under the stage transitions. The nonzero-K3-class surjectivity
criterion of Section~\ref{subsec:cosection-twisted-orientation} does
not assert surjectivity in curve class zero.

A perfect reduced deformation complex, a cosection-localized virtual
cycle and an oriented critical coefficient are different objects.
The orientation constructions of \cite{JoyceSong,JoyceUpmeier}
do not identify all three or supply every specialization above.
Theorem~\ref{thm:G-D0} describes a consequence of a supplied compatible
system. Failure of this proposed degeneration route does not rule out
an independently constructed reduced Hall theory.

\subsubsection*{Quotient orientation}

Let \(K^{\mathrm{red}}\) be the specified perfect reduced complex on
a geometric coefficient base \(B\), and write
\(\mathscr D^{\mathrm{red}}=\det K^{\mathrm{red}}\).
Condition \((O1)\) asks for a root pair
\[
 (o,\rho),\qquad
 \rho:o^{\otimes2}\xrightarrow{\sim}\mathscr D^{\mathrm{red}},
\]
on the required quotient stack, with coherent extension-flag
comparisons and stage transitions.
If a cosection correction is included, it must have a perfect model
\(C_{\mathrm{cs}}\). The actual line to be rooted is then
\[
 \mathscr D_{\mathrm{tw}}=
 \mathscr D^{\mathrm{red}}\otimes\det C_{\mathrm{cs}}.
\]
Its obstruction is the Kummer boundary
\[
 \alpha^{\mathrm{red}}=\delta(\mathscr D_{\mathrm{tw}})
       \in H^2(B,\mu_2).
\]
In analytic cohomology it is the mod-two reduction of the integral
first Chern class. Rational Grothendieck--Riemann--Roch does not
recover its possible two-torsion. The vanishing of
\(\alpha^{\mathrm{red}}\) permits a root pair; a choice of pair
and its extension-flag coherence remain necessary.

For an \(E\)-action, the equivariant root problem is on \([B/E]\).
On a free-action locus this is the quotient itself. There is no
universal obstruction obtained merely by writing
\[
 a_1u_1+a_2u_2\in H^2(BE;\mathbb F_2).
\]
Such a class contributes only through the actual pullback to the
chosen quotient problem; that pullback need not be injective.

On a finite-stabilizer stratum with stabilizer \(H\), one must
solve the equivariant Kummer problem
\(\widetilde\beta^H\in H^2_H(B;\mu_2)\).
For coarse descent, every stabilizer must also act trivially on the
chosen root fibre. Thus a zero gerbe obstruction does not remove a
nonzero degree-one sign character
\(\lambda^H\in H^1(BH;\mathbb F_2)\).
For \(H=E[2]\), the three order-two restrictions detect the three
coefficients of the degree-two class. For \(E[N]\), with
\(4\mid N\), the mixed class \(e_1e_2\) in the two-primary
cohomology is invisible on every cyclic subgroup. Its coefficient
\(A_{12}^{(N)}\) must be checked independently. These are the
finite-group calculations underlying Theorem~\ref{thm:G-O1}.

Orientation enters the critical coefficient through its sign
torsor and the corresponding constructible local system. An
ordinary tensor of a coherent root line with a constructible complex
is not the coefficient construction. Nor does descent of the
weight-zero endomorphism complex imply descent of the weight-one
universal sheaf.

\subsubsection*{Coherent Weyl transport}

Set \(W=W^{(2)}(\Lambda_{II}^{2,1})\).
Condition \((O1)^+\) supplies geometric maps between the required
strata and lifts of their action to the root pairs. At a finite stage
these maps may define only a partial action groupoid
\(\mathcal G_R\). Its composition defect belongs to
\[
 c_{o,R}\in Z^2(\mathcal G_R;\underline{\mu_2}),
\]
with the actual action on the sign coefficients. A one-cochain
\(\eta_R\) with \(\delta\eta_R=c_{o,R}\) rectifies the lifts.
One uses \(H^2(W;\mathbb F_2)\) only after identifying the relevant
complete action and constant coefficient system with that group
cohomology problem.

The corrected lifts must satisfy \(\tau_i^2=1\), preserve the
quotient root and its square map, and commute with stage
transitions. The torsor defect \(\omega_{i,\mathcal C}\) measuring
compatibility with the quotient orientation must vanish with a
specified trivialization. The free-product Coxeter presentation has
no finite braid relations between distinct generators. It therefore
simplifies the relations to verify, but it does not construct the
lifts or their scalar signs. The sign comparison in
Theorem~\ref{thm:G-O1-plus} additionally uses the local sections of
\((O2)\).

\subsubsection*{Pfaffian wall charts}

Condition \((O2)\) supplies the wall atlas of
Definition~\ref{def:appE-O2}. At each required simple wall the
geometric skew complex is compared with a tangential factor and the
normal rank-two block
\[
 [\mathcal O^2\xrightarrow{D_u}(\mathcal O^2)^\vee],
 \qquad D_u=\begin{pmatrix}0&u\\-u&0\end{pmatrix}.
\]
The normal Pfaffian is \(u\); its determinant is \(u^2\).
An invariant tangential Pfaffian unit and a section-compatible lift
of \(u\mapsto-u\) give the local residual sign \(-1\).
The comparison includes the required shift of the geometric
self-Ext complex. It is not inferred from the scalar complex
\([\mathcal O\xrightarrow{u}\mathcal O]\).

The charts must cover the geometric locus, with compatible
splittings, roots, sections and integration maps on components,
boundaries and overlaps. Their transitions must agree with the
Weyl lifts. Theorem~\ref{thm:G-O2-orbit-transport} transports such a
supplied system. The local matrix computation supplies none of its
global covering or gluing maps.

The set \(\{\pm1\}^{12}\) has \(4096\) elements. To regard it as
a geometric atlas count requires a separate bijection, including
stabilizer weights and overlap identifications. Its cardinality does
not determine a reflection character. The equality \(4096=64^2\)
in the scalar normalization
\(Z^X_{\mathrm{OP}}=-4096\Delta_5^{-2}\) is independent of that
geometric counting problem.

\subsubsection*{The joint conditional comparison}

\begin{theorem}[Pfaffian and primitive comparison from geometric data]
\label{thm:pfaffian-dirac-conditional}
Suppose the four geometric conditions above are supplied on a common
cofinal system, together with the following additional data.
\begin{enumerate}
\item The finite geometric Pfaffian datum of
  Definition~\ref{def:finite-geometric-pfaffian-datum}, including the
  parameter and line maps, section normalization and the effective
  formal or analytic limit required for the conclusion.
\item The primitive recognition data \textup{(S1)}--\textup{(S10)}
  of Theorem~\ref{thm:primitive-recognition}, or their compatible
  finite realization with every intermediate operation degree
  retained and the stated convergence hypotheses.
\item The associative coefficient source and the finite comparison
  of Definition~\ref{def:finite-koszul-comparison-datum}, on its
  declared tensor and support, with compatible inverse maps and the
  derived-limit hypotheses needed to pass to the chosen completion.
\end{enumerate}
Then the geometric Pfaffian is identified with \(\Delta_5\) by its
specified line comparison, and
\[
 \epsilon_o=\widetilde\nu_{\Delta_5}|_W,
 \qquad
 P_X^{\Pi,\mathrm{red}}\simeq\mathfrak g_{\Delta_5}.
\]
The third datum additionally identifies the source cobar with the
associative Ran target \(\mathcal A_{\Delta,E}\) of
Section~\ref{subsec:koszul-source}. Each conclusion is at the
coefficient, topology and operation scope of its own hypotheses.
\end{theorem}
\begin{proof}
The wall models and coherent lifts give the generator signs
\(-1\), hence the character comparison by
Theorem~\ref{thm:appE-local-orientation-match} and the Weyl action.
The finite skew Pfaffians multiply by the finite product formula.
The stipulated parameter, line and limit comparisons identify that
section with the normalized automorphic product.

For primitives, the geometric representatives define the
presentation map. Equality of its relation kernel and the defining
ideal gives injectivity; generation gives surjectivity, as in
Theorem~\ref{thm:primitive-recognition}. The independent coefficient
maps give the cobar comparison and preserve the required Weyl,
orientation, pairing and filtration data by
Proposition~\ref{prop:finite-koszul-comparison-compatibility}.
Passing to the specified limit uses its derived-limit hypotheses,
not just the existence of the finite objects.
\end{proof}

The cofinal collection is first a pro-object
\(\{\mathfrak D_{X,R}^{\mathrm{DI}}\}_R\) in the category where
its maps and operations have been defined. An inverse limit is an
additional construction in a specified complete category. The
four geometric conditions, primitive recognition and the
associative comparison are separate parts of that construction.
Automorphic catalogue entries fix target products and characters;
they do not supply the compact quotient orientation, Weyl lifts,
Pfaffian atlas or geometric representatives.


\section{Assembly of the geometric comparison}

The constructions above have supplied formal local models and
finite algebraic comparisons. Their compact extension can now be
stated as a conditional assembly theorem, with each operation in
its own recipient.

\begin{theorem}[The structured comparison system]
\label{thm:ch5-construction-of-dirac-igusa-pro-object}
Suppose a common cofiltered family of geometric stages carries
the following compatible structures.
\begin{enumerate}
\item[\textup{(L1)}] The local, wrapped and mixed carriers of
Section~\ref{subsec:hybrid-ran-carrier}, with actual anchors,
common-torsor quotients and all required coefficient and flag maps.
When an \(E_3\)-source and chiral specialization are asserted,
their field, operadic, anomaly, support and descent hypotheses
are supplied on that same system.
\item[\textup{(L2)}] Reduced orientation pairs, their quotient
and cross-determinant comparisons, coherent Weyl lifts and stage
maps, as in Section~\ref{subsec:cosection-twisted-orientation}.
Trivial stabilizer characters are required wherever coarse
descent is part of the comparison.
\item[\textup{(L3)}] The Hall product, coproduct, unit and counit
in their declared tensor, with the coefficient maps of
Section~\ref{subsec:hall-correspondences}. Primitive projection
and a full Hopf pairing are additional structures with their own
hypotheses. The invariant Lie form on primitives has its separate
radical quotient. A double comparison also requires an opposite
half, cross relations, integral form and completion; stable-envelope
transport is used only when its geometric maps have been supplied.
\item[\textup{(L4)}] The actual skew complexes, Pfaffian lines,
sections and parameter and automorphic-line comparisons of
Section~\ref{subsec:dirac-pfaffian}, including the boundary
normalization and the stated geometric or analytic effectivity.
\item[\textup{(L5)}] The source coalgebra and the associative
comparison of Section~\ref{subsec:koszul-source}, with the
finite comparison maps and their inverses. A bar of the full
factorization algebra is used only after its higher operations
and coherences exist. An ordinary bar and a current envelope
remain separate recipients.
\item[\textup{(L6)}] The compact source representatives, relation
kernel, generation, both parities, invariant form and PBW data of
\textup{(S1)}--\textup{(S10)}, or their exhaustive finite
realization in Definition~\ref{def:cofinal-primitive-recognition-datum}.
\item[\textup{(L7)}] Morphisms \textup{(M1)}--\textup{(M8)}
satisfying the identity and composition laws of
Definition~\ref{def:finite-stage-dirac-igusa-morphism}, with the
specified derived-limit and effectivity hypotheses wherever an
actual limit is claimed.
\item[\textup{(L8)}] The compatible degeneration, quotient
orientation, Weyl and wall-atlas data
\((D0)\), \((O1)\), \((O1)^+\), \((O2)\) and the finite Pfaffian datum
\(P_{\mathrm{fin}}\) of
Theorem~\ref{thm:G-four-obstruction-criterion}.
\end{enumerate}
Then the stages define the Dirac--Igusa pro-object. Their supplied
primitive maps identify \(P_X^{\Pi,\mathrm{red}}\) with
\(\mathfrak g_{\Delta_5}\) at the stated algebraic or completed
scope. Their associative source and Pfaffian comparisons have the
respective conclusions of their comparison theorems. An operator
trace requires the independent cyclic state datum of
Appendix~\ref{app:obstruction-discharges}.
\end{theorem}
\begin{proof}
The component maps in \textup{(L7)} define the cofiltered diagram
of structured finite stages. It is the asserted pro-object by
Proposition~\ref{prop:dirac-igusa-pro-object-universal-property}.
The geometric representatives of \textup{(L6)} define the
presentation map; its stated kernel equality and generation give
the primitive isomorphism. The maps and inverses in
\textup{(L5)} give the associative comparison. The skew sections,
line comparisons and normalizations in \textup{(L4)} and
\textup{(L8)} give the Pfaffian theorem. Each passage to a
completion uses that component's exactness and effectivity
hypotheses. None of these constructions selects the cyclic state
needed for the final trace assertion.
\end{proof}

The parameter datum \((S,H,\sigma,E,0_E)\) also requires its
actual base-change maps. Definitions
\ref{def:admissible-dirac-igusa-parameter-datum} and
\ref{def:polarization-origin-dependence} specify those maps and
the polarization and origin conventions. The origin change of
Proposition~\ref{prop:elliptic-origin-independence-character}
compares determinant Picard classes; an equivariant root character
requires its further lift. A scalar denominator determines none
of these base-change or pro-object morphisms.
